% Lesson 16 — Listening: Texts about requesting for services — Anglais 1ère A — Cours signé OnBuch+
% Programme officiel MINESEC. Compiler avec : tectonic ENA17-listening-texts-about-requesting-for-services.tex
\newcommand{\DOCMATIERE}{Anglais}
\newcommand{\DOCNIVEAU}{1ère A}
\newcommand{\DOCMODULE}{Chapter 2 : Economic Life and Occupations}
\newcommand{\DOCLECON}{ENA17}
\newcommand{\DOCTITRE}{Lesson 16 — Listening: Texts about requesting for services}
% OnBuch+ — préambule « cours signés » ENRICHI (compilé avec tectonic / XeLaTeX)
% Système visuel « Pop » d'OnBuch+ : fond crème (jamais blanc), bordures encre
% épaisses, ombres « dures » décalées, titres Archivo Black en capitales.
% Version MATHÉMATIQUES : graphes (pgfplots/tikz), boîtes pédagogiques
% (définition, propriété, méthode, attention, le savais-tu, exercice résolu,
% exercice, TP, bilan), page de garde et signature OnBuch+ ancrée en bas.
%
% Macros attendues dans main.tex AVANT \input{preamble} :
%   \DOCMATIERE  \DOCNIVEAU  \DOCMODULE  \DOCLECON (numéro)  \DOCTITRE
\documentclass[11pt]{article}

\usepackage{fontspec}
\usepackage[english,french]{babel} % français = langue principale ; l'anglais via \eng{..} / textebox
\usepackage[a4paper,margin=2cm,top=3.4cm,bottom=2.8cm,headheight=1.4cm,footskip=1.3cm]{geometry}
\usepackage{amsmath,amssymb,amsfonts}
\usepackage{mathtools} % \xmapsto, \coloneqq... souvent utilisés par les modèles
\usepackage{cancel} % simplifications barrées (bilans de forces)
% \abs{z} (module, valeur absolue) et \norm{u} (norme), utilisés par les modèles
\providecommand{\abs}{}\let\abs\relax\DeclarePairedDelimiter{\abs}{\lvert}{\rvert}
\providecommand{\norm}{}\let\norm\relax\DeclarePairedDelimiter{\norm}{\lVert}{\rVert}
\usepackage[table]{xcolor} % [table] : fournit aussi \rowcolors (alternance de lignes)
\usepackage{pagecolor}
\usepackage{tikz}
\usetikzlibrary{arrows.meta,positioning,calc,shapes.geometric,shapes.misc,shapes.arrows,decorations.pathmorphing,decorations.pathreplacing,decorations.markings,patterns,shadows,mindmap,trees,fit,backgrounds,matrix,intersections,angles,quotes}
\usepackage{pgfplots}
\usepackage[version=4]{mhchem} % équations nucléaires \ce{^{235}_{92}U}
\usepackage[european,siunitx]{circuitikz} % schémas électriques
\pgfplotsset{compat=1.18}
\usepgfplotslibrary{fillbetween,groupplots} % aires entre courbes (calcul intégral)
\usepackage{enumitem}
\usepackage{fancyhdr}
% [most] charge toutes les bibliothèques tcolorbox (listings, minted, raster,
% documentation...) et gonfle inutilement la mémoire TeX sur un document long
% (~300 boîtes) : on ne charge que ce qui est réellement utilisé.
\usepackage[skins,breakable]{tcolorbox}
\usepackage{graphicx}
\usepackage{colortbl}
\usepackage{booktabs}
\usepackage{tabularx}
\usepackage{diagbox} % case d'en-tête diagonale des tableaux à double entrée
% Colonnes extensibles centrée / gauche / droite (C, L, R), souvent utilisées par les modèles
\newcolumntype{C}{>{\centering\arraybackslash}X}
\newcolumntype{L}{>{\raggedright\arraybackslash}X}
\newcolumntype{R}{>{\raggedleft\arraybackslash}X}
\usepackage{array}
\usepackage{multicol}
\usepackage{multirow}
\usepackage{siunitx}
% parse-numbers=false : accepte aussi \SI{2,5 \times 10^{-3}}{...}
\sisetup{locale=FR,output-decimal-marker={,},group-minimum-digits=4,parse-numbers=false}
\DeclareSIUnit\ohm{\ensuremath{\symup{\Omega}}} % U+2126 absent de la police
\DeclareSIUnit\division{div} % réglages d'oscilloscope (ms/div, V/div)
\DeclareSIUnit\tr{tr} % tours (vitesse de rotation, ex. \SI{1500}{\tr\per\minute})
\DeclareSIUnit\molar{M} % concentration molaire (ex. \SI{3}{\molar}, \SI{1}{\micro\molar})
\DeclareSIUnit\an{an} % année (le modèle confond parfois avec \year, un compteur TeX) — ex. \SI{5}{\kilo\meter\per\an}
\DeclareSIUnit\FCFA{FCFA} % franc CFA (ex. \SI{500}{\FCFA\per\kilo\gram})
\DeclareSIUnit\degreeBrix{\textdegree Brix} % teneur en sucre d'un sirop/jus (réfractométrie)
\DeclareSIUnit\month{mois} % le modèle utilise parfois \month, un compteur TeX, comme unité de durée
\usepackage{qrcode}
% \degree : commande fréquemment utilisée par les modèles pour un angle ou une
% température en dehors de \SI{}{} (non fournie par les paquets chargés).
\providecommand{\degree}{^{\circ}}
% \qed (fin de démonstration), utilisé par les modèles sans amsthm
\providecommand{\qed}{\hfill$\square$}
% \barre : double barre des valeurs interdites dans un tableau de variations (tkz-tab)
\providecommand{\barre}{\ensuremath{\|}}
% environnement proof (amsthm non chargé) utilisé par les modèles
\ifdefined\proof\else\newenvironment{proof}[1][Démonstration]{\par\noindent\textit{#1.}\ }{\qed\par}\fi
\usepackage{lastpage}
\usepackage{hyperref}
\hypersetup{hidelinks}

% ============================================================
% Palette « Pop » OnBuch+ — mêmes hex que la classe P (plus_ui.dart)
% ============================================================
\definecolor{poporange}{HTML}{F59321}
\definecolor{popdark}{HTML}{E07A0C}
\definecolor{popcream}{HTML}{FFF6E8}
\definecolor{popink}{HTML}{1C1714}
\definecolor{poppurple}{HTML}{7A5AE0}
\definecolor{popgreen}{HTML}{2E9E6A}
\definecolor{poppink}{HTML}{E0457B}
\definecolor{popblue}{HTML}{2F7DE1}
\definecolor{popgold}{HTML}{C98A12}
\definecolor{popmuted}{HTML}{8A7F76}
\definecolor{popline}{HTML}{EADFD2}
\definecolor{popgray}{HTML}{8A7F76}
\colorlet{popgrey}{popgray}
% Alias tolérants (le modèle nomme parfois les couleurs différemment)
\colorlet{popwhite}{white}
\colorlet{popblack}{popink}
\colorlet{popred}{poppink}
\colorlet{popyellow}{popgold}
% Teintes claires (fonds de tableaux/encadrés) — le modèle nomme parfois
% les couleurs avec un suffixe L (ex. popblueL) sans les avoir définies.
\colorlet{poporangeL}{poporange!20}
\colorlet{popdarkL}{popdark!20}
\colorlet{poppurpleL}{poppurple!20}
\colorlet{popgreenL}{popgreen!20}
\colorlet{poppinkL}{poppink!20}
\colorlet{popblueL}{popblue!20}
\colorlet{popgoldL}{popgold!20}
\colorlet{popredL}{popred!20}
\colorlet{popgrayL}{popgray!20}
% Teintes claires pour fonds de tableaux / zones
\colorlet{popblueL}{popblue!10!white}
\colorlet{poporangeL}{poporange!14!white}
\colorlet{popgreenL}{popgreen!12!white}
\colorlet{poppurpleL}{poppurple!12!white}
\colorlet{poppinkL}{poppink!12!white}
\colorlet{popgoldL}{popgold!14!white}

\pagecolor{popcream}

% ============================================================
% Typographie
% ============================================================
\defaultfontfeatures{Ligatures=TeX}
\setmainfont{PlusJakartaSans-Medium.ttf}[Path=../../../fonts/,BoldFont=PlusJakartaSans-Bold.ttf]
% Symboles, API et emojis absents de la police principale : repli sur DejaVu Sans (caractères actifs)
\newfontface\symfont{DejaVuSans.ttf}[Path=../../../fonts/]
\makeatletter
\newcommand\symactive[1]{\catcode#1=13 \begingroup\lccode`\~=#1 \lowercase{\endgroup\def~}{{\symfont\char#1}}}
\newcommand\symblank[1]{\catcode#1=13 \begingroup\lccode`\~=#1 \lowercase{\endgroup\def~}{}}
\newcommand\symhyph[1]{\catcode#1=13 \begingroup\lccode`\~=#1 \lowercase{\endgroup\def~}{\mbox{-}}}
\newcount\symc
\newcommand\symrange[2]{\symc=#1 \@whilenum\symc<#2 \do{\expandafter\symactive\expandafter{\the\symc}\advance\symc by 1 }}
\newcommand\symblankrange[2]{\symc=#1 \@whilenum\symc<#2 \do{\expandafter\symblank\expandafter{\the\symc}\advance\symc by 1 }}
\makeatother
\symrange{"0250}{"02FF}\symactive{"2713}\symactive{"2714}\symactive{"2717}\symactive{"2718}\symactive{"2610}\symactive{"2611}\symactive{"2612}
\symhyph{"2011}\symhyph{"2E1A}\symblankrange{"1F300}{"1FAFF}\symblank{"FE0F}
% Maths en sans-serif (Fira Math) avec chiffres et lettres droites de Plus
% Jakarta, pour que les formules se fondent dans le texte.
\usepackage[math-style=ISO,bold-style=ISO]{unicode-math}
\setmathfont{FiraMath-Regular.otf}[Path=../../../fonts/]
\setmathfont{PlusJakartaSans-Medium.ttf}[Path=../../../fonts/,range={up/num,up/{Latin,latin},\mathup}]
\usepackage{icomma} % virgule décimale sans espace en mode math
% Caractères mathématiques Unicode absents des polices de texte (Plus Jakarta,
% Archivo Black) : rendus par leur équivalent mathématique, en texte comme en
% formule — sinon ils s'affichent en carré vide (ex. « Arithmétique dans ℤ »).
\usepackage{newunicodechar}
\newunicodechar{ℝ}{\ensuremath{\mathbb{R}}}
\newunicodechar{ℤ}{\ensuremath{\mathbb{Z}}}
\newunicodechar{ℂ}{\ensuremath{\mathbb{C}}}
\newunicodechar{ℕ}{\ensuremath{\mathbb{N}}}
\newunicodechar{ℚ}{\ensuremath{\mathbb{Q}}}
\newunicodechar{𝔻}{\ensuremath{\mathbb{D}}}
\newunicodechar{α}{\ensuremath{\alpha}}
\newunicodechar{β}{\ensuremath{\beta}}
\newunicodechar{γ}{\ensuremath{\gamma}}
\newunicodechar{δ}{\ensuremath{\delta}}
\newunicodechar{ε}{\ensuremath{\varepsilon}}
\newunicodechar{θ}{\ensuremath{\theta}}
\newunicodechar{λ}{\ensuremath{\lambda}}
\newunicodechar{μ}{\ensuremath{\mu}}
\newunicodechar{σ}{\ensuremath{\sigma}}
\newunicodechar{τ}{\ensuremath{\tau}}
\newunicodechar{φ}{\ensuremath{\varphi}}
\newunicodechar{ψ}{\ensuremath{\psi}}
\newunicodechar{ω}{\ensuremath{\omega}}
\newunicodechar{Σ}{\ensuremath{\Sigma}}
% majuscules produites par \MakeUppercase dans les titres (α-aminés -> Α)
\newunicodechar{Α}{\ensuremath{\alpha}}
\newunicodechar{Β}{\ensuremath{\beta}}
\newunicodechar{Γ}{\ensuremath{\Gamma}}
\newunicodechar{Δ}{\ensuremath{\Delta}}
\newunicodechar{Ε}{\ensuremath{\varepsilon}}
\newunicodechar{Θ}{\ensuremath{\Theta}}
\newunicodechar{Λ}{\ensuremath{\Lambda}}
\newunicodechar{Μ}{\ensuremath{\mu}}
\newunicodechar{Τ}{\ensuremath{\tau}}
\newunicodechar{Φ}{\ensuremath{\Phi}}
\newunicodechar{Ψ}{\ensuremath{\Psi}}
\newunicodechar{Ω}{\ensuremath{\Omega}}
\newunicodechar{↦}{\ensuremath{\mapsto}}
\newunicodechar{∈}{\ensuremath{\in}}
\newunicodechar{∉}{\ensuremath{\notin}}
\newunicodechar{∧}{\ensuremath{\wedge}}
\newunicodechar{⇔}{\ensuremath{\Leftrightarrow}}
\newunicodechar{⟺}{\ensuremath{\Longleftrightarrow}}
\newunicodechar{⇒}{\ensuremath{\Rightarrow}}
\newunicodechar{⟹}{\ensuremath{\Longrightarrow}}
\newunicodechar{∘}{\ensuremath{\circ}}
\newunicodechar{∪}{\ensuremath{\cup}}
\newunicodechar{∩}{\ensuremath{\cap}}
\newunicodechar{≡}{\ensuremath{\equiv}}
\newunicodechar{⊂}{\ensuremath{\subset}}
\newunicodechar{⊥}{\ensuremath{\perp}}
\newunicodechar{∃}{\ensuremath{\exists}}
\newunicodechar{∀}{\ensuremath{\forall}}
\newunicodechar{∛}{\ensuremath{\sqrt[3]{\,}}}
\newunicodechar{∜}{\ensuremath{\sqrt[4]{\,}}}
\newunicodechar{⃗}{\ensuremath{{}^{\to}}}
\newunicodechar{ⁿ}{\ifmmode^{n}\else\textsuperscript{\ensuremath{n}}\fi}
\newunicodechar{ˣ}{\ifmmode^{x}\else\textsuperscript{\ensuremath{x}}\fi}
\newunicodechar{ᵇ}{\ifmmode^{b}\else\textsuperscript{\ensuremath{b}}\fi}
\newunicodechar{ʳ}{\ifmmode^{r}\else\textsuperscript{\ensuremath{r}}\fi}
\newunicodechar{ᵘ}{\ifmmode^{u}\else\textsuperscript{\ensuremath{u}}\fi}
\newunicodechar{ʸ}{\ifmmode^{y}\else\textsuperscript{\ensuremath{y}}\fi}
\newunicodechar{ᵃ}{\ifmmode^{a}\else\textsuperscript{\ensuremath{a}}\fi}
\newunicodechar{ᵉ}{\ifmmode^{e}\else\textsuperscript{\ensuremath{e}}\fi}
\newunicodechar{ᵏ}{\ifmmode^{k}\else\textsuperscript{\ensuremath{k}}\fi}
\newunicodechar{ᵖ}{\ifmmode^{p}\else\textsuperscript{\ensuremath{p}}\fi}
\newunicodechar{ᴺ}{\ifmmode^{N}\else\textsuperscript{\ensuremath{N}}\fi}
\newunicodechar{ᶜ}{\ifmmode^{c}\else\textsuperscript{\ensuremath{c}}\fi}
\newunicodechar{ʰ}{\ifmmode^{h}\else\textsuperscript{\ensuremath{h}}\fi}
\newunicodechar{ᵠ}{\ifmmode^{q}\else\textsuperscript{\ensuremath{q}}\fi}
\newunicodechar{ₙ}{\ifmmode_{n}\else\textsubscript{\ensuremath{n}}\fi}
\newunicodechar{ₐ}{\ifmmode_{a}\else\textsubscript{\ensuremath{a}}\fi}
\newunicodechar{ᵢ}{\ifmmode_{i}\else\textsubscript{\ensuremath{i}}\fi}
\newunicodechar{ᵣ}{\ifmmode_{r}\else\textsubscript{\ensuremath{r}}\fi}
\newunicodechar{ₖ}{\ifmmode_{k}\else\textsubscript{\ensuremath{k}}\fi}
\newunicodechar{ᵦ}{\ifmmode_{\beta}\else\textsubscript{\ensuremath{\beta}}\fi}
\newunicodechar{ₓ}{\ifmmode_{x}\else\textsubscript{\ensuremath{x}}\fi}
\newfontfamily\popdisplay{ArchivoBlack-Regular.ttf}[Path=../../../fonts/]
\newfontfamily\poplogo{SpaceGrotesk-Bold.ttf}[Path=../../../fonts/]
\newfontfamily\popsemi{PlusJakartaSans-SemiBold.ttf}[Path=../../../fonts/]
\newfontfamily\popxbold{PlusJakartaSans-ExtraBold.ttf}[Path=../../../fonts/]
\newfontfamily\popxboldls{PlusJakartaSans-ExtraBold.ttf}[Path=../../../fonts/,LetterSpace=3.0]

\setlist[enumerate]{leftmargin=1.7em,itemsep=5pt,topsep=4pt}
\setlist[itemize]{leftmargin=1.7em,itemsep=4pt,topsep=4pt,label={\color{poporange}\textbullet}}
\setlength{\parindent}{0pt}
\setlength{\parskip}{5pt}
\renewcommand{\arraystretch}{1.25}

% pgfplots : style Pop par défaut
\pgfplotsset{
  every axis/.append style={axis line style={popink,line width=1pt},tick style={popink},
    grid style={popline,line width=0.5pt},label style={font=\small},tick label style={font=\footnotesize}},
  popcurve/.style={poporange,line width=1.8pt,no markers,smooth},
}
% Même style disponible dans un \draw TikZ simple (hors pgfplots)
\tikzset{popcurve/.style={poporange,line width=1.8pt}}
\tikzset{popgrid/.style={popline,very thin}}
\colorlet{popcurve}{poporange} % le nom du style est parfois utilisé comme couleur

% ============================================================
% Pill / squiggle
% ============================================================
\newcommand{\poppill}[3]{%
  \tikz[baseline=-0.55ex]\node[draw=popink,line width=1.2pt,rounded corners=2.6pt,%
    fill=#2,inner xsep=6.5pt,inner ysep=3pt]{%
    \popxboldls\fontsize{7}{7}\selectfont\color{#3}\MakeUppercase{#1}};%
}
\newcommand{\popsquiggle}[1]{%
  \tikz[baseline=-0.4ex]{
    \draw[#1,line width=2.6pt,line cap=round]
      plot [smooth] coordinates {(0,0.09) (0.34,0.01) (0.68,0.21) (1.02,0.01) (1.36,0.10)};
  }%
}
% Mot-clé surligné (marqueur orange sous le texte)
\newcommand{\cle}[1]{{\popxbold\color{popdark}#1}}

% ============================================================
% En-tête / pied
% ============================================================
\pagestyle{fancy}
\fancyhf{}
\renewcommand{\headrulewidth}{0pt}
\renewcommand{\footrulewidth}{0pt}
\lhead{\raisebox{-0.15ex}{\poplogo\fontsize{13}{13}\selectfont\color{popink}OnBuch\color{poporange}+}}
\rhead{\poppill{\DOCMATIERE\ \textbullet\ \DOCNIVEAU\ \textbullet\ Leçon \DOCLECON}{white}{popink}}
\lfoot{\popxbold\fontsize{7.6}{7.6}\selectfont\color{popmuted}\MakeUppercase{Cours signé OnBuch+}\ \textbullet\ {\popsemi\DOCTITRE}}
\rfoot{\tikz[baseline=-0.6ex]\node[rounded corners=8.5pt,fill=popink,minimum height=17pt,minimum width=17pt,inner xsep=5pt,inner sep=0]{\popxbold\fontsize{8}{8}\selectfont\color{popcream}\thepage\,/\,\pageref*{LastPage}};}

% ============================================================
% Titres
% ============================================================
\newcommand{\chaptitre}[2]{%
  \begin{tcolorbox}[enhanced,colback=poporange,colframe=popink,boxrule=2.6pt,arc=11pt,
      drop shadow=popink,left=15pt,right=15pt,top=12pt,bottom=11pt,before skip=3pt,after skip=18pt]
    \poppill{#2}{popink}{popcream}\par\vspace{9pt}
    {\raggedright\hyphenpenalty=10000\exhyphenpenalty=10000\popdisplay\fontsize{22}{24}\selectfont\color{popcream}\MakeUppercase{#1}\par}
  \end{tcolorbox}
}
% Section numérotée : pastille orange avec le numéro + titre Archivo
\newcounter{popsec}
\newcommand{\coursec}[1]{%
  \stepcounter{popsec}\setcounter{popsub}{0}%
  \par\vspace{16pt}\noindent
  \begin{minipage}{\linewidth}
  \tikz[baseline=-0.7ex]\node[draw=popink,line width=1.6pt,fill=poporange,rounded corners=4pt,minimum size=22pt,inner sep=0]{\popdisplay\fontsize{12}{12}\selectfont\color{popcream}\thepopsec};%
  \hspace{8pt}{\popdisplay\fontsize{15}{17}\selectfont\color{popink}\MakeUppercase{#1}}\par
  \vspace{4pt}\noindent{\color{poporange}\rule{36pt}{3.6pt}}
  \end{minipage}\par\nopagebreak\vspace{8pt}
}
\newcounter{popsub}
\newcommand{\courssub}[1]{%
  \stepcounter{popsub}%
  \par\vspace{9pt}\noindent{\popxbold\fontsize{11.5}{13}\selectfont\color{popink}{\color{poporange}\thepopsec.\thepopsub}\hspace{6pt}#1}\par\nopagebreak\vspace{4pt}
}

% ============================================================
% Boîtes pédagogiques — même recette : fond blanc, bordure épaisse colorée,
% ombre dure encre, badge plein. Titre optionnel [#1].
% ============================================================
\tcbset{popbase/.style={breakable,enhanced,colback=white,boxrule=2.3pt,arc=9pt,
  shadow={3pt}{-3pt}{0pt}{popink},left=14pt,right=14pt,top=11pt,bottom=11pt,before skip=11pt,after skip=15pt,
  fontupper=\popsemi\fontsize{10.6}{14.5}\selectfont\color{popink},
  fontlower=\popsemi\fontsize{10.6}{14.5}\selectfont\color{popink}}}
% Coche dessinée (le glyphe ✓ n'existe pas dans les polices du document)
\newcommand{\popcheck}[1]{\tikz[baseline=-0.5ex]\draw[#1,line width=1.6pt,line cap=round,line join=round](-0.13,0.0)--(-0.04,-0.09)--(0.13,0.1);}
\providecommand{\coche}{\popcheck{popgreen}}
\providecommand{\resultat}[1]{\par\smallskip\noindent\fcolorbox{popgreen}{popgreenL}{\parbox{\dimexpr\linewidth-2\fboxsep-2\fboxrule\relax}{\textbf{Résultat.} #1}}\par\smallskip}
\newcommand{\popboxhead}[3]{\poppill{#1}{#2}{white}\ifx\relax#3\relax\else\hspace{6pt}{\popxbold\fontsize{10.6}{12}\selectfont\color{popink}#3}\fi\par\vspace{7pt}}

\newtcolorbox{aretenir}[1][]{popbase,colframe=popgreen,before upper={\popboxhead{À retenir}{popgreen}{#1}}}
% Texte en anglais (document de lecture, dialogue, extrait) : cadre
% d'étude. Un titre optionnel (source, auteur) s'affiche dans l'en-tête.
\newtcolorbox{textebox}[1][]{popbase,colframe=popblue,before upper={\popboxhead{Text}{popblue}{#1}\begin{otherlanguage}{english}},after upper={\end{otherlanguage}}}
% Marques ✓ / ✗ (forme correcte / fautive) souvent utilisées par les modèles
\providecommand{\cross}{\textcolor{poppink}{\ensuremath{\times}}}
\providecommand{\xmark}{\cross}\providecommand{\crossmark}{\cross}\providecommand{\cmark}{\ensuremath{\checkmark}}
% Ligne de réponse / justification à compléter (inventée par les modèles)
\providecommand{\justification}[1]{\par\noindent\textit{Justification~:} \dotfill\par}
% \textipa (package tipa non chargé) : la police ne couvre pas l'API -> simple passage
\providecommand{\textipa}[1]{#1}
% Soulignement (ulem non chargé) : \uline, \ul
\providecommand{\uline}[1]{\underline{#1}}\providecommand{\ul}[1]{\underline{#1}}
% \glossaire{\item ...} inventée par les modèles : liste de glossaire
\providecommand{\glossaire}[1]{\begin{itemize}#1\end{itemize}}
% \alert (beamer) inventée par les modèles : texte mis en évidence
\providecommand{\alert}[1]{\textbf{\textcolor{poppink}{#1}}}
% variantes ASCII d'ordinaux écrites en macro
\providecommand{\iere}{\textsuperscript{re}}\providecommand{\ieme}{\textsuperscript{e}}\providecommand{\ere}{\textsuperscript{re}}\providecommand{\eme}{\textsuperscript{e}}\providecommand{\er}{\textsuperscript{er}}\providecommand{\ie}{\textsuperscript{e}}
% Ordinaux français écrits en macro par les modèles : 1\ière, 3\ième
\newcommand{\ière}{\textsuperscript{re}}\newcommand{\ième}{\textsuperscript{e}}
% \exemple (inventée par les modèles) : étiquette « Exemple : »
\providecommand{\exemple}{\textbf{Exemple~:}\ }
% \square utilisable aussi hors mode math (cases à cocher)
\AtBeginDocument{\let\squareMath\square\renewcommand{\square}{\ensuremath{\squareMath}}}
% Mot ou phrase en anglais à l'intérieur d'un paragraphe français
\newcommand{\eng}[1]{\ifmmode\text{\textit{#1}}\else\begingroup\selectlanguage{english}\itshape #1\endgroup\fi}
\let\esp\eng % alias toléré
% flèches d'intonation inventées par les modèles
\providecommand{\textsearrow}{}\renewcommand{\textsearrow}{\ensuremath{\searrow}}\providecommand{\textnearrow}{}\renewcommand{\textnearrow}{\ensuremath{\nearrow}}
% \fr{..} : mot français (inventé par les modèles)
\providecommand{\fr}[1]{\textit{#1}}
\newtcolorbox{exemplebox}[1][]{popbase,colframe=poporange,before upper={\popboxhead{Exemple}{poporange}{#1}}}
\newtcolorbox{definition}[1][]{popbase,colframe=popblue,before upper={\popboxhead{Définition}{popblue}{#1}}}
\newtcolorbox{propriete}[1][]{popbase,colframe=poppurple,before upper={\popboxhead{Propriété}{poppurple}{#1}}}
\newtcolorbox{methode}[1][]{popbase,colframe=popgold,before upper={\popboxhead{Méthode}{popgold}{#1}}}
\newtcolorbox{attention}[1][]{popbase,colframe=poppink,before upper={\popboxhead{Attention}{poppink}{#1}}}
\newtcolorbox{experience}[1][]{popbase,colframe=popblue,colback=popblueL,before upper={\popboxhead{Expérience}{popblue}{#1}}}
% « Le savais-tu ? » : carte violette pleine (contexte, culture, Cameroun)
\newtcolorbox{savaistu}[1][]{popbase,colback=poppurple,colframe=popink,
  fontupper=\popsemi\fontsize{10.4}{14.2}\selectfont\color{white},
  before upper={\poppill{Le savais-tu\,?}{white}{poppurple}\ifx\relax#1\relax\else\hspace{6pt}{\popxbold\color{white}#1}\fi\par\vspace{7pt}}}
% Exercice résolu : énoncé en haut, \tcblower puis la solution sur fond crème
\newcounter{popexo}\newcounter{popexr}
\newtcolorbox{exoresolu}[1][]{popbase,colframe=poporange,colbacklower=popcream,
  segmentation style={popink,line width=1.2pt,dashed},
  before upper={\stepcounter{popexr}\popboxhead{Exercice résolu \thepopexr}{poporange}{#1}},
  before lower={\poppill{Solution}{popink}{popcream}\par\vspace{6pt}}}
% Exercice (non résolu en place) — la correction est dans la partie Corrigés
\newtcolorbox{exercice}[1][]{popbase,colframe=popink,
  before upper={\stepcounter{popexo}\popboxhead{Exercice \thepopexo}{popink}{#1}}}
% Corrigé
\newtcolorbox{corrige}[1][]{popbase,colframe=popgreen,colback=popgreenL,
  before upper={\popboxhead{Corrigé}{popgreen}{#1}}}
% Objectifs / prérequis / bilan
\newtcolorbox{objectifs}{popbase,colframe=popink,colback=poporangeL,
  before upper={\popboxhead{Objectifs de la leçon}{popink}{}}}
\newtcolorbox{prerequis}{popbase,colframe=popink,colback=popblueL,
  before upper={\popboxhead{Prérequis}{popblue}{}}}
\newtcolorbox{bilan}{popbase,colframe=popink,colback=white,boxrule=2.8pt,
  before upper={\popboxhead{Fiche bilan}{poporange}{L'essentiel à connaître pour l'examen}}}
% Figure encadrée avec légende
\newtcolorbox{popfigure}[1][]{enhanced,breakable=false,colback=white,colframe=popline,boxrule=1.4pt,arc=8pt,
  left=10pt,right=10pt,top=10pt,bottom=8pt,before skip=10pt,after skip=12pt,halign=center,
  lower separated=false,halign lower=center,
  fontlower=\popsemi\fontsize{9}{11.5}\selectfont\color{popmuted},
  #1}
\newcommand{\legende}[1]{\par\vspace{6pt}{\popsemi\fontsize{9}{11.5}\selectfont\color{popmuted}#1}}

% ============================================================
% Signature OnBuch+
% ============================================================
\newcommand{\poplogoinline}[1]{{\poplogo\fontsize{#1}{#1}\selectfont\color{popink}OnBuch\color{poporange}+}}
\newcommand{\signatureonbuch}{%
  \begin{tcolorbox}[enhanced,colback=popink,colframe=popink,boxrule=2.2pt,arc=10pt,
      drop shadow=poporange,left=14pt,right=14pt,top=10pt,bottom=10pt,before skip=0pt,after skip=0pt]
    \tikz[baseline=-0.7ex]\node[circle,fill=popgreen,draw=popcream,line width=1.3pt,minimum size=22pt,inner sep=0]{\popcheck{white}};%
    \hspace{9pt}\begin{minipage}[c]{0.62\linewidth}
      {\popxbold\fontsize{12}{13}\selectfont\color{popcream}Signé\ \textbullet\ {\poplogo OnBuch\color{poporange}+}}\par\vspace{2pt}
      {\raggedright\popsemi\fontsize{8}{10.5}\selectfont\color{popline}\MakeUppercase{Équipe pédagogique OnBuch+}\\\MakeUppercase{Conforme au programme officiel MINESEC}\par}
    \end{minipage}\hfill
    {\poplogo\fontsize{22}{22}\selectfont\color{popcream}OnBuch\color{poporange}+}
  \end{tcolorbox}%
}

% ============================================================
% Page de garde
% #1 titre · #2 matière · #3 niveau · #4 module · #5 n° leçon · #6 durée ·
% #7 description / objectifs (texte ou itemize)
% ============================================================
\newcommand{\pagegarde}[7]{%
  \thispagestyle{empty}
  \newgeometry{a4paper,margin=1.7cm,top=1.6cm,bottom=1.4cm}%
  \begin{tikzpicture}[remember picture,overlay]
    \fill[poporange!18] ([xshift=-3.2cm,yshift=-3.6cm]current page.north east) circle (4.6cm);
    \fill[poppurple!14] ([xshift=2.4cm,yshift=7.6cm]current page.south west) circle (3.8cm);
  \end{tikzpicture}%
  {\poplogo\fontsize{30}{30}\selectfont\color{popink}OnBuch\color{poporange}+}\hfill
  \raisebox{0.6ex}{\poppill{Cours signé \textbullet\ Premium}{popink}{popcream}}\par
  \vspace{1pt}\popsquiggle{poppurple}\par
  \vspace{8pt}
  \begin{tcolorbox}[enhanced,colback=poporange,colframe=popink,boxrule=2.8pt,arc=13pt,
      drop shadow=popink,left=18pt,right=18pt,top=12pt,bottom=13pt,after skip=10pt]
    \poppill{#2 \textbullet\ #3}{popink}{popcream}\hspace{5pt}\poppill{#4}{white}{popink}\par\vspace{8pt}
    {\popdisplay\fontsize{14}{14}\selectfont\color{popink}LEÇON #5}\par\vspace{4pt}
    {\raggedright\hyphenpenalty=10000\exhyphenpenalty=10000\popdisplay\fontsize{25}{27}\selectfont\color{popcream}\MakeUppercase{#1}\par}
    \vspace{8pt}
    {\popxbold\fontsize{9.5}{9.5}\selectfont\color{popink}\MakeUppercase{Durée conseillée : #6}}
  \end{tcolorbox}
  \begin{tcolorbox}[enhanced,colback=white,colframe=popink,boxrule=2.2pt,arc=10pt,
      drop shadow=popink,left=16pt,right=16pt,top=10pt,bottom=10pt,after skip=8pt]
    \poppill{Dans cette leçon}{poppurple}{white}\par\vspace{5pt}
    \popsemi\fontsize{9.3}{12.6}\selectfont\color{popink}#7
  \end{tcolorbox}
  \noindent\begin{minipage}[t]{0.487\linewidth}
    \begin{tcolorbox}[enhanced,colback=popgreen,colframe=popink,boxrule=2.2pt,arc=10pt,
        drop shadow=popink,left=11pt,right=11pt,top=10pt,bottom=10pt,height=3.1cm,valign=center]
      \tcbox[colback=white,colframe=popink,boxrule=1.3pt,arc=5pt,boxsep=0pt,nobeforeafter,
        left=3pt,right=3pt,top=3pt,bottom=3pt,valign=center]{\qrcode[height=1.9cm]{https://play.google.com/store/apps/details?id=com.luvvix.onbuch&hl=fr}}%
      \hspace{8pt}\begin{minipage}[c]{0.5\linewidth}
        {\popdisplay\fontsize{11}{12.5}\selectfont\color{white}\MakeUppercase{Télécharge OnBuch}}\par\vspace{4pt}
        {\raggedright\popsemi\fontsize{8.4}{11}\selectfont\color{white}Gratuit \textbullet\ Google Play\\Tuteur IA, annales, concours, résultats\par}
      \end{minipage}
    \end{tcolorbox}
  \end{minipage}\hfill
  \begin{minipage}[t]{0.487\linewidth}
    \begin{tcolorbox}[enhanced,colback=poppurple,colframe=popink,boxrule=2.2pt,arc=10pt,
        drop shadow=popink,left=11pt,right=11pt,top=10pt,bottom=10pt,height=3.1cm,valign=center]
      \tikz[baseline=-0.7ex]\node[circle,fill=white,draw=popink,line width=1.3pt,minimum size=1.6cm,inner sep=0]{\popdisplay\fontsize{24}{24}\selectfont\color{poppurple}+};%
      \hspace{8pt}\begin{minipage}[c]{0.6\linewidth}
        {\popdisplay\fontsize{11}{12.5}\selectfont\color{white}\MakeUppercase{Passe à OnBuch+}}\par\vspace{4pt}
        {\raggedright\popsemi\fontsize{8.4}{11}\selectfont\color{white}Tous les cours signés, Tuteur IA illimité, fiches PDF — onglet OnBuch+ de l'app\par}
      \end{minipage}
    \end{tcolorbox}
  \end{minipage}
  \vspace{8pt}
  \popcontenu
  \vfill
  \signatureonbuch
  \restoregeometry
  \newpage
}

% Rangée « Ce cours contient » (page de garde)
\newcommand{\poptile}[3]{% #1 couleur #2 gros texte #3 légende
  \begin{tcolorbox}[enhanced,colback=#1,colframe=popink,boxrule=2pt,arc=9pt,shadow={2.5pt}{-2.5pt}{0pt}{popink},
      left=6pt,right=6pt,top=9pt,bottom=9pt,halign=center,valign=center,height=1.75cm,width=\linewidth,nobeforeafter]
    {\popdisplay\fontsize{12.5}{13}\selectfont\color{white}\MakeUppercase{#2}}\par\vspace{3pt}
    {\centering\popsemi\fontsize{7.8}{9.5}\selectfont\color{white}#3\par}
  \end{tcolorbox}}
\newcommand{\popcontenu}{%
  \par\noindent{\popxboldls\fontsize{7.5}{7.5}\selectfont\color{popmuted}CE COURS SIGNÉ CONTIENT}\par\vspace{7pt}
  \noindent\begin{minipage}[t]{0.235\linewidth}\poptile{popblue}{Cours}{complet, illustré, pas à pas}\end{minipage}\hfill
  \begin{minipage}[t]{0.235\linewidth}\poptile{poporange}{Méthodes}{et exercices résolus}\end{minipage}\hfill
  \begin{minipage}[t]{0.235\linewidth}\poptile{poppink}{Exercices}{type examen + corrigés}\end{minipage}\hfill
  \begin{minipage}[t]{0.235\linewidth}\poptile{popgreen}{Fiche bilan}{l'essentiel à retenir}\end{minipage}\par}

% Sommaire
\newtcolorbox{sommaire}{enhanced,colback=white,colframe=popink,boxrule=2.2pt,arc=10pt,
  drop shadow=popink,left=18pt,right=18pt,top=13pt,bottom=13pt,before skip=6pt,after skip=20pt,
  before upper={\poppill{Sommaire}{poppurple}{white}\par\vspace{9pt}\popsemi\fontsize{10.6}{14.5}\selectfont\color{popink}}}

% Signature de fin de cours — poussée tout en bas de la dernière page
\newcommand{\findecours}{%
  \par\vfill\nopagebreak
  \begin{center}\popsquiggle{poporange}\end{center}\vspace{4pt}
  \signatureonbuch
}
\AtBeginDocument{\def\llbracket{\mathopen{[\mkern-3.5mu[}}\def\rrbracket{\mathclose{]\mkern-3.5mu]}}} % intervalles d'entiers (glyphe ⟦ absent de la police)
% Glyphes absents de Fira Math : redéfinis à partir de glyphes disponibles
\AtBeginDocument{%
  \def\setminus{\mathbin{\backslash}}\let\smallsetminus\setminus
  \def\vdots{\mathord{\vbox{\baselineskip3.2pt\lineskiplimit0pt\kern4pt\hbox{.}\hbox{.}\hbox{.}}}}%
  \def\ddots{\mathinner{\mkern1mu\raise6.5pt\hbox{.}\mkern2mu\raise3.6pt\hbox{.}\mkern2mu\raise0.7pt\hbox{.}\mkern1mu}}%
  \def\triangle{\mathord{\scalebox{0.9}{$\symup{\Delta}$}}}%
  \def\nabla{\mathord{\raisebox{\depth}{\scalebox{1}[-1]{$\symup{\Delta}$}}}}%
  \def\checkmark{\popcheck{popdark}}%
  \def\star{\mathbin{\ast}}\def\bigstar{\mathord{\ast}}%
  \def\diamond{\mathbin{\scalebox{0.6}{$\Diamond$}}}%
  \def\bigcup{\mathop{\mathchoice{\raisebox{-0.3em}{\scalebox{1.7}{$\cup$}}}{\scalebox{1.25}{$\cup$}}{\cup}{\cup}}\displaylimits}%
  \def\bigcap{\mathop{\mathchoice{\raisebox{-0.3em}{\scalebox{1.7}{$\cap$}}}{\scalebox{1.25}{$\cap$}}{\cap}{\cap}}\displaylimits}%
  \def\lesseqqgtr{\mathrel{\lessgtr}}\def\bigoplus{\mathop{\oplus}\displaylimits}%
}
\providecommand{\ssi}{si et seulement si} % abréviation fréquente
\AtBeginDocument{\providecommand{\wideparen}{\overparen}} % arc de cercle

%% FIGFIX-BEGIN v5 — correctifs globaux des figures (docs/onbuch-generation/tools/figfix)
\usepackage{adjustbox}\usepackage{etoolbox}\tcbuselibrary{raster}
% toute figure TikZ/pgfplots plus large que la ligne est réduite à la largeur disponible
\BeforeBeginEnvironment{tikzpicture}{\begin{adjustbox}{max width=\linewidth}}
\AfterEndEnvironment{tikzpicture}{\end{adjustbox}}
% légendes pgfplots lisibles par-dessus les courbes
\pgfplotsset{every axis/.append style={legend style={fill=white,fill opacity=0.92,text opacity=1,draw=black!15,inner sep=3pt}}}
% étiquettes d'arêtes (syntaxe edge["..."]) avec fond blanc
\tikzset{every edge quotes/.append style={fill=white,inner sep=1.2pt}}
%% FIGFIX-END
\begin{document}

\pagegarde{Lesson 16 — Listening: Texts about requesting for services}{Anglais}{1ère A}{Chapter 2 : Economic Life and Occupations}{ENA17}{2 h}{Leçon de compréhension orale centrée sur les demandes de services (banque, hôtel, téléphonie, réparations). L'élève apprend à écouter avec stratégie — repérer le thème, les mots-clés, les chiffres et les noms propres — et à réemployer le lexique et les formules de politesse dans ses propres demandes.
\vspace{4pt}
\begin{multicols}{2}\small
\begin{itemize}
\item À la fin de cette leçon, je sais identifier le thème général d'un enregistrement sur les services dès la première écoute.
\item Je sais repérer les mots-clés, les chiffres (prix, numéros, dates) et les noms propres dans un dialogue.
\item Je sais répondre à des questions de compréhension de types variés (vrai/faux, QCM, réponses courtes).
\item Je connais le lexique essentiel des services : banque, hôtel, téléphonie, réparations.
\item Je sais utiliser les formules polies de demande : 'Could you...?', 'I'd like...', 'I was wondering if...'.
\item Je sais distinguer une demande acceptée d'une demande refusée dans un dialogue.
\end{itemize}\end{multicols}}

\begin{sommaire}
\begin{multicols}{2}
\begin{enumerate}
\item Mise en route et stratégies d'écoute
\item Script d'écoute : At the service desk
\item Compréhension : questions variées
\item Lexique thématique : requesting services
\item Production guidée : je demande un service
\item Activité d'intégration
\item Méthodes et exercices résolus
\item Exercices
\item Corrigés des exercices
\item Fiche bilan
\end{enumerate}
\end{multicols}
\end{sommaire}

\begin{objectifs}
\begin{itemize}
\item À la fin de cette leçon, je sais identifier le thème général d'un enregistrement sur les services dès la première écoute.
\item Je sais repérer les mots-clés, les chiffres (prix, numéros, dates) et les noms propres dans un dialogue.
\item Je sais répondre à des questions de compréhension de types variés (vrai/faux, QCM, réponses courtes).
\item Je connais le lexique essentiel des services : banque, hôtel, téléphonie, réparations.
\item Je sais utiliser les formules polies de demande : 'Could you...?', 'I'd like...', 'I was wondering if...'.
\item Je sais distinguer une demande acceptée d'une demande refusée dans un dialogue.
\item Je sais rédiger de courtes phrases réemployant le lexique du thème.
\item Je sais prendre des notes efficaces pendant l'écoute (tableau, abréviations).
\end{itemize}
\end{objectifs}

\begin{prerequis}
\begin{itemize}
\item Les auxiliaires modaux de base : can, could, would (Seconde).
\item Le présent simple et le présent continu pour décrire des situations.
\item Les nombres, les prix et les heures en anglais.
\item Le vocabulaire général des métiers et des lieux (Chapter 2, leçons précédentes).
\item Les questions en wh- (what, where, when, how much).
\end{itemize}
\end{prerequis}

\newpage

\coursec{Lesson 16 — Listening: Texts about Requesting Services}

\courssub{Mise en route et stratégies d'écoute}

Imagine : tu es à l'agence MTN de Buea pour réclamer un forfait internet défectueux, ou à la réception d'un hôtel à Limbe pour demander une chambre. En anglais, demander un service exige des formules précises et polies : \eng{Could you...?}, \eng{I'd like to...}, \eng{I was wondering if...}. Aujourd'hui, tu vas écouter des conversations authentiques entre clients et employés — à la banque, à l'hôtel, au garage, à l'agence de téléphonie — et apprendre à repérer rapidement les informations essentielles : qui demande quoi, où, quand, et à quel prix.

\textbf{Problématique :} comment comprendre efficacement un dialogue de demande de service en anglais, sans tout comprendre mot à mot, et comment noter les informations utiles pour répondre aux questions de compréhension ?

\courssub{Brainstorming : quels services demande-t-on au quotidien ?}

\begin{experience}[Warm-up : Services in My Town]
En deux minutes, liste avec ton voisin tous les services que tu as demandés (ou que ta famille a demandés) ce mois-ci. Puis classe-les par lieu. Exemples attendus :
\begin{itemize}
\item \eng{At the bank}: opening an account, withdrawing money, exchanging currency.
\item \eng{At the hotel}: booking a room, asking for an extra blanket, ordering breakfast.
\item \eng{At the hospital / clinic}: making an appointment, asking for a prescription, seeing a specialist.
\item \eng{At the phone agency} (MTN, Orange): buying airtime, reporting a faulty bundle, registering a SIM card.
\item \eng{At the garage}: repairing a flat tyre, servicing a car, changing the oil.
\item \eng{At the hairdresser's / barber's}: asking for a haircut, braiding hair.
\item \eng{At the post office}: sending a parcel, buying stamps.
\end{itemize}
Partage tes idées avec la classe. Le professeur construit un tableau au tableau : \textbf{Place} | \textbf{Service} | \textbf{Key Vocabulary}.
\end{experience}

\begin{definition}[What is a Service?]
Un \cle{service} est une prestation qu'on demande à un professionnel, généralement contre paiement. En anglais : \eng{a service}, \eng{to request a service} (demander un service — plus formel que \eng{ask for}), \eng{a customer} (un client dans un commerce), \eng{a client} (un client d'un professionnel : avocat, banque, coiffeur), \eng{a guest} (un client d'hôtel ou un invité), \eng{a service provider} (un prestataire de services).
\end{definition}

\begin{center}
\begin{tabularx}{\linewidth}{|X|X|X|}
\hline
\rowcolor{popblueL} Français & English & Remarque / Prononciation \\
\hline
demander un service & to request a service / to ask for a service & \eng{request} + objet direct ; \eng{ask for} + objet \\
\hline
un client (commerce) & a customer & « koss-te-meur » \\
\hline
un client (hôtel / pro) & a guest / a client & \eng{guest} = hôtel ; \eng{client} = avocat, banque \\
\hline
un employé au guichet & a receptionist / a clerk / an agent / a teller (banque) & \eng{teller} = guichetier bancaire \\
\hline
un compte bancaire & a bank account & \eng{current account} (compte courant), \eng{savings account} (compte épargne) \\
\hline
réserver & to book & \textbf{Faux ami :} \eng{book} $\neq$ livre (\eng{book} n.) \\
\hline
une facture / l'addition & a bill (GB) / a check (US) & \eng{The bill, please.} \\
\hline
une réparation & a repair & verbe : \eng{to repair} / \eng{to fix} \\
\hline
un rendez-vous & an appointment & \textbf{Jamais} \eng{a rendezvous} (rencontre secrète/amoureuse) \\
\hline
un forfait (téléphonie) & a bundle / a plan / a package & \eng{data bundle}, \eng{monthly plan} \\
\hline
une carte SIM & a SIM card & \eng{to register a SIM card} \\
\hline
recharger (crédit) & to top up / to recharge & \eng{I'd like to top up my account.} \\
\hline
\end{tabularx}
\end{center}

\begin{attention}[Faux amis et pièges de vocabulaire fréquents]
\begin{itemize}
\item \eng{to book} = réserver. \eng{I'd like to book a double room.} (Je voudrais réserver une chambre double.) — Ne pas confondre avec \eng{a book} (un livre).
\item \eng{a bill} = une facture, l'addition (GB). US : \eng{a check}. \eng{The bill, please.}
\item \eng{to assist} = aider. \eng{Can I assist you?} (Puis-je vous aider ?) — Ne pas confondre avec « assister à » (\eng{attend} / \eng{watch}).
\item \eng{actually} = en fait, en réalité. \eng{Actually, I have an appointment.} — Ne pas confondre avec « actuellement » (\eng{currently} / \eng{at the moment}).
\item \eng{an appointment} = un rendez-vous (professionnel). \eng{a date} = un rendez-vous amoureux.
\item \eng{a subscription} = un abonnement. \eng{to subscribe} = s'abonner.
\end{itemize}
\end{attention}

\courssub{Point de grammaire : Formules de politesse pour demander un service (Polite Requests)}

Demander un service en anglais repose sur l'usage des \cle{modaux} et de structures \cle{indirectes} pour adoucir la demande. Plus la situation est formelle (banque, administration), plus la structure est complexe.

\begin{propriete}[Les degrés de politesse dans la demande]
\begin{center}
\begin{tabular}{|l|l|l|}
\hline
\rowcolor{popblueL} \textbf{Structure} & \textbf{Exemple} & \textbf{Niveau / Contexte} \\
\hline
\eng{Can I / Can you...?} & \eng{Can I have a receipt?} & Standard, neutre, courant. \\
\hline
\eng{Could I / Could you...?} & \eng{Could you check my balance, please?} & Plus poli, recommandé au guichet. \\
\hline
\eng{May I...?} & \eng{May I see your ID card?} & Formel, permission (souvent employé par l'employé). \\
\hline
\eng{I'd like to...} (\eng{I would like to}) & \eng{I'd like to open an account.} & Très courant, professionnel, clair. \\
\hline
\eng{I was wondering if I could...} & \eng{I was wondering if I could book a room for Friday.} & Très poli, hésitant, "au cas où". \\
\hline
\eng{Would you mind... -ing?} & \eng{Would you mind filling in this form?} & Poli, demande un effort à l'interlocuteur. \\
\hline
\end{tabular}
\end{center}

\textbf{Règles clés :}
\begin{enumerate}
\item \eng{Would like} est suivi de l'\textbf{infinitif} (\eng{to + BV}). Jamais \eng{I would like opening}.
\item \eng{Wondering if} introduit une \textbf{question indirecte} : l'ordre des mots est \textbf{affirmatif} (sujet + verbe). \eng{I was wondering if you \textbf{could help} me.} (Pas \eng{...if could you help}).
\item \eng{Mind} est suivi du \textbf{gérondif} (\eng{-ing}). \eng{Would you mind \textbf{waiting}?}
\item \eng{Can/Could/May} sont suivis de la \textbf{base verbale} (sans \eng{to}).
\end{enumerate}
\end{propriete}

\begin{exemplebox}[Exemples traduits]
\begin{itemize}
\item \eng{Could you tell me the price, please?} « Pourriez-vous me dire le prix, s'il vous plaît ? »
\item \eng{I'd like to report a faulty bundle.} « Je voudrais signaler un forfait défectueux. »
\item \eng{I was wondering if I could get a receipt.} « Je me demandais si je pouvais obtenir un reçu. »
\item \eng{Would you mind signing here?} « Voulez-vous bien signer ici ? / Est-ce que ça vous dérange de signer ici ? »
\end{itemize}
\end{exemplebox}

\begin{attention}[Erreurs typiques des francophones]
\begin{itemize}
\item \eng{$\times$ I would like to \textbf{opening} an account.} $\rightarrow$ \eng{\checkmark I would like to \textbf{open} an account.} (Infinitif après \eng{would like}).
\item \eng{$\times$ Could you \textbf{to help} me?} $\rightarrow$ \eng{\checkmark Could you \textbf{help} me?} (Base verbale après modal).
\item \eng{$\times$ I was wondering if \textbf{can I} have...} $\rightarrow$ \eng{\checkmark I was wondering if \textbf{I could} have...} (Ordre affirmatif + past modal).
\item \eng{$\times$ I want a room.} $\rightarrow$ \eng{\checkmark I'd like a room.} (\eng{I want} est trop direct, voire impoli dans un contexte de service).
\end{itemize}
\end{attention}

\courssub{Stratégies d'écoute : repérer l'essentiel}

\begin{methode}[Écouter efficacement en quatre étapes]
\begin{enumerate}
\item \textbf{Lis les questions AVANT l'écoute.} Souligne les mots interrogatifs (\eng{where, when, how much, who, what}) : ils te disent \textbf{quoi} chercher.
\item \textbf{Première écoute = le thème général (le \cle{gist}).} Ne note rien ou presque. Réponds mentalement : « Où sommes-nous ? Qui parle ? De quoi s'agit-il ? »
\item \textbf{Deuxième écoute = les détails.} Note les \textbf{chiffres} (prix, numéros, dates, heures), les \textbf{noms propres} (souvent épelés) et les \textbf{mots-clés} du thème.
\item \textbf{Troisième écoute / Relecture du script = vérification.} Contrôle chaque réponse, vérifie l'orthographe des noms et l'accord de tes phrases.
\end{enumerate}
\end{methode}

\courssub{Les quatre stratégies d'écoute en détail}

\textbf{Stratégie 1 : Identifier le thème général (Gist).}\par
À la première écoute, oublie les détails. Les formules d'ouverture trahissent presque toujours le lieu :
\begin{itemize}
\item \eng{Good morning, how can I help you?} / \eng{Next, please!} $\rightarrow$ Guichet, agence, banque, poste.
\item \eng{Reception, good evening. How may I assist you?} $\rightarrow$ Hôtel.
\item \eng{Your car is ready, sir. The repair cost is...} $\rightarrow$ Garage.
\item \eng{Welcome to MTN. How can I help you today?} $\rightarrow$ Agence de téléphonie.
\end{itemize}
Un seul indice suffit souvent à identifier le contexte.

\textbf{Stratégie 2 : Repérer les mots-clés (champ lexical).}\par
Chaque lieu possède son vocabulaire porteur. Entends-tu ces mots ? $\rightarrow$ Tu sais où tu es.
\begin{center}
\begin{tabularx}{\linewidth}{|X|X|}
\hline
\rowcolor{popblueL} \textbf{Lieu} & \textbf{Mots-clés attendus (Keywords)} \\
\hline
Banque & \eng{account, deposit, withdraw, balance, ID card, signature, current, savings, transfer} \\
\hline
Hôtel & \eng{room, single, double, night, key, booking, reservation, check-in, check-out, breakfast, bill} \\
\hline
Garage & \eng{repair, fix, engine, tyre, oil, brake, bill, estimate, spare part, mechanic} \\
\hline
Téléphonie (MTN/Orange) & \eng{bundle, airtime, data, SIM card, register, top up, recharge, network, promo} \\
\hline
Hôpital / Clinique & \eng{appointment, doctor, prescription, specialist, symptom, pharmacy, ward, nurse} \\
\hline
\end{tabularx}
\end{center}

\textbf{Stratégie 3 : Noter les chiffres et les noms propres.}\par
Les questions d'examen portent très souvent sur des prix, des numéros de chambre, des dates ou des noms de clients.
\begin{itemize}
\item \textbf{Prix :} \eng{25 000 francs CFA} $\rightarrow$ note \textbf{25 000 F} ou \textbf{25k}. Attention aux confusions auditives : \eng{fifteen} (15) vs \eng{fifty} (50) ; \eng{thirteen} (13) vs \eng{thirty} (30) ; \eng{nineteen} (19) vs \eng{ninety} (90). \textbf{Astuce :} écoute l'accent tonique (\eng{fif\underline{teen}} vs \eng{\underline{fif}ty}).
\item \textbf{Noms propres :} Souvent épelés lettre par lettre. \eng{My name is Ngo Bassa — that's N-G-O, B-A-S-S-A.} Note-les en \textbf{majuscules} immédiatement.
\item \textbf{Dates/Heures :} \eng{On Friday the 12th at 10 a.m.} $\rightarrow$ \textbf{Fri 12, 10h}.
\end{itemize}

\textbf{Stratégie 4 : Écouter les formules d'ouverture et de clôture.}\par
Le début et la fin sont prévisibles et riches.
\begin{itemize}
\item \textbf{Ouverture employé :} \eng{Good morning, madam/sir. How may I help you? / What can I do for you?}
\item \textbf{Ouverture client :} \eng{I'd like to... / I was wondering if... / Could you...?}
\item \textbf{Clôture :} \eng{Thank you for your help. / You're welcome. Have a nice day! / Here is your receipt/key.}
\end{itemize}
Ces formules confirment le thème et le degré de formalité.

\begin{center}
\begin{tabularx}{\linewidth}{|X|X|X|}
\hline
\rowcolor{popblueL} Stratégie & Ce que tu écoutes / notes & Exemple de signal \\
\hline
1. Gist (Thème) & Formule d'ouverture, ton général, lieu & \eng{How can I help you?} $\rightarrow$ Guichet \\
\hline
2. Mots-clés & Champ lexical du lieu (noms, verbes) & \eng{deposit, bundle, tyre, prescription} \\
\hline
3. Chiffres \& Noms & Prix, numéros, dates, noms épelés & \eng{Room 204, 30 000 F, Mr Etame, N-G-O} \\
\hline
4. Ouverture / Clôture & Début et fin du dialogue & \eng{I'd like to book...} / \eng{Have a nice day!} \\
\hline
\end{tabularx}
\end{center}

\begin{popfigure}
\begin{center}\tcbox[colback=popblue,colframe=popblue,coltext=white,arc=9pt,boxsep=4pt,left=6pt,right=6pt]{\bfseries\small LISTENING STRATEGIES}\end{center}
\vspace{-2pt}
\begin{tcbraster}[raster columns=3,raster equal height=rows,raster column skip=6pt,raster row skip=6pt,enhanced,blankest]
\begin{tcolorbox}[enhanced,colback=white,colframe=poporange!70,boxrule=0.9pt,arc=3pt,title={1. Gist (1st listening)},fonttitle=\bfseries\scriptsize,coltitle=white,colbacktitle=poporange!70,left=4pt,right=4pt,top=2pt,bottom=2pt,boxsep=1pt,toptitle=1.5pt,bottomtitle=1.5pt,halign=left]
\begin{itemize}[nosep,leftmargin=*,itemsep=1pt,font=\scriptsize]
\item Opening lines
\item Tone / Speakers
\end{itemize}
\end{tcolorbox}
\begin{tcolorbox}[enhanced,colback=white,colframe=popgreen!70,boxrule=0.9pt,arc=3pt,title={2. Keywords (Lexical field)},fonttitle=\bfseries\scriptsize,coltitle=white,colbacktitle=popgreen!70,left=4pt,right=4pt,top=2pt,bottom=2pt,boxsep=1pt,toptitle=1.5pt,bottomtitle=1.5pt,halign=left]
\begin{itemize}[nosep,leftmargin=*,itemsep=1pt,font=\scriptsize]
\item Bank: account, deposit
\item Hotel: room, night
\item Garage: repair, bill
\item Phone: bundle, SIM
\end{itemize}
\end{tcolorbox}
\begin{tcolorbox}[enhanced,colback=white,colframe=poppurple!70,boxrule=0.9pt,arc=3pt,title={3. Numbers \& Names (Details)},fonttitle=\bfseries\scriptsize,coltitle=white,colbacktitle=poppurple!70,left=4pt,right=4pt,top=2pt,bottom=2pt,boxsep=1pt,toptitle=1.5pt,bottomtitle=1.5pt,halign=left]
\begin{itemize}[nosep,leftmargin=*,itemsep=1pt,font=\scriptsize]
\item Prices (CFA)
\item Room / Account \#
\item Spelling names
\item Dates / Times
\end{itemize}
\end{tcolorbox}
\begin{tcolorbox}[enhanced,colback=white,colframe=poppink!70,boxrule=0.9pt,arc=3pt,title={4. Opening \& Closing (Routines)},fonttitle=\bfseries\scriptsize,coltitle=white,colbacktitle=poppink!70,left=4pt,right=4pt,top=2pt,bottom=2pt,boxsep=1pt,toptitle=1.5pt,bottomtitle=1.5pt,halign=left]
\begin{itemize}[nosep,leftmargin=*,itemsep=1pt,font=\scriptsize]
\item Could you...?
\item I'd like to...
\item Thank you / Welcome
\end{itemize}
\end{tcolorbox}
\begin{tcolorbox}[enhanced,colback=white,colframe=popgold!80,boxrule=0.9pt,arc=3pt,title={5. Note-taking (2-column table)},fonttitle=\bfseries\scriptsize,coltitle=white,colbacktitle=popgold!80,left=4pt,right=4pt,top=2pt,bottom=2pt,boxsep=1pt,toptitle=1.5pt,bottomtitle=1.5pt,halign=left]
\begin{itemize}[nosep,leftmargin=*,itemsep=1pt,font=\scriptsize]
\item Customer / Employee
\item Abbreviations
\end{itemize}
\end{tcolorbox}
\end{tcbraster}

\legende{Carte mentale : les cinq stratégies d'écoute pour un dialogue de demande de service.}
\end{popfigure}

\courssub{Technique de prise de notes : le tableau à deux colonnes}

Pendant la deuxième écoute, ne rédige pas de phrases complètes. Utilise un tableau à deux colonnes — une pour le client (\textbf{C}), une pour l'employé (\textbf{E}) — et des abréviations personnelles.

\begin{exemplebox}[Exemple de prise de notes efficace]
\textbf{Script entendu (extrait banque) :}
\eng{— Good morning. I'd like to open a savings account, please.}
\eng{— Certainly, madam. Could I have your ID card? The minimum deposit is 10 000 francs.}
\eng{— Here it is. I have 50 000 francs to deposit.}
\eng{— Thank you. Please sign here. Your account number is 204551.}

\begin{center}
\begin{tabular}{|l|l|}
\hline
\rowcolor{popblueL} \textbf{Customer (C)} & \textbf{Employee (E)} \\
\hline
wants to open savings acc. & asks for ID card \\
\hline
has 50 000 F to dep. & min. dep. = 10 000 F \\
\hline
— & sign here \\
\hline
— & acc. no: 204551 \\
\hline
\end{tabular}
\end{center}

\textbf{Abréviations recommandées :}
\begin{itemize}
\item \textbf{acc.} = account / \textbf{acc. no} = account number
\item \textbf{rm} = room / \textbf{rm no} = room number
\item \textbf{dep.} = deposit / \textbf{withdr.} = withdrawal
\item \textbf{F} = francs CFA / \textbf{k} = mille (ex: 25k = 25 000)
\item \textbf{bkfst} = breakfast / \textbf{appt} = appointment
\item \textbf{=} pour les prix / \textbf{?} pour info manquée / \textbf{$\rightarrow$} pour action suivante
\end{itemize}
\end{exemplebox}

\courssub{Entraînement : Scripts d'écoute et compréhension}

Voici trois dialogues types. Ton professeur les lira (ou les fera écouter) deux fois. Applique la méthode : \textbf{1. Lis les questions $\rightarrow$ 2. Écoute gist $\rightarrow$ 3. Écoute détails $\rightarrow$ 4. Vérifie.}

\begin{textebox}[Script 1 : At the Bank — Opening an Account]
\textbf{Context:} A customer enters a bank in Douala to open a savings account.

\eng{— Good morning, madam. How can I help you today?}
\eng{— Good morning. I'd like to open a savings account, please.}
\eng{— Certainly. Could I have your national ID card and a proof of address?}
\eng{— Yes, here is my ID card. My name is Ngo Bassa — that's N-G-O, B-A-S-S-A. And here is my electricity bill for the address.}
\eng{— Thank you. The minimum opening deposit is 10 000 francs CFA. How much would you like to deposit today?}
\eng{— I have 50 000 francs.}
\eng{— Very good. Please sign this form here, and here. Your account number is 2-0-4-5-5-1. Here is your passbook.}
\eng{— Thank you very much. Have a nice day.}
\eng{— You're welcome. Goodbye.}
\end{textebox}

\begin{exercice}[Compréhension : Script 1 — At the Bank]
\textbf{Answer the following questions in English.}
\begin{enumerate}
\item Where does the conversation take place?
\item What service does the customer request?
\item What is the customer's full name? (Spell it)
\item What two documents does the employee ask for?
\item What is the minimum deposit required?
\item How much does the customer actually deposit?
\item What is the account number given?
\end{enumerate}
\end{exercice}


\begin{corrige}[Exercice Script 1]
\begin{enumerate}
\item At a bank (in Douala).
\item She wants to open a savings account.
\item Ngo Bassa (N-G-O, B-A-S-S-A).
\item National ID card and a proof of address (electricity bill).
\item 10 000 francs CFA.
\item 50 000 francs CFA.
\item 204551.
\end{enumerate}
\end{corrige}

\begin{textebox}[Script 2 : At the Hotel — Booking a Room]
\textbf{Context:} A businessman arrives at a hotel in Yaoundé.

\eng{— Good evening, sir. Welcome to Hotel Mont Fébé. Do you have a reservation?}
\eng{— No, I don't. I was wondering if you have a single room available for tonight.}
\eng{— Let me check... Yes, we have a single room on the 3rd floor. Room 304. It's 35 000 francs CFA per night, breakfast included.}
\eng{— That sounds fine. I'll take it.}
\eng{— Great. Could you fill in this registration form, please? I'll also need to see your passport.}
\eng{— Here you are. My name is Etame — E-T-A-M-E.}
\eng{— Thank you, Mr Etame. Here is your key. Breakfast is served from 6:30 to 9:30. Enjoy your stay!}
\eng{— Thank you. Good night.}
\end{textebox}

\begin{exercice}[Compréhension : Script 2 — At the Hotel]
\textbf{True or False? Justify with a sentence from the text.}
\begin{enumerate}
\item The customer has a reservation.
\item The room is a double room on the second floor.
\item The price is 35 000 francs CFA per night.
\item Breakfast is not included in the price.
\item The customer's name is spelled E-T-A-M-E.
\item Breakfast starts at 6:00 a.m.
\end{enumerate}
\end{exercice}


\begin{corrige}[Exercice Script 2]
\begin{enumerate}
\item False. \eng{"No, I don't."} / \eng{"I was wondering if you have a single room available..."}
\item False. It is a \textbf{single} room on the \textbf{3rd} floor (Room 304).
\item True. \eng{"It's 35 000 francs CFA per night..."}
\item False. \eng{"...breakfast included."}
\item True. \eng{"My name is Etame — E-T-A-M-E."}
\item False. \eng{"Breakfast is served from 6:30 to 9:30."}
\end{enumerate}
\end{corrige}

\begin{textebox}[Script 3 : At the Phone Agency (MTN) — Reporting a Faulty Bundle]
\textbf{Context:} A student at the MTN agency in Buea complains about a data bundle.

\eng{— Good morning. Welcome to MTN. How may I assist you?}
\eng{— Good morning. I bought a 1 000 francs data bundle yesterday, but it's not working. I can't browse.}
\eng{— I'm sorry to hear that. Let me check your number. What is your phone number?}
\eng{— It's 6 77 12 34 56.}
\eng{— And your name, please?}
\eng{— F-O-N-G-A-N-G, Fonjang.}
\eng{— Thank you. I see the bundle was activated but there is a network issue in your area. We will fix it within 2 hours. As compensation, you will receive 500 MB free.}
\eng{— Okay, thank you. Will I receive an SMS confirmation?}
\eng{— Yes, you will. Is there anything else I can help you with?}
\eng{— No, that's all. Thanks. Bye.}
\eng{— You're welcome. Have a good day.}
\end{textebox}

\begin{exercice}[Compréhension : Script 3 — At the Phone Agency]
\textbf{Multiple Choice. Choose the correct answer (A, B or C).}
\begin{enumerate}
\item What problem does the customer report?
A) He lost his phone. \quad B) His data bundle is not working. \quad C) He wants to buy a new SIM card.
\item How much did the bundle cost?
A) 500 F \quad B) 1 000 F \quad C) 2 000 F
\item What is the customer's phone number?
A) 6 77 12 34 56 \quad B) 6 77 21 43 65 \quad C) 6 78 12 34 56
\item How long will the repair take?
A) 30 minutes \quad B) 1 hour \quad C) 2 hours
\item What compensation is offered?
A) 1 000 F credit \quad B) 500 MB free data \quad C) A new SIM card
\end{enumerate}
\end{exercice}


\begin{corrige}[Exercice Script 3]
\begin{enumerate}
\item B (\eng{"I bought a 1 000 francs data bundle yesterday, but it's not working."})
\item B (\eng{"1 000 francs"})
\item A (\eng{"It's 6 77 12 34 56."})
\item C (\eng{"We will fix it within 2 hours."})
\item B (\eng{"As compensation, you will receive 500 MB free."})
\end{enumerate}
\end{corrige}

\courssub{Lexique thématique récapitulatif : Requesting Services}

\begin{center}
\begin{tabularx}{\linewidth}{|X|X|X|}
\hline
\rowcolor{popblueL} \textbf{Verbes d'action} & \textbf{Noms utiles} & \textbf{Expressions clés (Client $\rightarrow$ Employé)} \\
\hline
to book / to reserve & a booking / a reservation & \eng{I'd like to book...} \\
to open (an account) & an account (current/savings) & \eng{I'd like to open an account.} \\
to deposit / to withdraw & a deposit / a withdrawal & \eng{I want to deposit / withdraw...} \\
to top up / to recharge & airtime / credit / a bundle & \eng{I'd like to top up my account.} \\
to register & a SIM card / a subscription & \eng{I need to register my SIM card.} \\
to report / to complain & a fault / an issue / a problem & \eng{I'd like to report a faulty bundle.} \\
to repair / to fix & a repair / a spare part / a bill & \eng{My car needs a repair.} \\
to check in / to check out & a key / a receipt / a bill & \eng{I'd like to check in / out.} \\
to fill in / to complete & a form / a registration form & \eng{Could you fill in this form?} \\
to sign & a signature & \eng{Please sign here.} \\
\hline
\end{tabularx}
\end{center}

\begin{attention}[Prononciation : Pièges auditifs (Minimal Pairs)]
Entraîne-toi à distinguer ces paires, cruciales pour les chiffres et les noms :
\begin{center}
\begin{tabular}{|c|c|c|}
\hline
\rowcolor{popblueL} \textbf{Mot 1} & \textbf{Mot 2} & \textbf{Astuce} \\
\hline
\eng{fifteen} (15) /fɪfˈtiːn/ & \eng{fifty} (50) /ˈfɪfti/ & Accent sur 2e syllabe vs 1re ; \eng{teen} = long. \\
\hline
\eng{thirteen} (13) /θɜːˈtiːn/ & \eng{thirty} (30) /ˈθɜːti/ & Idem. \eng{thir-teen} vs \eng{thir-ty}. \\
\hline
\eng{nineteen} (19) /naɪnˈtiːn/ & \eng{ninety} (90) /ˈnaɪnti/ & Idem. \\
\hline
\eng{A} /eɪ/ (« éï ») & \eng{E} /iː/ (« i ») & \eng{A} = nom de la lettre A ; \eng{E} = nom de la lettre E. \\
\hline
\eng{I} /aɪ/ (« aï ») & \eng{Y} /waɪ/ (« waï ») & \eng{I} = diphtongue ; \eng{Y} = commence par w. \\
\hline
\eng{G} /dʒiː/ (« dji ») & \eng{J} /dʒeɪ/ (« djéï ») & \eng{G} rime avec \eng{see} ; \eng{J} rime avec \eng{day}. \\
\hline
\end{tabular}
\end{center}
\textbf{Exercice :} Le professeur épelle : \eng{N-G-O, B-A-S-S-A} ; \eng{E-T-A-M-E} ; \eng{F-O-N-G-A-N-G}. Écris les noms.
\end{attention}

\courssub{Production guidée : Je demande un service (Speaking \& Writing)}

\begin{experience}[Role-play : At the Service Desk]
\textbf{Consigne :} En binôme. L'un joue le \textbf{Client}, l'autre l'\textbf{Employé}. Choisis une situation ci-dessous. Prépare 2 minutes (notes autorisées, pas de texte lu). Joue le dialogue (1 à 2 min). Change de rôle.

\textbf{Situation A (Banque - Yaoundé) :} Tu veux ouvrir un \eng{current account}. Dépôt initial : 25 000 F. Tu donnes ton nom (épelé) et montre ta CNI.
\textbf{Situation B (Hôtel - Limbe) :} Tu arrives sans réservation. Tu veux une \eng{double room} pour 2 nuits. Prix ? Petit-déj inclus ? Tu donnes ton passeport.
\textbf{Situation C (Garage - Douala) :} Tu récupères ta voiture. Réparation : plaquettes de frein (\eng{brake pads}). Prix total ? Tu paies par carte (\eng{by card}).
\textbf{Situation D (Agence Orange - Buea) :} Ton forfait \eng{social bundle} ne marche pas. Tu donnes ton numéro. L'agent vérifie : panne réseau, réparation 1h. Compensation : 200 Mo.

\textbf{Critères de réussite :}
\begin{itemize}
\item Utilisation d'au moins 2 formules polies (\eng{Could you...}, \eng{I'd like to...}, \eng{I was wondering if...}).
\item Vocabulaire précis du lieu (account, room, brake pads, bundle).
\item Communication du prix, nom, numéro sans hésitation majeure.
\item Clôture polie (\eng{Thank you / Have a nice day}).
\end{itemize}
\end{experience}

\begin{exercice}[Writing : Formal Email / Letter of Request]
\textbf{Consigne :} Rédige un email formel (80-100 mots) dans l'une des situations ci-dessous. Respecte la structure : \textbf{Objet} — \textbf{Formule d'appel} — \textbf{Contexte} — \textbf{Demande précise (polie)} — \textbf{Infos utiles (nom, numéro, date)} — \textbf{Formule de politesse finale} — \textbf{Signature}.

\textbf{Option 1 :} Tu es étudiant à l'Université de Buea. Écris au \textbf{Registrar} (Secrétaire académique) pour demander un \eng{transcript of records} (relevé de notes) pour une bourse. Donne ton \eng{student ID number}.
\textbf{Option 2 :} Tu es client à la \textbf{BCEAO} (banque). Écris au \textbf{Branch Manager} pour demander l'augmentation de ton plafond de retrait journalier sur ta carte Visa. Donne ton \eng{account number}.
\textbf{Option 3 :} Tu as séjourné à l'\textbf{Hôtel Akwa Palace} (Douala). Écris au \textbf{Manager} pour te plaindre du bruit (travaux) et demander un geste commercial (réduction ou nuit offerte). Donne ton \eng{booking reference}.

\textbf{Modèle de phrases utiles :}
\begin{itemize}
\item \eng{Subject: Request for... / Complaint regarding...}
\item \eng{Dear Sir / Madam,}
\item \eng{I am writing to request... / to complain about...}
\item \eng{I would be grateful if you could...}
\item \eng{I was wondering if it would be possible to...}
\item \eng{My details are as follows: Name: ... / Account No: ... / Booking Ref: ...}
\item \eng{I look forward to your reply.}
\item \eng{Yours faithfully,} (si \eng{Dear Sir/Madam}) / \eng{Yours sincerely,} (si nom connu)
\item \eng{[Ton Nom Complet]}
\end{itemize}
\end{exercice}


\begin{corrige}[Exercice Writing — Modèle type (Option 1)]
\eng{Subject: Request for Academic Transcript — Student ID 2024/ENG/045}

\eng{Dear Sir / Madam,}

\eng{I am writing to request an official transcript of my academic records for the 2023-2024 academic year. I need this document to apply for a Master's scholarship abroad.}

\eng{I would be grateful if you could issue the transcript as soon as possible and let me know when I can collect it from the registry.}

\eng{My details are as follows:}
\eng{Full Name: Ngo Bassa Marie}
\eng{Student ID Number: 2024/ENG/045}
\eng{Department: English Modern Letters}
\eng{Phone: 6 77 12 34 56}

\eng{I look forward to your favourable reply.}

\eng{Yours faithfully,}

\eng{Ngo Bassa Marie}
\end{corrige}

\begin{savaistu}[Culture \& Exam Tips : Listening au Baccalauréat]
\begin{itemize}
\item \textbf{Format officiel :} L'enregistrement est lu \textbf{deux fois}. Les questions (Vrai/Faux, QCM, Réponses courtes) sont distribuées \textbf{avant} la 1ère écoute. \textbf{Profite de ce temps de lecture !}
\item \textbf{Statistiques examinateurs :} Les candidats qui soulignent les mots-clés des questions (\eng{where, how much, name}) et qui préparent un tableau \textbf{Client / Employé} gagnent en moyenne 3 à 4 points de plus sur 20.
\item \textbf{Épellation (Spelling) :} Au Cameroun, les noms propres sont \textbf{quasiment toujours épelés} dans les scripts officiels (Baccalauréat, Probatoire). Entraîne-toi à reconnaître l'alphabet anglais \textbf{à l'oreille} : \eng{A (éï), E (i), I (aï), G (dji), J (djéï), H (eïtch), R (are), Y (waï)}. Ce sont les pièges n°1 pour les francophones.
\item \textbf{Devises :} Les prix sont en \textbf{francs CFA} (souvent abrégé \eng{frs} ou \eng{F}). Note \eng{25 000 F} ou \eng{25k}. Ne convertis pas en euros/dollars.
\item \textbf{Accents :} Les speakers peuvent avoir un accent camerounais, nigérian, ghanéen ou britannique standard. Habitue-toi à la variété (BBC Learning English, VOA Learning English, Cameroon Radio Television - CRTV News en anglais).
\end{itemize}
\end{savaistu}

\begin{exoresolu}[Bilan : Simulation d'examen (Mini-Test)]
Énoncé. \textbf{Écoute le script suivant (lu 2 fois par le prof) et réponds aux questions.}

\eng{— Good morning. Welcome to the General Hospital.}
\eng{— Good morning. I'd like to make an appointment with a cardiologist, please.}
\eng{— Do you have a referral letter from your GP?}
\eng{— Yes, here it is. Dr Mbarga referred me.}
\eng{— Okay. The earliest slot is Thursday 15th at 10:00 a.m. with Dr Nkomo.}
\eng{— That works. My name is T-A-B-I, Tabi.}
\eng{— Thank you, Mr Tabi. Your appointment is confirmed for Thursday 15th, 10 a.m., Cardiology Department, Room 12. Please bring your ID and the referral letter.}
\eng{— Thank you. Goodbye.}

\textbf{Questions :}
1. Where does the conversation take place?
2. What specialist does the patient want to see?
3. Does the patient have a referral letter?
4. What is the date and time of the appointment?
5. What is the patient's name (spelled)?
6. Which doctor will he see?
7. What two documents must he bring?

\tcblower
\textbf{Solution détaillée.}
1. \textbf{At the General Hospital} (Ouverture : \eng{Welcome to the General Hospital}).
2. \textbf{A cardiologist} (Mots-clés : \eng{cardiologist, Cardiology Department}).
3. \textbf{Yes}, he has a referral letter from Dr Mbarga.
4. \textbf{Thursday 15th at 10:00 a.m.} (Chiffres + jour + heure).
5. \textbf{T-A-B-I, Tabi} (Épellation explicite).
6. \textbf{Dr Nkomo}.
7. \textbf{His ID and the referral letter} (Clôture : \eng{Please bring your ID and the referral letter}).

\textbf{Analyse des stratégies utilisées :}
\begin{itemize}
\item \textbf{Gist :} Hôpital, rendez-vous médical (Stratégie 1).
\item \textbf{Keywords :} \eng{cardiologist, referral, appointment, department} (Stratégie 2).
\item \textbf{Numbers/Names :} \eng{15th, 10:00, Room 12, Tabi (spelled)} (Stratégie 3).
\item \textbf{Opening/Closing :} Confirmation des détails à la fin (Stratégie 4).
\item \textbf{Note-taking :} Tableau C/E efficace pour séparer infos patient / infos réceptionniste.
\end{itemize}
\end{exoresolu}

\coursec{Script d'écoute : At the service desk}

Dans la section précédente, nous avons préparé notre écoute : identifier le lieu, les personnages et le type de demande \emph{avant} d'écouter, puis repérer les mots-clés, les chiffres et les noms propres \emph{pendant} l'écoute. Mettons maintenant ces stratégies en pratique avec trois dialogues authentiques. Chaque scène se déroule à un \cle{service desk} (un guichet, un comptoir d'accueil) où un client formule une demande et où un employé y répond : à la banque, à l'hôtel, puis au garage. Lis chaque script à voix haute avec ton professeur, puis ferme le manuel et écoute : sauras-tu relever les informations essentielles ?

\courssub{Script 1 : At the bank — opening a savings account}

\begin{textebox}[Script 1 — At the bank]
\textbf{Clerk:} Good morning, sir. How may I help you?\par
\textbf{Customer:} Good morning. I'd like to open a savings account, please.\par
\textbf{Clerk:} Certainly. Could you fill in this form and show me your ID card?\par
\textbf{Customer:} Here you are. How much is the minimum deposit?\par
\textbf{Clerk:} It's 10,000 francs CFA.\par
\textbf{Customer:} Fine. And could I also get a cheque book?\par
\textbf{Clerk:} Of course. It will be ready in two days.
\end{textebox}

\textbf{Personnages et situation.}\par
Un \cle{clerk} (employé de banque) accueille un \cle{customer} (client) qui souhaite ouvrir un \cle{savings account} (compte épargne). Observe la structure typique de tout dialogue de service :

\begin{enumerate}
\item \textbf{Greeting} (salutation) : \eng{Good morning, sir.}
\item \textbf{Offer of help} (offre d'aide) : \eng{How may I help you?} « Comment puis-je vous aider ? »
\item \textbf{Request} (demande du client) : \eng{I'd like to open a savings account, please.}
\item \textbf{Steps and details} (étapes et détails) : formulaire, pièce d'identité, dépôt minimum.
\item \textbf{Extra request} (demande supplémentaire) : un chéquier.
\item \textbf{Closing information} (information finale) : \eng{It will be ready in two days.}
\end{enumerate}

\begin{definition}[Glossaire du Script 1]
\begin{itemize}
\item \eng{savings account} (n.) : compte épargne
\item \eng{to fill in a form} (v.) : remplir un formulaire (anglais britannique)
\item \eng{ID card} (n.) : carte d'identité
\item \eng{deposit} (n.) : dépôt (somme versée) ; \eng{minimum deposit} : dépôt minimum
\item \eng{cheque book} (n.) : chéquier (orthographe britannique : \eng{cheque})
\item \eng{Here you are.} : « Voilà. » (formule employée en tendant un objet à quelqu'un)
\end{itemize}
\end{definition}

\begin{attention}[Faux amis et pièges]
\eng{Actually} ne signifie pas « actuellement » mais « en fait » ; « actuellement » se dit \eng{currently}. De même, \eng{a deposit} est un « dépôt » (argent versé), pas un « dépôt » au sens de magasin ou d'entrepôt. Enfin, on écrit \eng{10,000 francs} avec une virgule en anglais pour séparer les milliers, là où le français utilise un espace insécable (« 10 000 francs »).
\end{attention}

\courssub{Script 2 : At the hotel — booking a room with a view}

\begin{textebox}[Script 2 — At the hotel]
\textbf{Guest:} Good evening. I was wondering if you have a double room with a sea view.\par
\textbf{Receptionist:} Let me check... Yes, room 204 is free. It's 35,000 francs a night, breakfast included.\par
\textbf{Guest:} Perfect. Could you also wake me up at 6 a.m.?\par
\textbf{Receptionist:} No problem, madam.
\end{textebox}

\textbf{Personnages et situation.}\par
Une \cle{guest} (cliente de l'hôtel) s'adresse au \cle{receptionist} (réceptionniste). Elle formule deux demandes : une \cle{double room with a sea view} (chambre double avec vue sur la mer) et un réveil à 6 heures du matin. Le réceptionniste accepte et précise le prix : \eng{35,000 francs a night, breakfast included} « 35 000 francs la nuit, petit-déjeuner compris ».

\begin{definition}[Glossaire du Script 2]
\begin{itemize}
\item \eng{guest} (n.) : client(e) d'un hôtel (ou invité)
\item \eng{double room} (n.) : chambre double ; \eng{single room} : chambre simple
\item \eng{sea view} (n.) : vue sur la mer
\item \eng{breakfast included} : petit-déjeuner compris
\item \eng{to wake somebody up} (v.) : réveiller quelqu'un (verbe à particule) ; \eng{wake me up at 6 a.m.} : réveillez-moi à 6 heures du matin
\item \eng{Let me check.} : « Laissez-moi vérifier. / Je vais voir. »
\end{itemize}
\end{definition}

\begin{propriete}[La demande très polie : I was wondering if...]
\eng{I was wondering if you have a double room...} est une formule de demande \emph{très polie et indirecte}, plus élégante que \eng{Do you have a double room?}. Structure : \textbf{I was wondering if + sujet + verbe}. Le passé (\eng{was wondering}) adoucit la demande — c'est le même principe que le conditionnel français « Je me demandais si vous auriez... ».
\begin{itemize}
\item On utilise souvent le présent après \eng{if} pour une disponibilité actuelle : \eng{I was wondering if you \emph{have} a room.}
\item Le prétérit \eng{had} est aussi possible (backshifting) : \eng{I was wondering if you \emph{had} a room.}
\item \textbf{Attention :} on n'utilise \textbf{pas} le futur (\eng{will have}) après \eng{if} dans cette structure.
\end{itemize}
\end{propriete}

\courssub{Script 3 : At the garage — asking for a repair and an estimate}

\begin{textebox}[Script 3 — At the garage]
\textbf{Customer:} My motorbike won't start. Could you have a look at it?\par
\textbf{Mechanic:} Sure. I'm afraid the engine needs a full repair.\par
\textbf{Customer:} How much would it cost?\par
\textbf{Mechanic:} About 25,000 francs. I can give you a written estimate.\par
\textbf{Customer:} All right. When will it be ready?\par
\textbf{Mechanic:} By Friday afternoon.
\end{textebox}

\textbf{Personnages et situation.}\par
Un client arrive chez le \cle{mechanic} (mécanicien) parce que sa moto ne démarre pas : \eng{My motorbike won't start.} Il demande un examen, puis un prix, puis un \cle{written estimate} (devis écrit). Le mécanicien annonce une mauvaise nouvelle avec la formule d'excuse \eng{I'm afraid...} et donne un délai : \eng{By Friday afternoon} « d'ici vendredi après-midi / au plus tard vendredi après-midi ».

\begin{definition}[Glossaire du Script 3]
\begin{itemize}
\item \eng{motorbike} (n.) : moto (à Douala, on pense aux fameux \emph{bendskins} !)
\item \eng{won't start} : ne démarre pas (\eng{won't} = \eng{will not} ; refus de fonctionner)
\item \eng{to have a look at} (v.) : jeter un coup d'œil à, examiner
\item \eng{engine} (n.) : moteur
\item \eng{repair} (n. et v.) : réparation ; réparer
\item \eng{estimate} (n.) : devis ; \eng{written estimate} : devis écrit
\item \eng{by Friday} : d'ici vendredi, au plus tard vendredi (échéance)
\end{itemize}
\end{definition}

\begin{attention}[\eng{By} ou \eng{until} ?]
\eng{By Friday} signifie « au plus tard vendredi » (limite, échéance). \eng{Until Friday} signifie « jusqu'à vendredi » (durée continue). \eng{It will be ready by Friday.} « Ce sera prêt d'ici vendredi. » \eng{The garage is open until Friday.} « Le garage est ouvert jusqu'à vendredi. »
Autre piège : \eng{won't start} — les francophones disent souvent \emph{doesn't want to start} (calque de « ne veut pas démarrer ») ; en anglais, \eng{will not} exprime ici le refus de fonctionner, très naturel pour une machine ou un objet inanimé.
\end{attention}

\courssub{Vue d'ensemble : lieux, demandes et réponses}

\begin{popfigure}
\begin{tikzpicture}[every node/.style={font=\small}]
\node[draw=popblue, fill=popblueL, rounded corners, align=center, minimum width=3.1cm, minimum height=1.6cm] (b1) at (0,0) {\textbf{Bank}\\ open a savings\\ account};
\node[draw=popgreen, fill=popgreenL, rounded corners, align=center, minimum width=3.1cm, minimum height=1.6cm] (b2) at (5.2,0) {\textbf{Hotel}\\ double room\\ with a sea view};
\node[draw=poporange, fill=poporangeL, rounded corners, align=center, minimum width=3.1cm, minimum height=1.6cm] (b3) at (10.4,0) {\textbf{Garage}\\ repair a motorbike\\ + estimate};
\node[draw=popblue, fill=white, rounded corners, align=center, minimum width=3.1cm, minimum height=1.6cm] (r1) at (0,-2.6) {\textbf{Clerk}\\ ``Certainly.''\\ form + ID card};
\node[draw=popgreen, fill=white, rounded corners, align=center, minimum width=3.1cm, minimum height=1.6cm] (r2) at (5.2,-2.6) {\textbf{Receptionist}\\ ``No problem.''\\ room 204, 35,000 F};
\node[draw=poporange, fill=white, rounded corners, align=center, minimum width=3.1cm, minimum height=1.6cm] (r3) at (10.4,-2.6) {\textbf{Mechanic}\\ ``I'm afraid...''\\ 25,000 F, by Friday};
\draw[-{Stealth}, thick, popblue] (b1) -- (r1);
\draw[-{Stealth}, thick, popgreen] (b2) -- (r2);
\draw[-{Stealth}, thick, poporange] (b3) -- (r3);
\node[font=\small\itshape, popmuted] at (0,1.3) {Customer's request};
\node[font=\small\itshape, popmuted] at (0,-3.9) {Employee's answer};
\end{tikzpicture}
\legende{Les trois dialogues : la demande du client (en haut) et la réponse de l'employé (en bas).}
\end{popfigure}

\courssub{Les formules de demande polie entendues dans les scripts}

Observons les demandes des trois clients :

\begin{exemplebox}[Les demandes des scripts]
\begin{itemize}
\item \eng{I'd like to open a savings account, please.} « Je voudrais ouvrir un compte épargne, s'il vous plaît. »
\item \eng{Could you fill in this form?} « Pourriez-vous remplir ce formulaire ? »
\item \eng{I was wondering if you have a double room with a sea view.} « Je me demandais si vous aviez une chambre double avec vue sur la mer. »
\item \eng{Could you also wake me up at 6 a.m.?} « Pourriez-vous aussi me réveiller à 6 heures ? »
\item \eng{Could you have a look at it?} « Pourriez-vous y jeter un coup d'œil ? »
\item \eng{How much would it cost?} « Combien cela coûterait-il ? »
\end{itemize}
\end{exemplebox}

\begin{propriete}[Former une demande polie en anglais]
\begin{itemize}
\item \textbf{I'd like + nom / base verbale} : \eng{I'd like a room.} / \eng{I'd like to book a room.} (\eng{I'd like} = \eng{I would like}, équivalent du conditionnel français « je voudrais »).
\item \textbf{Could you + base verbale ?} : plus poli que \eng{Can you...?}. Jamais de \emph{-s} ni de \emph{to} après \eng{could} : \eng{Could you help me?}
\item \textbf{I was wondering if + sujet + verbe (présent ou prétérit)} : le degré de politesse le plus élevé.
\item \textbf{How much / When + would / will...?} pour demander un prix ou un délai : \eng{How much is the deposit?} \eng{When will it be ready?}
\item Le mot \cle{please} se place généralement en fin de phrase, précédé d'une virgule : \eng{A single room, please.}
\end{itemize}
\end{propriete}

\begin{center}
\begin{tabularx}{\linewidth}{|X|X|X|}
\hline
\rowcolor{popblueL} Français & English & Degré de politesse \\
\hline
Je veux un chéquier. & I want a cheque book. & direct (à éviter en contexte de service) \\
\hline
Je voudrais un chéquier. & I'd like a cheque book, please. & poli (standard) \\
\hline
Pourriez-vous me réveiller à 6 h ? & Could you wake me up at 6 a.m.? & poli \\
\hline
Je me demandais si vous aviez une chambre. & I was wondering if you have a room. & très poli (indirect) \\
\hline
Combien coûte le dépôt minimum ? & How much is the minimum deposit? & neutre (information) \\
\hline
\end{tabularx}
\end{center}

\courssub{Accepter et refuser poliment : les réponses de l'employé}

\begin{propriete}[Formules d'acceptation et de refus poli]
\textbf{Accepter :} \eng{Certainly, sir.} « Bien sûr, monsieur. » ; \eng{Of course.} « Bien entendu. » ; \eng{No problem, madam.} « Pas de problème, madame. » ; \eng{Sure.} « Avec plaisir / D'accord. » (plus familier).\par
\textbf{Refuser ou annoncer une mauvaise nouvelle :} on commence toujours par \eng{I'm afraid...} « Je crains que... / J'ai bien peur que... / Malheureusement... » : \eng{I'm afraid we can't.} « Je crains que nous ne puissions pas. » ; \eng{I'm afraid the engine needs a full repair.} ; \eng{I'm afraid the hotel is full.} « Je crains que l'hôtel ne soit complet. » On adoucit souvent avec \eng{sorry} : \eng{I'm sorry, I'm afraid room 204 is taken.}
\end{propriete}

\begin{attention}[\eng{I'm afraid} ne veut pas dire « j'ai peur »]
Faux ami classique : \eng{I'm afraid we can't} ne signifie pas « j'ai peur que nous ne puissions pas » au sens de la frayeur, mais « je crains que / malheureusement ». C'est une formule de politesse figée pour annoncer un refus ou une mauvaise nouvelle. De même, \eng{No problem} est une réponse positive (« avec plaisir »), pas une question.
\end{attention}

\begin{exoresolu}[Repérage d'informations dans les scripts]
Énoncé — Réponds en anglais par une phrase courte :
\begin{enumerate}
\item What does the customer want to open at the bank?
\item How much is the minimum deposit?
\item Which room is free at the hotel, and what is the price?
\item What time does the guest want to wake up?
\item Why does the customer go to the garage?
\item When will the motorbike be ready?
\end{enumerate}
\tcblower
Solution détaillée :
\begin{enumerate}
\item \eng{He wants to open a savings account.} (Script 1, première demande du client.)
\item \eng{It is 10,000 francs CFA.} (Réponse du clerk ; relever un chiffre est une stratégie clé.)
\item \eng{Room 204 is free; it is 35,000 francs a night, breakfast included.} (Script 2.)
\item \eng{She wants to wake up at 6 a.m.} (Attention : \eng{wake up}, pas \emph{wake up myself}.)
\item \eng{Because his motorbike won't start.} (Script 3, première réplique.)
\item \eng{It will be ready by Friday afternoon.} (\eng{by} = au plus tard.)
\end{enumerate}
\textbf{Méthode :} souligne d'abord le mot interrogatif (\eng{what}, \eng{how much}, \eng{when}, \eng{why}) pour savoir quel type d'information chercher, puis localise la réplique correspondante dans le script.
\end{exoresolu}

\begin{exercice}[À toi de jouer]
Reconstitue un mini-dialogue poli au guichet d'une agence de voyage à Yaoundé : le client veut réserver deux billets de bus pour Buea (\eng{I'd like to book...}), demande le prix (\eng{How much...?}) ; l'employé accepte (\eng{Certainly.}) puis refuse poliment le paiement par chèque (\eng{I'm afraid we don't accept cheques.}). Écris six répliques.
\end{exercice}

\begin{aretenir}[L'essentiel de la section]
Trois lieux, trois demandes : à la banque (\eng{open a savings account}, dépôt de 10 000 F), à l'hôtel (\eng{a double room with a sea view}, 35 000 F la nuit), au garage (\eng{repair a motorbike}, devis de 25 000 F). Pour demander poliment : \eng{I'd like...}, \eng{Could you...?}, \eng{I was wondering if...}. Pour accepter : \eng{Certainly.}, \eng{No problem.} Pour refuser ou annoncer une mauvaise nouvelle : \eng{I'm afraid...}. À l'écoute, vise toujours : le lieu, la demande, les chiffres (prix, heures, numéros) et la réponse finale.
\end{aretenir}

\coursec{Compréhension : questions variées}

Dans la section précédente, nous avons écouté et lu les trois scripts de la leçon : \emph{Booking a hotel room} (à la réception d'un hôtel de Buea), \emph{Opening a bank account} (à la banque, avec un dépôt minimum) et \emph{At the repair shop} (le mécanicien qui répare une moto). Nous avons repéré le thème général — \cle{requesting services}, « demander un service » — et relevé les mots-clés. Il est maintenant temps de vérifier notre compréhension fine : c'est exactement le type d'épreuve que tu rencontreras au contrôle et à l'examen. Nous allons travailler quatre formats de questions : le \cle{vrai/faux} justifié, le \cle{QCM}, les \cle{réponses courtes} et la \cle{fiche à compléter}, avant de terminer par une question de synthèse.

\courssub{Rappel : les trois scripts en bref}

Avant de répondre aux questions, revoici l'essentiel des trois dialogues, condensé en un court texte de révision. Lis-le attentivement : toutes les réponses aux exercices s'y trouvent.

\begin{textebox}[Summary of the three listening scripts]
At the Mountain View Hotel in Buea, a traveller books a double room for two nights. The receptionist tells him: “It's 35,000 francs a night, breakfast included.” The guest asks for a quiet room and is given room 12 on the first floor. He pays in cash and receives his key.

At the bank in Yaoundé, a student wants to open a savings account. The clerk explains that the minimum deposit is 10,000 francs. The student fills in a form, shows his identity card and hands over the money. The clerk says the account will be active the next day and that the bank card will be ready in one week.

At a repair shop in Douala, a young woman brings her motorbike because the brakes do not work well. The mechanic checks the bike and says he must change the brake pads. The repair will cost 15,000 francs. “Come back tomorrow at four o'clock,” he says. “Your motorbike will be ready.”
\end{textebox}

\begin{definition}[Glossaire utile]
\begin{itemize}
\item \eng{double room} (n.) : chambre double
\item \eng{breakfast included} : petit-déjeuner compris
\item \eng{minimum deposit} (n.) : dépôt minimum
\item \eng{to fill in a form} : remplir un formulaire
\item \eng{identity card} (n.) : carte d'identité
\item \eng{brakes} (n. pl.) : freins ; \eng{brake pads} : plaquettes de frein
\item \eng{repair} (n.) : réparation ; \eng{to be ready} : être prêt
\end{itemize}
\end{definition}

\courssub{Exercice 1 : Vrai ou faux ? (justifie ta réponse)}

\begin{methode}[Comment réussir un vrai/faux en compréhension]
\begin{enumerate}
\item Lis d'abord toutes les affirmations avant de réécouter ou de relire le script.
\item Pour chaque affirmation, repère les mots-clés (chiffres, noms, lieux) et cherche-les dans le script.
\item \textbf{Souligne dans le script la phrase exacte qui prouve ta réponse} et recopie-la comme justification. Au Bac, une réponse non justifiée peut perdre des points, même si elle est correcte.
\item Attention aux affirmations « presque vraies » : un seul détail faux (le prix, le jour, le numéro) rend l'affirmation fausse.
\end{enumerate}
\end{methode}

\begin{exercice}[True or false? Justify with a sentence from the script]
\begin{enumerate}
\item The traveller books a single room at the hotel.
\item Breakfast is included in the price of the room.
\item The minimum deposit at the bank is 25,000 francs.
\item The bank card will be ready in one week.
\item The motorbike will be ready in two days.
\end{enumerate}
\end{exercice}

\begin{corrige}[Exercice 1]
\begin{enumerate}
\item \textbf{False.} The script says: “a traveller books a \emph{double} room for two nights.” (Il réserve une chambre double, pas simple.)
\item \textbf{True.} The receptionist says: “It's 35,000 francs a night, \emph{breakfast included}.”
\item \textbf{False.} The clerk says: “the minimum deposit is \emph{10,000} francs.” (10 000 F, pas 25 000 F.)
\item \textbf{True.} The clerk says: “the bank card will be ready \emph{in one week}.”
\item \textbf{False.} The mechanic says: “Come back \emph{tomorrow at four o'clock}.” (Le lendemain, pas dans deux jours.)
\end{enumerate}
\end{corrige}

\courssub{Exercice 2 : QCM — choisir la bonne réponse}

Le QCM (\eng{multiple choice questions}) porte sur les détails précis : prix, numéros, délais. La stratégie est la même : repérer le mot-clé de la question, retrouver la phrase du script, éliminer les distracteurs.

\begin{exemplebox}[Exemple traité en classe]
\textbf{Question :} The double room costs: a) 10,000 F \quad b) 25,000 F \quad c) 35,000 F

\textbf{Réponse : c)}, car la réceptionniste dit : \eng{“It's 35,000 francs a night.”} « C'est 35 000 francs la nuit. » Les options a) et b) sont des distracteurs : 10,000 F est le dépôt minimum de la banque, et 25,000 F n'apparaît nulle part — c'est un piège classique qui mélange les chiffres des trois scripts.
\end{exemplebox}

\begin{exercice}[Choose the correct answer]
\begin{enumerate}
\item The traveller's room number is: a) 2 \quad b) 12 \quad c) 21
\item The minimum deposit at the bank is: a) 5,000 F \quad b) 10,000 F \quad c) 15,000 F
\item The motorbike repair costs: a) 10,000 F \quad b) 15,000 F \quad c) 35,000 F
\item The motorbike will be ready: a) today at 4 p.m. \quad b) tomorrow at 4 p.m. \quad c) in one week
\end{enumerate}
\end{exercice}

\begin{corrige}[Exercice 2]
\begin{enumerate}
\item \textbf{b) 12} — “room 12 on the first floor.”
\item \textbf{b) 10,000 F} — “the minimum deposit is 10,000 francs.”
\item \textbf{b) 15,000 F} — “The repair will cost 15,000 francs.”
\item \textbf{b) tomorrow at 4 p.m.} — “Come back tomorrow at four o'clock.”
\end{enumerate}
\end{corrige}

\courssub{Exercice 3 : Réponses courtes (short answers)}

\begin{attention}[How much ou when ? Ne confonds pas les questions !]
\eng{How much...?} demande un \textbf{prix} ou une quantité : \eng{How much is the deposit? — 10,000 francs.} « Combien coûte le dépôt ? — 10 000 francs. » \eng{When...?} demande un \textbf{moment} : \eng{When will the motorbike be ready? — Tomorrow at 4 p.m.} « Quand la moto sera-t-elle prête ? — Demain à 16 h. » Beaucoup de francophones répondent un prix à une question \eng{when} parce qu'ils entendent un chiffre dans le script. Identifie d'abord le mot interrogatif, puis cherche le type d'information correspondant.
\end{attention}

\begin{center}
\begin{tabularx}{\linewidth}{|X|X|X|}
\hline
\rowcolor{popblueL} Mot interrogatif & Type de réponse & Exemple \\
\hline
How much...? & un prix, une quantité & How much is the room? — 35,000 francs. \\
\hline
When...? & un moment, une date & When will it be ready? — Tomorrow. \\
\hline
Where...? & un lieu & Where is the hotel? — In Buea. \\
\hline
What...? & une chose, un service & What does he want? — A double room. \\
\hline
How long...? & une durée & How long will he stay? — Two nights. \\
\hline
\end{tabularx}
\end{center}

\begin{exoresolu}[Short answers]
\textbf{Énoncé.} Réponds par une phrase courte et complète :
\begin{enumerate}
\item How much is the minimum deposit at the bank?
\item When will the motorbike be ready?
\item How long does the traveller stay at the hotel?
\end{enumerate}
\tcblower
\textbf{Solution détaillée.}
\begin{enumerate}
\item \eng{The minimum deposit is 10,000 francs.} — La question commence par \eng{how much} : on cherche un prix dans le script de la banque (« the minimum deposit is 10,000 francs »).
\item \eng{It will be ready tomorrow at four o'clock.} — La question commence par \eng{when} : on cherche un moment (« Come back tomorrow at four o'clock »). On reprend le futur \eng{will be ready} de la question.
\item \eng{He stays for two nights.} — \eng{How long} demande une durée : « a double room for two nights ».
\end{enumerate}
Remarque la méthode : on reprend la structure de la question (sujet + verbe) et on remplace seulement l'information demandée. Évite les réponses d'un seul mot à l'écrit : une phrase courte mais complète vaut toujours plus de points.
\end{exoresolu}

\courssub{Exercice 4 : Recherche d'informations — la fiche client}

Au restaurant, à l'hôtel ou dans une administration, on remplit souvent une fiche (\eng{a form}) avec les informations du client. C'est un excellent exercice d'écoute ciblée : tu ne cherches que quatre informations précises.

\begin{popfigure}
\begin{tikzpicture}
\node[draw=popblue, thick, rounded corners, fill=popblueL, minimum width=11cm, minimum height=0.9cm, font=\bfseries] at (0,3.4) {CLIENT FORM — Mountain Services};
\node[draw=popline, fill=white, minimum width=4.5cm, minimum height=0.8cm, anchor=west, font=\bfseries] at (-5.5,2.3) {Name};
\node[draw=popline, fill=popcream, minimum width=6cm, minimum height=0.8cm, anchor=west] at (-0.8,2.3) {};
\node[draw=popline, fill=white, minimum width=4.5cm, minimum height=0.8cm, anchor=west, font=\bfseries] at (-5.5,1.3) {Service requested};
\node[draw=popline, fill=popcream, minimum width=6cm, minimum height=0.8cm, anchor=west] at (-0.8,1.3) {};
\node[draw=popline, fill=white, minimum width=4.5cm, minimum height=0.8cm, anchor=west, font=\bfseries] at (-5.5,0.3) {Price};
\node[draw=popline, fill=popcream, minimum width=6cm, minimum height=0.8cm, anchor=west] at (-0.8,0.3) {};
\node[draw=popline, fill=white, minimum width=4.5cm, minimum height=0.8cm, anchor=west, font=\bfseries] at (-5.5,-0.7) {Time or deadline};
\node[draw=popline, fill=popcream, minimum width=6cm, minimum height=0.8cm, anchor=west] at (-0.8,-0.7) {};
\end{tikzpicture}
\legende{Fiche client vierge : à compléter en écoutant le script}
\end{popfigure}

\begin{exercice}[Complete the client form]
Réécoute (ou relis) le script du garage de Douala et complète la fiche client : Name (invente le nom de la cliente : \eng{Miss Ngono}), Service requested, Price, Time or deadline.
\end{exercice}

\begin{corrige}[Exercice 4]
\begin{itemize}
\item \textbf{Name :} Miss Ngono
\item \textbf{Service requested :} motorbike repair (changing the brake pads)
\item \textbf{Price :} 15,000 francs
\item \textbf{Time or deadline :} tomorrow at 4 p.m.
\end{itemize}
\end{corrige}

\courssub{Exercice 5 : Question de synthèse}

La question de synthèse oblige à comparer les informations des trois scripts, et non à recopier une seule phrase. C'est l'étape la plus difficile : tu dois rassembler les prix, les classer, puis rédiger une phrase de comparaison.

\begin{exoresolu}[Synthesis question]
\textbf{Énoncé.} Which service is the most expensive? Compare the three services and justify your answer.
\tcblower
\textbf{Solution détaillée.} On relève d'abord les trois prix : hotel room = 35,000 F per night ; bank deposit = 10,000 F ; motorbike repair = 15,000 F. On compare : 35,000 > 15,000 > 10,000.

Réponse rédigée : \eng{The hotel room is the most expensive service: it costs 35,000 francs a night. The motorbike repair is cheaper (15,000 francs), and the bank deposit is the cheapest (10,000 francs).} — « La chambre d'hôtel est le service le plus cher : 35 000 F la nuit. La réparation de la moto est moins chère (15 000 F) et le dépôt bancaire est le moins cher (10 000 F). »

Note l'emploi du superlatif \eng{the most expensive} (le plus cher) et de \eng{cheaper / the cheapest} (moins cher / le moins cher).
\end{exoresolu}

\begin{aretenir}[L'essentiel de la section]
\begin{itemize}
\item Vrai/faux : toujours \textbf{justifier} en citant la phrase du script.
\item QCM : repérer le mot-clé, retrouver la phrase, éliminer les distracteurs (souvent des chiffres venant d'un autre script).
\item Réponses courtes : identifier le mot interrogatif (\eng{how much} = prix, \eng{when} = moment) et répondre par une phrase complète.
\item Fiche client : écoute ciblée de quatre informations (nom, service, prix, délai).
\item Synthèse : comparer avec \eng{more... than}, \eng{the most...}, \eng{cheaper / the cheapest}.
\end{itemize}
\end{aretenir}

\begin{experience}[Correction collective en classe]
Par groupes de quatre, corrigez les exercices 1 à 5. Pour chaque réponse, un élève lit la phrase du script qui justifie la réponse à voix haute (\eng{“The script says: ...”}), et le groupe valide ou discute. Ensuite, chaque groupe invente une nouvelle affirmation vraie ou fausse sur les scripts et la soumet à la classe.
\end{experience}

Maintenant que tu sais vérifier ta compréhension d'un dialogue de service, la section suivante te donnera tout le lexique thématique pour demander toi-même un service en anglais.

\coursec{Mise en route et stratégies d'écoute}

\begin{experience}[Brainstorming : les services au quotidien]
En classe, avec ton professeur, dresse la liste des services que tu utilises ou que ta famille utilise régulièrement à Yaoundé, Douala, Buea ou Bamenda.
\begin{enumerate}
\item Où vas-tu pour retirer de l'argent ? Pour ouvrir un compte ?
\item Comment réserves-tu une chambre d'hôtel pour un parent qui vient de l'intérieur du pays ?
\item Que fais-tu quand ta moto ou ton téléphone est en panne ?
\item Comment achètes-tu du crédit ou un forfait internet pour ton téléphone ?
\end{enumerate}
Note le vocabulaire anglais qui te vient à l'esprit pour chaque situation. Ce brainstorming active tes connaissances antérieures et prépare l'écoute.
\end{experience}

\begin{methode}[Stratégies d'écoute : comprendre un dialogue de service]
Pour réussir une compréhension orale sur le thème « demander un service », suis ces étapes :
\begin{enumerate}
\item \textbf{Avant l'écoute (Prédiction) :} Lis le titre, observe les images ou les mots-clés donnés. Devine le lieu (banque, hôtel, garage, agence télécom), les interlocuteurs (client / employé) et le but de la conversation (retirer, réserver, réparer, recharger).
\item \textbf{Première écoute (Globalité) :} Ne t'arrête pas aux mots inconnus. Repère le \cle{thème} général et l'\cle{intention} de chaque locuteur : qui demande ? qui répond ? Le ton est-il poli ? Le service est-il accordé ou refusé ?
\item \textbf{Deuxième écoute (Détails) :} Concentre-toi sur les \cle{mots-clés} (noms de lieux, montants, numéros de chambre, noms de pièces, type de forfait), les \cle{chiffres} (prix, heure, numéro de compte, quantité de data) et les \cle{noms propres}.
\item \textbf{Troisième écoute (Vérification) :} Confirme tes réponses. Attention aux \cle{marqueurs de relation logique} (\eng{but, however, unfortunately, certainly}) qui signalent un changement ou une réponse négative polie.
\item \textbf{Après l'écoute :} Réemploie immédiatement le lexique noté pour résumer la situation à l'oral ou par écrit.
\end{enumerate}
\end{methode}

\begin{attention}[Pièges fréquents à l'écoute]
\begin{itemize}
\item \textbf{Les nombres :} \eng{fifteen} (15) vs \eng{fifty} (50) ; \eng{thirteen} (13) vs \eng{thirty} (30). Écoute l'accent tonique : sur le second syllabe pour les « -teen », sur le premier pour les « -ty ».
\item \textbf{La politesse masquée :} \eng{I'm afraid the network is down} = refus poli (le réseau est en panne), pas une peur réelle.
\item \textbf{L'inversion sujet-verbe dans les questions polies :} \eng{Could you...?}, \eng{Would you mind...?} — ne confonds pas avec l'affirmation.
\end{itemize}
\end{attention}

\coursec{Script d'écoute : At the service desk}

\begin{textebox}[Script d'écoute : At the service desk (Multi-service Centre)]
\textbf{Scene 1 : At the Bank Counter}
\eng{Clerk:} Good morning, Madam. How can I help you today?
\eng{Mrs Ngono:} Good morning. I'd like to make a withdrawal of 50,000 francs from my savings account, please. And I also need a new cheque book.
\eng{Clerk:} Certainly. Could you fill in this withdrawal slip and show me your ID card?
\eng{Mrs Ngono:} Of course. Here is my ID. ... Here is the slip.
\eng{Clerk:} Thank you. Would you like the money in 10,000-franc notes?
\eng{Mrs Ngono:} Yes, please. And could you transfer 20,000 francs to my sister's account at the same bank? Her account number is 334455.
\eng{Clerk:} No problem, Madam. I'll process the transfer now. Here is your receipt.

\textbf{Scene 2 : At the Hotel Reception}
\eng{Receptionist:} Good evening, Sir. Welcome to Hotel Mont Fébé. Do you have a reservation?
\eng{Mr Talla:} Good evening. Yes, I booked a double room for three nights under the name Talla. T-A-L-L-A.
\eng{Receptionist:} Let me check... Ah yes, Mr Talla. Double room, sea view, from today until Friday. That's 45,000 francs per night, breakfast included.
\eng{Mr Talla:} Perfect. I was wondering if late check-out is possible on Friday? My flight is at 6 p.m.
\eng{Receptionist:} I'm afraid check-out is at 12 noon, but we can keep your luggage safely until your departure.
\eng{Mr Talla:} That will be fine. Could I have the Wi-Fi password, please?
\eng{Receptionist:} Certainly. It's written on the key card envelope. Enjoy your stay!

\textbf{Scene 3 : At the Repair Shop}
\eng{Technician:} Good morning. What seems to be the problem with your motorbike?
\eng{Young Man:} Good morning. The engine makes a strange noise and it won't start easily. I think the spark plug is broken.
\eng{Technician:} Let me check... Yes, the spark plug needs replacing. And the brake pads are worn out too.
\eng{Young Man:} Can you fix both today?
\eng{Technician:} I have the spark plug in stock, but I need to order the brake pads. They'll arrive tomorrow morning.
\eng{Young Man:} Okay. Before you start, can you give me an estimate for the parts and labour?
\eng{Technician:} Sure. The spark plug is 2,000 francs, the brake pads 8,000, and labour 5,000. Total: 15,000 francs.
\eng{Young Man:} Agreed. Please go ahead.

\textbf{Scene 4 : At the Telecom Agency}
\eng{Agent:} Welcome to MTN. How may I assist you?
\eng{Student:} Hello. I'd like to buy airtime and a data bundle for my phone.
\eng{Agent:} Do you have a registered SIM card?
\eng{Student:} Yes, here it is. I want the 1,000 francs daily bundle.
\eng{Agent:} That gives you 500 MB for 24 hours. The network is a bit slow in this area today, but it should work.
\eng{Student:} Also, I want to make a complaint. Since yesterday, I cannot receive calls. People say my number is "not reachable".
\eng{Agent:} I see. Let me check your line... There seems to be a technical issue on the mast near your quarter. Our team is working on it. It should be fixed by tomorrow.
\eng{Student:} Thank you. I hope so. Here is 2,000 francs for the airtime and the bundle.
\eng{Agent:} Thank you. Here is your receipt. Dial *157\# to check your balance.
\end{textebox}

\coursec{Compréhension : questions variées}

\begin{exercice}[Compréhension de l'écoute : At the service desk]
Réponds aux questions d'après le script ci-dessus. Écoute ou relis attentivement.

\textbf{A. Vrai ou Faux ? Justifie par une phrase du texte.}
\begin{enumerate}
\item Mrs Ngono wants to deposit money at the bank.
\item Mr Talla has a reservation for a single room.
\item The technician has the brake pads in stock.
\item The student's SIM card is not registered.
\item The network problem will be fixed today.
\end{enumerate}

\textbf{B. Questions à choix multiples (QCM). Choisis la bonne réponse.}
\begin{enumerate}
\item How much money does Mrs Ngono withdraw?
\begin{itemize}
\item[a)] 20,000 francs
\item[b)] 50,000 francs
\item[c)] 70,000 francs
\end{itemize}
\item What time is check-out at Hotel Mont Fébé?
\begin{itemize}
\item[a)] 6 p.m.
\item[b)] 12 noon
\item[c)] 10 a.m.
\end{itemize}
\item How much is the total estimate for the motorbike repair?
\begin{itemize}
\item[a)] 10,000 francs
\item[b)] 13,000 francs
\item[c)] 15,000 francs
\end{itemize}
\item What data bundle does the student buy?
\begin{itemize}
\item[a)] 1 GB for a week
\item[b)] 500 MB for 24 hours
\item[c)] Unlimited for a month
\end{itemize}
\end{enumerate}

\textbf{C. Réponses courtes (en anglais).}
\begin{enumerate}
\item What two services does Mrs Ngono request at the bank?
\item Why can't Mr Talla have a late check-out?
\item What two parts does the motorbike need?
\item Why does the student make a complaint?
\end{enumerate}
\end{exercice}

\begin{corrige}[Exercice : Compréhension de l'écoute]
\textbf{A. Vrai ou Faux ?}
1. \textbf{False.} « I'd like to make a \textbf{withdrawal} of 50,000 francs... » (Elle veut faire un retrait, pas un dépôt).
2. \textbf{False.} « I booked a \textbf{double room}... » (Chambre double, pas simple).
3. \textbf{False.} « I need to \textbf{order} the brake pads. They'll arrive tomorrow morning. » (Il doit les commander, pas en stock).
4. \textbf{False.} « Do you have a \textbf{registered} SIM card? » — « Yes, here it is. » (Elle est enregistrée).
5. \textbf{False.} « It should be fixed by \textbf{tomorrow}. » (Pas aujourd'hui).

\textbf{B. QCM.}
1. b) 50,000 francs. 2. b) 12 noon. 3. c) 15,000 francs. 4. b) 500 MB for 24 hours.

\textbf{C. Réponses courtes.}
1. She requests a withdrawal and a transfer (and a new cheque book).
2. Because check-out is at 12 noon (policy).
3. A spark plug and brake pads.
4. Because he cannot receive calls (his number is "not reachable").
\end{corrige}

\coursec{Lexique thématique : requesting services}

Dans la section précédente, tu as répondu à des questions de compréhension sur le script \eng{At the service desk} : vrai ou faux, QCM, réponses courtes. Tu as donc déjà rencontré, en contexte, des mots comme \eng{withdrawal}, \eng{reservation} ou \eng{complaint}. Le moment est venu de les organiser, de les mémoriser et d'apprendre à les réemployer. Un bon lexique ne s'apprend pas mot par mot : il s'apprend par \cle{champs lexicaux}, c'est-à-dire par familles de mots liés à une même situation. Ici, quatre situations de la vie quotidienne : la banque, l'hôtel, les réparations, le téléphone et internet.

\courssub{Vue d'ensemble : la carte mentale du lexique}

Observe d'abord cette carte mentale. Elle résume tout le vocabulaire de la leçon. Recopie-la dans ton cahier et complète-la au fur et à mesure.

\begin{popfigure}
\begin{center}\tcbox[colback=popblue,colframe=popblue,coltext=white,arc=9pt,boxsep=4pt,left=6pt,right=6pt]{\bfseries\small SERVICES}\end{center}
\vspace{-2pt}
\begin{tcbraster}[raster columns=2,raster equal height=rows,raster column skip=6pt,raster row skip=6pt,enhanced,blankest]
\begin{tcolorbox}[enhanced,colback=white,colframe=poporange,boxrule=0.9pt,arc=3pt,title={Bank},fonttitle=\bfseries\scriptsize,coltitle=white,colbacktitle=poporange,left=4pt,right=4pt,top=2pt,bottom=2pt,boxsep=1pt,toptitle=1.5pt,bottomtitle=1.5pt,halign=left]
\begin{itemize}[nosep,leftmargin=*,itemsep=1pt,font=\scriptsize]
\item account deposit
\item withdrawal loan
\item cheque book transfer
\end{itemize}
\end{tcolorbox}
\begin{tcolorbox}[enhanced,colback=white,colframe=popgreen,boxrule=0.9pt,arc=3pt,title={Hotel},fonttitle=\bfseries\scriptsize,coltitle=white,colbacktitle=popgreen,left=4pt,right=4pt,top=2pt,bottom=2pt,boxsep=1pt,toptitle=1.5pt,bottomtitle=1.5pt,halign=left]
\begin{itemize}[nosep,leftmargin=*,itemsep=1pt,font=\scriptsize]
\item reservation single / double
\item receptionist bill
\item check in check out
\end{itemize}
\end{tcolorbox}
\begin{tcolorbox}[enhanced,colback=white,colframe=poppurple,boxrule=0.9pt,arc=3pt,title={Repairs},fonttitle=\bfseries\scriptsize,coltitle=white,colbacktitle=poppurple,left=4pt,right=4pt,top=2pt,bottom=2pt,boxsep=1pt,toptitle=1.5pt,bottomtitle=1.5pt,halign=left]
\begin{itemize}[nosep,leftmargin=*,itemsep=1pt,font=\scriptsize]
\item repair fix
\item mechanic spare part
\item estimate broken
\end{itemize}
\end{tcolorbox}
\begin{tcolorbox}[enhanced,colback=white,colframe=poppink,boxrule=0.9pt,arc=3pt,title={Phone \& Internet},fonttitle=\bfseries\scriptsize,coltitle=white,colbacktitle=poppink,left=4pt,right=4pt,top=2pt,bottom=2pt,boxsep=1pt,toptitle=1.5pt,bottomtitle=1.5pt,halign=left]
\begin{itemize}[nosep,leftmargin=*,itemsep=1pt,font=\scriptsize]
\item airtime data bundle
\item SIM card network
\item complaint
\end{itemize}
\end{tcolorbox}
\end{tcbraster}

\legende{Carte mentale : les quatre champs lexicaux de la demande de services.}
\end{popfigure}

\courssub{Champ 1 — At the bank (à la banque)}

\begin{definition}[Vocabulaire de la banque]
\begin{itemize}
\item \eng{account} (nom) : compte bancaire. \eng{to open an account} = ouvrir un compte. \eng{savings account} = compte d'épargne ; \eng{current account} = compte courant.
\item \eng{deposit} (nom et verbe) : dépôt ; déposer de l'argent. \eng{to make a deposit} = faire un dépôt.
\item \eng{withdrawal} (nom) : retrait d'argent. Le verbe est \eng{to withdraw} (irrégulier : \eng{withdraw -- withdrew -- withdrawn}).
\item \eng{loan} (nom) : prêt, emprunt. \eng{to apply for a loan} = demander un prêt.
\item \eng{cheque book} (nom) : chéquier, carnet de chèques. (Orthographe UK : \eng{cheque} ; US : \eng{check}).
\item \eng{transfer} (nom et verbe) : virement ; virer de l'argent. \eng{bank transfer} = virement bancaire.
\item \eng{slip} (nom) : bordereau, formulaire pré-imprimé. \eng{withdrawal slip}, \eng{pay-in slip}.
\item \eng{ID card} (nom) : carte d'identité. \eng{receipt} (nom) : reçu.
\end{itemize}
\end{definition}

\begin{exemplebox}[À la banque de Buea]
\eng{I'd like to withdraw some money from my account.} « Je voudrais retirer de l'argent de mon compte. »

\eng{Could you transfer 50,000 francs to my brother's account?} « Pourriez-vous virer 50 000 francs sur le compte de mon frère ? »

\eng{I need a new cheque book, please.} « J'ai besoin d'un nouveau chéquier, s'il vous plaît. »

\eng{Here is the receipt for your deposit.} « Voici le reçu pour votre dépôt. »
\end{exemplebox}

\begin{attention}[withdraw / withdrawal : attention à la famille de mots]
Le verbe \eng{to withdraw} est irrégulier : \eng{withdraw -- withdrew -- withdrawn}. Le nom est \eng{withdrawal} (avec un seul \eng{l} au milieu, deux à la fin). Ne dis jamais \eng{to do a withdrawal} : dis \eng{to make a withdrawal} ou simplement \eng{to withdraw money}. De même, on dit \eng{to make a deposit} (et non \eng{to do a deposit}). Attention à la prononciation : \eng{withdrawal} se prononce « wi-th-drô-al » (le \eng{w} initial ne se prononce pas).
\end{attention}

\courssub{Champ 2 — At the hotel (à l'hôtel)}

\begin{definition}[Vocabulaire de l'hôtel]
\begin{itemize}
\item \eng{reservation} (nom) : réservation. Verbe : \eng{to book a room} ou \eng{to reserve a room} = réserver une chambre. \eng{to make a reservation} = faire une réservation.
\item \eng{single room} : chambre simple (une personne) ; \eng{double room} : chambre double (deux personnes) ; \eng{twin room} : chambre à deux lits simples.
\item \eng{receptionist} (nom) : réceptionniste ; \eng{reception} = la réception (le lieu).
\item \eng{bill} (nom) : facture, note. \eng{Can I have the bill, please?} = L'addition / la facture, s'il vous plaît.
\item \eng{to check in} (verbe) : s'enregistrer à l'arrivée ; \eng{to check out} : régler sa note et partir. Noms : \eng{check-in}, \eng{check-out} (avec trait d'union).
\item \eng{key card} (nom) : carte magnétique (clé électronique).
\item \eng{luggage} (nom indénombrable) : bagages. Synonyme US : \eng{baggage}.
\end{itemize}
\end{definition}

\begin{exemplebox}[À l'hôtel de Kribi]
\eng{Good evening. I have a reservation for a double room.} « Bonsoir. J'ai une réservation pour une chambre double. »

\eng{What time is check-out?} « À quelle heure doit-on libérer la chambre ? »

\eng{Could I have the bill, please? I'm checking out this morning.} « Pourrais-je avoir la facture, s'il vous plaît ? Je pars ce matin. »

\eng{Can you keep my luggage until 6 p.m.?} « Pouvez-vous garder mes bagages jusqu'à 18 h ? »
\end{exemplebox}

\courssub{Champ 3 — Repairs (les réparations)}

\begin{definition}[Vocabulaire des réparations]
\begin{itemize}
\item \eng{repair} (verbe et nom) : réparer ; réparation. \eng{fix} (verbe) est le synonyme plus courant à l'oral (surtout pour objets, appareils, véhicules).
\item \eng{mechanic} (nom) : mécanicien (pour les véhicules). Prononciation : « mé-ka-nik » (\eng{/mɪˈkænɪk/}).
\item \eng{technician} (nom) : technicien (appareils électroniques, électroménager).
\item \eng{spare part} (nom) : pièce détachée. Pluriel : \eng{spare parts}.
\item \eng{estimate} (nom) : devis, estimation du prix. \eng{Can you give me an estimate?} = Pouvez-vous me faire un devis ? Verbe : \eng{to estimate}.
\item \eng{broken} (adjectif) : cassé, en panne, hors service. \eng{My motorbike is broken.} = Ma moto est en panne.
\item \eng{worn out} (adjectif) : usé. \eng{The brake pads are worn out.} = Les plaquettes de frein sont usées.
\item \eng{in stock} (locution) : en stock, disponible. \eng{out of stock} : en rupture de stock.
\item \eng{labour} (nom indénombrable) : main-d'œuvre (coût de la main-d'œuvre). US : \eng{labor}.
\end{itemize}
\end{definition}

\begin{exemplebox}[Chez le mécanicien de Douala]
\eng{Could you fix my phone? The screen is broken.} « Pourriez-vous réparer mon téléphone ? L'écran est cassé. »

\eng{Do you have the spare parts for this model?} « Avez-vous les pièces détachées pour ce modèle ? »

\eng{Before you start, can you give me an estimate?} « Avant de commencer, pouvez-vous me faire un devis ? »

\eng{The parts are in stock, but the labour will cost 5,000 francs.} « Les pièces sont en stock, mais la main-d'œuvre coûtera 5 000 francs. »
\end{exemplebox}

\courssub{Champ 4 — Phone and internet (téléphonie et internet)}

\begin{definition}[Vocabulaire du téléphone et d'internet]
\begin{itemize}
\item \eng{airtime} (nom indénombrable) : crédit téléphonique (pour appeler). \eng{to buy airtime} = acheter du crédit. \eng{top up} (verbe/nom) : recharger / recharge.
\item \eng{data bundle} (nom) : forfait internet (données mobiles). \eng{My data bundle is finished.} = Mon forfait est épuisé. \eng{to subscribe to a bundle} = souscrire un forfait.
\item \eng{SIM card} (nom) : carte SIM. \eng{to register a SIM card} = enregistrer une carte SIM (obligatoire au Cameroun).
\item \eng{network} (nom) : réseau. \eng{The network is down.} = Le réseau est en panne / ne passe pas. \eng{network coverage} = couverture réseau.
\item \eng{complaint} (nom) : réclamation. Verbe : \eng{to complain} = se plaindre, faire une réclamation. \eng{I'd like to make a complaint.} = Je voudrais faire une réclamation.
\item \eng{receipt} (nom) : reçu (preuve de paiement). \eng{balance} (nom) : solde (crédit restant). \eng{to check one's balance} = vérifier son solde.
\end{itemize}
\end{definition}

\begin{exemplebox}[À l'agence télécom de Bamenda]
\eng{I'd like to buy airtime and a data bundle, please.} « Je voudrais acheter du crédit et un forfait internet, s'il vous plaît. »

\eng{I'm afraid the network is down in our quarter.} « Je crains que le réseau ne soit en panne dans notre quartier. »

\eng{I want to make a complaint about my SIM card.} « Je veux faire une réclamation au sujet de ma carte SIM. »

\eng{Dial *157\# to check your balance.} « Composez le *157\# pour vérifier votre solde. »
\end{exemplebox}

\courssub{Formules de demande polie et formules de réponse}

En anglais, la politesse passe par des \cle{formules fixes} qu'il faut connaître par cœur. Plus la formule est longue, plus elle est polie. Au Cameroun, dans un contexte administratif ou commercial, le registre poli (\eng{polite}) est la norme.

\begin{propriete}[Demander poliment un service]
\begin{itemize}
\item \eng{I'd like + nom / to + V} (=\eng{I would like}) : Je voudrais... (poli, standard, le plus utilisé). \eng{I'd like to open an account.}
\item \eng{Could you + V (base verbale)...?} : Pourriez-vous...? (poli, demande d'action à l'interlocuteur). \eng{Could you check my balance?}
\item \eng{I was wondering if you could + V} : Je me demandais si vous pourriez... (très poli, indirect, hésitant).
\item \eng{Would you mind + V-ing (gérondif)...?} : Cela vous dérangerait-il de...? (le plus poli, formel). \textbf{Attention : le verbe est en \eng{-ing} !}
\end{itemize}
\textbf{Réponses types de l'employé :}
\begin{itemize}
\item \eng{Certainly.} / \eng{Of course.} / \eng{No problem.} (Accord).
\item \eng{Let me check.} / \eng{Just a moment, please.} (Attente / vérification).
\item \eng{I'm afraid + phrase} (Je crains que... : refus poli ou mauvaise nouvelle). \eng{I'm afraid the network is down.}
\item \eng{Would you mind waiting a moment?} (Demande d'attente polie).
\end{itemize}
\end{propriete}

\begin{savaistu}[La formule la plus polie de l'anglais : la logique inversée]
\eng{Would you mind + V-ing} est le sommet de la politesse anglaise. \eng{Would you mind checking my account?} = « Cela vous dérangerait-il de vérifier mon compte ? »
\textbf{Piège majeur :} La réponse attendue suit la logique anglaise (répondre à « est-ce que cela vous dérange ? »).
\begin{itemize}
\item \eng{No, not at all.} / \eng{No, I don't mind.} = « Non, cela ne me dérange pas » $\rightarrow$ \textbf{Je le fais volontiers (OUI)}.
\item \eng{Yes, I would mind.} / \eng{Yes, I do.} = « Oui, cela me dérange » $\rightarrow$ \textbf{Je refuse (NON)}.
\end{itemize}
Au français, on répond « Oui, je veux bien » à « Cela te dérange ? ». En anglais, on répond « No » pour accepter. Entraîne-toi à cette inversion !
\end{savaistu}

\begin{attention}[Faux amis : to demand et deception]
\eng{To demand} ne signifie pas « demander » mais « exiger avec autorité » : c'est très impoli, voire agressif ! Pour « demander poliment », utilise \eng{to request} (formel), \eng{to ask for} (neutre) ou les formules \eng{I'd like...} / \eng{Could you...?}.
\begin{center}
\begin{tabular}{|l|l|l|}
\hline
\rowcolor{popblueL} Anglais & Sens réel & Français correct pour « demander » \\
\hline
\eng{He demanded a refund.} & Il a \textbf{exigé} un remboursement. & \eng{He requested a refund.} / \eng{He asked for a refund.} \\
\hline
\end{tabular}
\end{center}

Autre piège classique : \eng{deception} signifie « tromperie, duperie, fraude », et non « déception ».
\begin{itemize}
\item \eng{Disappointment} = déception. \eng{To my great disappointment, the hotel was full.}
\item \eng{Deception} = tromperie. \eng{He was accused of deception.}
\end{itemize}
\end{attention}

\courssub{Texte de consolidation et tableau récapitulatif}

\begin{textebox}[A busy morning in Yaoundé]
Yesterday morning, Mrs Ngono had many things to do in town. First, she went to the bank on Kennedy Avenue. “Good morning,” she said to the clerk. “I'd like to make a withdrawal and order a new cheque book, please.” “Certainly, Madam. Could you fill in this form first?” the clerk replied. After that, Mrs Ngono walked to a small hotel to book a double room for her sister, who was arriving from Bafoussam. “I was wondering if you could keep the reservation until Friday,” she said. “No problem,” answered the receptionist with a smile. Then she took her old radio to a repair shop near the market. “It's broken. Can you fix it?” she asked. “Let me check... I'll need a spare part. I can give you an estimate tomorrow,” said the technician. Finally, she stopped at a phone shop to buy airtime and a data bundle. “I'm afraid the network is down at the moment,” the seller explained, “but your bundle will work as soon as the connection is back.” Mrs Ngono thanked him politely and went home, happy that all her requests had been answered.

\textbf{Glossaire} : clerk (nom) : employé de guichet ; to fill in (verbe) : remplir (un formulaire) ; repair shop (nom) : atelier de réparation ; as soon as (locution) : dès que ; connection (nom) : connexion (réseau).
\end{textebox}

Questions de vérification rapide : (1) What did Mrs Ngono do at the bank? (2) Why did she go to the hotel? (3) What was broken? (4) Why couldn't her data bundle work immediately?

\begin{center}
\begin{tabularx}{\linewidth}{|X|X|X|}
\hline
\rowcolor{popblueL} English & Français & Remarque \\
\hline
to request a service & demander un service & plus poli que \eng{to demand} \\
\hline
to fill in a form & remplir un formulaire & \eng{fill in} (GB) = \eng{fill out} (US) \\
\hline
airtime & crédit téléphonique & indénombrable : pas de pluriel \\
\hline
data bundle & forfait internet & \eng{a bundle} = un forfait / un paquet \\
\hline
spare part & pièce détachée & deux mots séparés ; pl. \eng{spare parts} \\
\hline
bill & facture, note & à l'hôtel et au restaurant \\
\hline
withdrawal & retrait (bancaire) & verbe : \eng{to withdraw} (irr.) \\
\hline
complaint & réclamation & verbe : \eng{to complain} \\
\hline
estimate & devis & \eng{to give an estimate} \\
\hline
check-in / check-out & enregistrement / départ & noms avec trait d'union \\
\hline
\end{tabularx}
\end{center}

\begin{exoresolu}[Complète avec le mot du champ lexical correct]
Complète chaque phrase avec un mot de la leçon : \eng{estimate}, \eng{airtime}, \eng{reservation}, \eng{withdrawal}, \eng{spare part}.

1. I don't have enough credit to call you; I need to buy some \ldots
2. The mechanic says my car needs a new \ldots for the engine.
3. She went to the bank to make a \ldots of 20,000 francs.
4. We have a \ldots for a single room at the Hilton Hotel.
5. Before repairing my laptop, please give me an \ldots.
\tcblower
\textbf{Solution détaillée.} 1. \eng{airtime} : on parle de crédit téléphonique (« credit to call »), indénombrable. 2. \eng{spare part} : le mécanicien remplace une pièce détachée du moteur. 3. \eng{withdrawal} : à la banque, « retirer » 20 000 francs = faire un retrait (nom). 4. \eng{reservation} : à l'hôtel, on a une réservation pour une chambre simple. 5. \eng{estimate} : avant une réparation, on demande un devis (estimation du coût).
\end{exoresolu}

\begin{exercice}[À toi de jouer : Politesse et reformulation]
Réécris ces demandes directes pour les rendre plus polies, en utilisant la formule indiquée entre parenthèses.

1. Give me the bill. (\eng{Could you...?})
2. I want to check out. (\eng{I'd like...})
3. Check my account. (\eng{Would you mind...?})
4. I want to make a complaint. (\eng{I was wondering if...})
5. Fix my motorbike today. (\eng{Could you...?})
6. I need a new SIM card. (\eng{I'd like...})
\end{exercice}

\begin{corrige}[Exercice : À toi de jouer]
1. \eng{Could you give me the bill, please?}
2. \eng{I'd like to check out, please.}
3. \eng{Would you mind checking my account?} (attention au \eng{-ing} après \eng{mind}).
4. \eng{I was wondering if I could make a complaint.}
5. \eng{Could you fix my motorbike today?}
6. \eng{I'd like to buy a new SIM card, please.} (ou \eng{I'd like a new SIM card.})
\end{corrige}

\coursec{Production guidée : je demande un service}

\begin{methode}[Rédiger un dialogue de demande de service]
Pour écrire un dialogue réaliste (examen, composition), suis ce canevas :
\begin{enumerate}
\item \textbf{Salutation polie} : \eng{Good morning / afternoon / evening.}
\item \textbf{Identification du besoin} : \eng{I'd like to...} / \eng{Could you...?} + action précise.
\item \textbf{Échange d'informations} : L'employé demande des détails (nom, numéro, montant, modèle). Le client répond.
\item \textbf{Problème éventuel} : \eng{I'm afraid...} (rupture de stock, panne réseau, horaire non dispo).
\item \textbf{Solution / Accord} : \eng{Let me check...} / \eng{Certainly.} / \eng{It will be ready by...}
\item \textbf{Remerciements et congé} : \eng{Thank you. Goodbye.}
\end{enumerate}
Utilise le vocabulaire des 4 champs lexicaux et au moins 2 formules de politesse différentes.
\end{methode}

\begin{experience}[Jeu de rôle : Au guichet unique (One-stop service)]
\textbf{Consigne :} En binôme. L'un joue le client, l'autre l'employé d'un centre multi-services (banque, poste, télécom, billetterie). Tire une carte « situation » au hasard. Joue le dialogue 2 minutes. Change de rôle.
\textbf{Cartes situations (exemples) :}
\begin{itemize}
\item Carte A : Tu veux virer 100 000 FCFA à ton oncle à Douala. Tu as ton ID et son numéro de compte.
\item Carte B : Tu veux réserver une chambre simple pour 2 nuits à l'hôtel de la mairie. Tu demandes le prix et le Wi-Fi.
\item Carte C : Ton téléphone ne charge plus. Tu demandes un devis pour changer le port de charge.
\item Carte D : Tu veux acheter 2000 F d'airtime et le forfait « 2 Go / 30 jours ». Tu signales que le réseau est lent chez toi.
\end{itemize}
\textbf{Évaluation par les pairs :} Le client a-t-il utilisé \eng{I'd like / Could you / Would you mind} ? L'employé a-t-il répondu \eng{Certainly / I'm afraid / Let me check} ?
\end{experience}

\begin{exercice}[Production écrite : Lettre de réclamation formelle]
Tu as acheté un téléphone neuf dans une grande surface à Yaoundé. Au bout de trois jours, l'écran se fend sans choc. Le vendeur refuse de l'échanger, disant que c'est de ta faute. Écris une lettre de réclamation (\eng{formal letter of complaint}) au directeur du magasin (120--150 mots).
\textbf{Structure imposée :}
\begin{enumerate}
\item Tes coordonnées (en haut à gauche) + Date.
\item Coordonnées du destinataire (Directeur, Nom du magasin, Adresse).
\item Objet : \eng{Subject: Complaint about defective phone / Request for replacement}.
\item Corps : 1. Faits (date d'achat, modèle, prix). 2. Problème (écran fissuré, usage normal). 3. Réponse insatisfaisante du vendeur. 4. Demande précise (échange ou remboursement) + référence à la loi / garantie.
\item Formule de politesse finale : \eng{Yours faithfully,} + Signature + Nom imprimé.
\end{enumerate}
\end{exercice}

\begin{corrige}[Exercice : Lettre de réclamation (Modèle)]
\begin{textebox}[Modèle de lettre formelle]
\textbf{Martin Ngono} \\
12, Rue de l'Université, Yaoundé \\
Cameroon \\
\texttt{martin.ngono@email.com} \\
+237 6XX XX XX XX \\

\textbf{15 October 2025} \\

\textbf{The Manager} \\
\textbf{TechWorld Superstore} \\
Boulevard de la Liberté, Yaoundé \\
Cameroon \\

\textbf{Subject: Complaint about defective smartphone -- Request for replacement} \\

Dear Sir or Madam, \\

I am writing to complain about a smartphone I purchased at your Yaoundé branch on 12 October 2025. The item is a \eng{TechPro X5}, 128 GB, black, bought for 185,000 FCFA (receipt n° 4455 attached). \\

Three days after purchase, a large crack appeared on the screen spontaneously, without any drop or impact on my part. I returned to the store on 15 October to request an exchange under the manufacturer's warranty. However, your sales assistant refused to help, claiming the damage was due to ``misuse'' and that no refund or replacement was possible. \\

I find this response unacceptable. Under Cameroonian consumer protection law, a product must be fit for purpose and durable. A screen cracking under normal use indicates a manufacturing defect. \\

Therefore, I formally request a full replacement of the device or a complete refund within 14 days. I have attached a copy of the receipt and photos of the defect. \\

I look forward to your prompt resolution of this matter. \\

Yours faithfully, \\

\textit{(Signature)} \\

\textbf{Martin Ngono}
\end{textebox}
\end{corrige}

\begin{aretenir}[L'essentiel de la leçon 16]
\begin{itemize}
\item \textbf{4 champs lexicaux :} Banque (\eng{account, deposit, withdrawal, loan, cheque book, transfer, slip, receipt}), Hôtel (\eng{reservation, single/double/twin room, receptionist, bill, check-in/out, key card, luggage}), Réparations (\eng{repair, fix, mechanic, technician, spare part, estimate, broken, worn out, in stock, labour}), Télécom (\eng{airtime, data bundle, SIM card, network, complaint, receipt, balance, top up}).
\item \textbf{Politesse (ordre croissant) :} \eng{I'd like...} < \eng{Could you...?} < \eng{I was wondering if...} < \eng{Would you mind + V-ing}.
\item \textbf{Réponses types :} \eng{Certainly.} / \eng{Let me check.} / \eng{I'm afraid...} (refus poli).
\item \textbf{Faux amis :} \eng{to demand} = exiger (impoli) $\neq$ demander ; \eng{deception} = tromperie $\neq$ déception (\eng{disappointment}).
\item \textbf{Grammaire clé :} \eng{Would you mind + V-ing} (gérondif) ; \eng{make a withdrawal / a deposit / a reservation / a complaint} (collocations avec \eng{make}) ; verbes irréguliers \eng{withdraw -- withdrew -- withdrawn}.
\end{itemize}
\end{aretenir}

\coursec{Production guidée : je demande un service}

Dans la section précédente, nous avons constitué le lexique thématique de la demande de services (\eng{request, service, charge, assistance, available, appointment, data bundle, airtime…}) et les formules de politesse qui l'accompagnent (\eng{I'd like…, Could you…?, I'm afraid…}). Il est maintenant temps de \cle{réemployer} ce vocabulaire dans une production orale et écrite authentique : demander un service, comme un client réel le ferait à Douala, à Buea ou à Londres. Cette section te guide pas à pas, en cinq étapes, du choix de la situation jusqu'au jeu de rôle complet.

\courssub{Observer avant de produire : un modèle de dialogue}

Avant de parler, observons un dialogue complet. Repère les \cle{étapes} de la conversation : salutation, demande polie, détails, acceptation ou refus, remerciement.

\begin{textebox}[At the mobile phone agency — a model dialogue]
\textbf{Employee:} Good morning, madam. Welcome to Camtel. How may I help you today?

\textbf{Customer:} Good morning. I'd like to buy a data bundle, please.

\textbf{Employee:} Certainly. We have daily, weekly and monthly bundles. Which one would you prefer?

\textbf{Customer:} Could you tell me how much the monthly bundle costs?

\textbf{Employee:} The monthly bundle is five thousand francs for ten gigabytes.

\textbf{Customer:} I'm afraid that's a little expensive for me. Do you have anything cheaper?

\textbf{Employee:} Yes, of course. The weekly bundle costs two thousand francs for three gigabytes.

\textbf{Customer:} That's perfect. I'd like the weekly bundle, please.

\textbf{Employee:} Very well. Could you give me your phone number, please?

\textbf{Customer:} It's six nine zero, twelve, thirty-four, fifty-six.

\textbf{Employee:} Thank you. The bundle is now active. Is there anything else I can do for you?

\textbf{Customer:} No, thank you. Actually, could you also check my airtime balance, please?

\textbf{Employee:} Certainly. You have one hundred and fifty francs of airtime left.

\textbf{Customer:} Oh, I'm afraid I don't have enough airtime. Could I top up one thousand francs, please?

\textbf{Employee:} Of course. Here is your receipt. Thank you for choosing Camtel.

\textbf{Customer:} Thank you very much for your help. Goodbye.

\textbf{Employee:} Goodbye, madam. Have a nice day!
\end{textebox}

\begin{definition}[Glossaire du dialogue]
\begin{itemize}
\item \eng{bundle} (n.) : forfait (de données ou de communication)
\item \eng{airtime} (n.) : crédit téléphonique
\item \eng{to top up} (v.) : recharger (son crédit)
\item \eng{balance} (n.) : solde
\item \eng{receipt} (n.) : reçu — prononcé « ri-site » (le \eng{p} ne se prononce pas !)
\item \eng{I'm afraid…} : « j'ai bien peur que… » — formule polie pour annoncer un problème
\end{itemize}
\end{definition}

\textbf{Questions d'observation.} \emph{Réponds en anglais, en phrases courtes.}
\begin{enumerate}
\item What service does the customer request first?
\item How much does the weekly bundle cost?
\item Why does the customer refuse the monthly bundle?
\item What second service does the customer ask for at the end?
\item Underline all the polite expressions in the dialogue. How many are there?
\end{enumerate}

\courssub{La structure du dialogue : cinq étapes obligatoires}

Tout dialogue de demande de service suit la même architecture. La mémoriser, c'est avoir un plan prêt pour n'importe quelle situation : banque, hôtel, garage, agence de téléphonie, agence de voyage.

\begin{popfigure}
\begin{center}
\begin{tikzpicture}[
node distance=0.55cm and 0.9cm,
etape/.style={rectangle, rounded corners=4pt, draw=popblue, fill=popblueL, text width=3.4cm, align=center, font=\small, inner sep=5pt},
fleche/.style={-{Stealth[length=2.5mm]}, thick, popdark}
]
\node[etape, align=center] (a){\textbf{1. Greeting}\\ \eng{Good morning!}};
\node[etape, right=of a, fill=popgreenL, draw=popgreen, align=center] (b){\textbf{2. Request}\\ \eng{I'd like…, please.}};
\node[etape, right=of b, fill=popgoldL, draw=popgold, align=center] (c){\textbf{3. Details}\\ \eng{How much…? When…?}};
\node[etape, below=of c, fill=poporangeL, draw=poporange, align=center] (d){\textbf{4. Acceptance / Refusal}\\ \eng{Certainly. / I'm afraid…}};
\node[etape, left=of d, fill=poppinkL, draw=poppink, align=center] (e){\textbf{5. Thanks \& closing}\\ \eng{Thank you. Goodbye.}};
\draw[fleche] (a) -- (b);
\draw[fleche] (b) -- (c);
\draw[fleche] (c) -- (d);
\draw[fleche] (d) -- (e);
\end{tikzpicture}
\end{center}
\legende{Organigramme du dialogue de demande de service : les cinq étapes dans l'ordre.}
\end{popfigure}

\begin{methode}[Réussir le jeu de rôle en cinq gestes]
\begin{enumerate}
\item \textbf{Salue d'abord.} On ne demande jamais un service sans dire \eng{Good morning} ou \eng{Hello}. En anglais, sauter la salutation est perçu comme impoli.
\item \textbf{Formule ta demande poliment} avec \eng{I'd like…}, \eng{Could I…?} ou \eng{Could you…?} — jamais \eng{I want…}, trop direct.
\item \textbf{Donne ou demande les détails} : prix (\eng{How much is it?}), horaires (\eng{What time do you open?}), quantité, numéro.
\item \textbf{Réagis} : accepte (\eng{That's perfect.}) ou négocie poliment (\eng{I'm afraid that's too expensive. Do you have anything cheaper?}).
\item \textbf{Remercie et prends congé} : \eng{Thank you for your help. Goodbye.}
\end{enumerate}
\end{methode}

\courssub{Étape 1 : choisir une situation}

Choisis \textbf{un} des quatre contextes suivants. Chacun appartient au chapitre \emph{Economic Life and Occupations} et mobilise le lexique de la leçon.

\begin{center}
\begin{tabularx}{\linewidth}{|l|X|X|}
\hline
\rowcolor{popblueL} \textbf{Situation} & \textbf{Services possibles (English)} & \textbf{Lexique utile} \\
\hline
At the bank & open an account, withdraw money, change money & \eng{account, deposit, withdrawal, form, signature} \\
\hline
At the hotel & book a room, ask for breakfast, request a wake-up call & \eng{single/double room, reservation, available, per night} \\
\hline
At the garage & repair a tyre, check the engine, change the oil & \eng{repair, flat tyre, engine, mechanic, bill} \\
\hline
At the phone agency & buy a data bundle, top up airtime, check the balance & \eng{bundle, airtime, balance, SIM card, receipt} \\
\hline
\end{tabularx}
\end{center}

\courssub{Étape 2 : préparer sa demande avec la formule polie adaptée}

\begin{propriete}[Les formules de demande, de la plus directe à la plus polie]
\begin{itemize}
\item \eng{I'd like + nom / to + verbe} : la formule standard, toujours correcte. \eng{I'd like to open an account, please.} « Je voudrais ouvrir un compte, s'il vous plaît. »
\item \eng{Could I + verbe…?} : demande polie de permission ou de service. \eng{Could I book a double room for two nights?} « Pourrais-je réserver une chambre double pour deux nuits ? »
\item \eng{Could you + verbe…?} : demande polie adressée à l'employé. \eng{Could you check my balance, please?} « Pourriez-vous vérifier mon solde, s'il vous plaît ? »
\item \eng{I need + nom / to + verbe} : exprime un besoin, un peu plus direct. \eng{I need to repair a flat tyre.} « J'ai besoin de réparer un pneu crevé. »
\end{itemize}
\textbf{Refus poli côté employé} : \eng{I'm afraid + phrase} ou \eng{I'm sorry, but…} + justification. \eng{I'm afraid we don't have any rooms available tonight.} « J'ai bien peur que nous n'ayons aucune chambre disponible ce soir. »
\end{propriete}

\begin{attention}[Le piège numéro un des francophones : oublier la politesse]
L'omission de \eng{please} et de \eng{thank you} est l'erreur la plus fréquente des apprenants francophones à l'oral. En français, le ton de la voix suffit souvent ; en anglais, les mots de politesse sont \textbf{obligatoires}. Compare :
\begin{itemize}
\item Incorrect ou impoli : \eng{Give me a data bundle.} (« Donnez-moi un forfait. » — ordre sec !)
\item Correct : \eng{I'd like a data bundle, please.}
\item Impoli : \eng{How much?} — Correct : \eng{Could you tell me how much it costs, please?}
\end{itemize}
Autre piège : ne traduis pas « je veux » par \eng{I want} : c'est trop direct, presque agressif. Dis toujours \eng{I'd like} (= \eng{I would like}).
\end{attention}

\courssub{Étape 3 : rédiger cinq phrases avec le lexique de la leçon}

Avant de jouer le dialogue, rédige par écrit \textbf{cinq phrases} qui réemploient le vocabulaire de la section précédente. Cette préparation écrite garantit la fluidité de l'oral.

\begin{exemplebox}[Exemples de production guidée — situation « agence de téléphonie »]
\begin{enumerate}
\item \eng{Good morning. I'd like to buy a data bundle, please.} « Bonjour. Je voudrais acheter un forfait internet, s'il vous plaît. »
\item \eng{Could you tell me how much it costs?} « Pourriez-vous me dire combien cela coûte ? »
\item \eng{I'm afraid I don't have enough airtime.} « J'ai bien peur de ne pas avoir assez de crédit. »
\item \eng{Could I top up one thousand francs, please?} « Pourrais-je recharger mille francs, s'il vous plaît ? »
\item \eng{Thank you for your help. Goodbye.} « Merci pour votre aide. Au revoir. »
\end{enumerate}
\end{exemplebox}

\begin{exoresolu}[Rédiger les cinq phrases — situation « hôtel »]
Énoncé : Tu arrives dans un hôtel de Buea. Rédige cinq phrases pour demander une chambre, en utilisant : \eng{I'd like, available, Could you…?, per night, thank you}.
\tcblower
\textbf{Solution détaillée.}
\begin{enumerate}
\item Salutation + demande : \eng{Good evening. I'd like to book a single room, please.}
\item Question sur la disponibilité : \eng{Do you have any rooms available for tonight?}
\item Question sur le prix : \eng{Could you tell me how much a single room costs per night?}
\item Réaction : \eng{That's fine. I'll take it for two nights.}
\item Clôture : \eng{Thank you very much for your help. Goodbye.}
\end{enumerate}
Chaque phrase contient un mot du lexique imposé, une formule polie, et suit l'ordre de l'organigramme : salutation, demande, détails, décision, remerciement.
\end{exoresolu}

\courssub{Étapes 4 et 5 : jouer le dialogue en binôme}

\begin{experience}[Jeu de rôle : client / employé]
Travaillez par deux. L'élève A est le \cle{client}, l'élève B est l'\cle{employé}.
\begin{enumerate}
\item Choisissez ensemble une situation (banque, hôtel, garage, agence de téléphonie).
\item Le client prépare ses cinq phrases par écrit ; l'employé prépare deux réponses d'acceptation (\eng{Certainly. / Of course.}) et un refus justifié (\eng{I'm afraid… because…}).
\item Jouez le dialogue \textbf{sans regarder vos notes} autant que possible. L'employé doit refuser poliment au moins une demande, avec une justification (\eng{I'm afraid the mechanic is busy this morning. Could you come back at two o'clock?}).
\item Le client réagit au refus : négocier (\eng{Do you have anything cheaper?}) ou accepter une alternative (\eng{That would be fine, thank you.}).
\item Échangez les rôles et changez de situation.
\end{enumerate}
\end{experience}

\begin{exemplebox}[Refuser poliment avec justification — modèles côté employé]
\begin{itemize}
\item \eng{I'm afraid that service is not available today. The system is down.} « J'ai bien peur que ce service ne soit pas disponible aujourd'hui. Le système est en panne. »
\item \eng{I'm sorry, but all our rooms are fully booked this weekend.} « Je suis désolé(e), mais toutes nos chambres sont réservées ce week-end. »
\item \eng{I'm afraid we don't repair this type of engine here.} « J'ai bien peur que nous ne réparions pas ce type de moteur ici. »
\end{itemize}
Remarque : après \eng{I'm afraid}, on utilise une phrase complète (sujet + verbe), jamais un nom seul.
\end{exemplebox}

\begin{attention}[Erreurs fréquentes à éviter pendant le jeu de rôle]
\begin{itemize}
\item Ordre des mots dans les questions indirectes : \eng{Could you tell me how much it costs?} et non \eng{how much does it cost} après \eng{Could you tell me}.
\item \eng{How much} pour le prix (indénombrable) : \eng{How much is it?} — pas \eng{How many is it?}
\item Ne dis pas \eng{I am agree} : « je suis d'accord » se dit \eng{I agree} (calque syntaxique du français « je suis d'accord »).
\item \eng{Actually} ne veut pas dire « actuellement » mais « en fait ». Pour « actuellement », dis \eng{at the moment}.
\end{itemize}
\end{attention}

\courssub{Grille d'auto-évaluation}

Après le jeu de rôle, évalue-toi honnêtement (ou évalue ton binôme) avec cette grille. Mets une croix dans la colonne correspondante.

\begin{center}
\begin{tabularx}{\linewidth}{|X|c|c|c|}
\hline
\rowcolor{popgreenL} \textbf{Critère} & \textbf{Oui} & \textbf{Un peu} & \textbf{Non} \\
\hline
J'ai salué au début et pris congé à la fin (\eng{Good morning… Goodbye}) & & & \\
\hline
J'ai utilisé \eng{please} et \eng{thank you} au moins deux fois & & & \\
\hline
J'ai formulé ma demande avec \eng{I'd like…} ou \eng{Could you…?} & & & \\
\hline
J'ai réemployé au moins quatre mots du lexique de la leçon & & & \\
\hline
J'ai donné ou demandé des détails (prix, horaire, quantité) & & & \\
\hline
Le refus (côté employé) était poli et justifié (\eng{I'm afraid… because…}) & & & \\
\hline
Ma prononciation était claire et mon débit naturel & & & \\
\hline
\end{tabularx}
\end{center}

\begin{aretenir}[L'essentiel de la production guidée]
\begin{itemize}
\item Un dialogue de service suit toujours cinq étapes : \textbf{salutation, demande polie, détails, acceptation ou refus, remerciement}.
\item Les formules clés : \eng{I'd like…, Could I…?, Could you…?} pour demander ; \eng{I'm afraid… / I'm sorry, but…} pour refuser poliment.
\item \eng{Please} et \eng{thank you} ne sont jamais optionnels en anglais : c'est la marque de politesse la plus importante.
\item Préparer cinq phrases par écrit avant de parler est la meilleure stratégie pour gagner en fluidité et en confiance.
\end{itemize}
\end{aretenir}

\begin{exercice}[À toi de jouer]
\begin{enumerate}
\item Choisis la situation « garage ». Rédige cinq phrases du client en utilisant : \eng{I'd like, Could you…?, repair, How much, thank you}.
\item Rédige deux réponses de l'employé : une acceptation et un refus poli justifié.
\item Joue le dialogue avec ton voisin, puis remplis la grille d'auto-évaluation.
\end{enumerate}
\end{exercice}

\begin{corrige}[Exercice « À toi de jouer »]
\textbf{1. Phrases du client (modèle) :}
\eng{Good morning. I'd like to repair a flat tyre, please. / Could you also check the engine, please? / How much will the repair cost? / I'm afraid that's too expensive. Could you do it cheaper? / Thank you for your help. Goodbye.}

\textbf{2. Réponses de l'employé (modèle) :}
Acceptation : \eng{Certainly, sir. Bring your car in and we will repair it this afternoon.}
Refus justifié : \eng{I'm afraid the mechanic is busy until three o'clock. Could you come back this evening?}

\textbf{3. Grille :} vérifie en particulier la présence de \eng{please} et \eng{thank you}, l'utilisation de \eng{I'd like / Could you}, et la justification après \eng{I'm afraid}.
\end{corrige}

\newpage

\coursec{Activité d'intégration}

\coursec{Lesson 16 — Listening: Texts about requesting for services}

\courssub{Mise en route et stratégies d'écoute}

\begin{methode}[Avant l'écoute : Se préparer efficacement]
Pour réussir une activité de compréhension de l'oral (\eng{Listening}), suivez ces trois étapes clés :
\begin{enumerate}
    \item \textbf{Analyser le contexte et les consignes} : Lisez le titre, regardez les images éventuelles, identifiez le type de document (dialogue au guichet, appel téléphonique, annonce publique) et la nature des questions (Vrai/Faux, QCM, complétion, réponse courte). Cela active votre \cle{connaissance du monde} (\eng{world knowledge}) sur le thème des services.
    \item \textbf{Prédire le lexique et les structures} : Anticipez les mots-clés liés au thème (ex. \eng{transfer, booking, quote, subscription, fees, reference number}) et les formules de politesse (\eng{Could you..., I would like to..., How much..., How long...}).
    \item \textbf{Repérer les indices sonores} : Pendant l'écoute, soyez attentif à l'intonation (montante pour les questions, descendante pour les affirmations), aux temps forts (\eng{stress}) sur les mots porteurs de sens (noms, verbes, chiffres, noms propres) et aux marqueurs de relation (\eng{however, unfortunately, certainly, right away}).
\end{enumerate}
\end{methode}

\begin{aretenir}[Stratégies pendant l'écoute (Double écoute standard)]
\begin{itemize}
    \item \textbf{Première écoute (Globale) :} Ne cherchez pas à tout comprendre. Identifiez : \textbf{Qui parle ?} (Client / Employé), \textbf{Où ?} (Banque, Hôtel, Garage, Agence télécom), \textbf{Quel service ?} (Virement, Réservation, Réparation, Forfait). Cochez les cases correspondantes sur votre fiche.
    \item \textbf{Deuxième écoute (Détaillée) :} Concentrez-vous sur les \textbf{informations chiffrées} (prix, dates, numéros, durées) et les \textbf{mots-clés} exacts. Notez-les en chiffres ou en abréviations (ex. \eng{2,500 FCFA}, \eng{48h}, \eng{Ref: 8842}). Ne réécrivez pas des phrases complètes.
    \item \textbf{Troisième écoute (si dispo / Vérification) :} Validez vos réponses, vérifiez les pièges (négation, correction du locuteur : \eng{Actually, it's 3,000 not 2,000}).
\end{itemize}
\end{aretenir}

\courssub{Script d'écoute : At the service desk}

\begin{textebox}[Script d'écoute original — 4 situations de service (Enregistrements simulés)]
\textbf{Situation 1 : Agence MTN — Bamenda (Quarter Mile 3)}
\eng{Agent:} Good morning. Welcome to MTN Mile 3. How may I assist you?
\eng{Client:} Good morning. I would like to migrate my number to the 4G network and subscribe to the MTN Pulse bundle. Could you please tell me how much the migration costs and when the service will be active?
\eng{Agent:} Certainly. The migration to 4G is free of charge. The Pulse bundle with 5 GB of data valid for 30 days costs 2,500 FCFA per month. Activation is usually instant once the payment is confirmed. Your transaction reference is MTN-BDA-8842. Please dial *157\# to confirm.
\eng{Client:} Thank you very much. Goodbye.
\eng{Agent:} You're welcome. Enjoy your day.

\vspace{0.5cm}
\textbf{Situation 2 : Hôtel La Plage — Kribi (Appel téléphonique)}
\eng{Receptionist:} Hotel La Plage, good afternoon. How can I help you?
\eng{Client:} Good afternoon. I would like to book a double room with a sea view for the weekend of July 14th to 16th, that is two nights. Could you tell me the total price including breakfast?
\eng{Receptionist:} Of course. A double room with sea view is 22,500 FCFA per night, breakfast included. The total for two nights comes to 45,000 FCFA. We require a 50\% deposit to secure the booking.
\eng{Client:} And what is your cancellation policy?
\eng{Receptionist:} Free cancellation up to 48 hours before check-in. After that, the first night is charged. Your booking reference is HOT-KRI-559.
\eng{Client:} Perfect. I will send the deposit today. Thank you.
\eng{Receptionist:} You're welcome, sir. We look forward to welcoming you.

\vspace{0.5cm}
\textbf{Situation 3 : Agence UBA — Yaoundé (Carrefour Warda)}
\eng{Teller:} Good morning, sir. Welcome to UBA. What can I do for you today?
\eng{Client:} Good morning. I need to make an urgent SWIFT transfer of 250,000 FCFA to a beneficiary in France for tuition fees. Could you tell me the transfer charges and the processing time?
\eng{Teller:} Right away. The flat transfer fee is 12,500 FCFA. Please note that correspondent bank charges may apply and are deducted from the amount received. The standard processing time is 2 to 3 working days.
\eng{Client:} I see. And the exchange rate applied today?
\eng{Teller:} The current rate is 1 Euro = 655.957 FCFA. Your transaction reference is SWF-YDE-1102. Here is your receipt.
\eng{Client:} Thank you. Have a nice day.
\eng{Teller:} You too, sir.

\vspace{0.5cm}
\textbf{Situation 4 : Garage Auto Rapide — Douala (Akwa)}
\eng{Mechanic:} Good morning. Auto Rapide, how can I help?
\eng{Client:} Good morning. My Toyota Corolla 2015 makes a grinding noise when I brake. I would like a quote for replacing the front brake pads and discs. How long will the repair take?
\eng{Mechanic:} We need to inspect the car first, but for a Corolla 2015, genuine pads and discs cost around 35,000 FCFA. Labour is 15,000 FCFA. The total estimate is 50,000 FCFA.
\eng{Client:} That sounds reasonable. When will the car be ready?
\eng{Mechanic:} If you bring it in now, it will take about 3 hours. It should be ready by 2 PM. We offer a 6-month warranty on parts and labour.
\eng{Client:} Excellent. I'll bring it right away. Thank you.
\eng{Mechanic:} You're welcome. See you soon.
\end{textebox}

\courssub{Compréhension de l'oral : Exercices variés}

\begin{exercice}[Exercice 1 : Vrai ou Faux — Justifiez par une phrase du texte]
Écoutez les quatre situations et dites si les affirmations sont \eng{True} ou \eng{False}. Citez la phrase exacte (en anglais) qui justifie votre réponse.
\begin{enumerate}
    \item The migration to 4G at MTN costs 2,500 FCFA.
    \item The hotel in Kribi requires a full payment in advance to confirm the booking.
    \item The SWIFT transfer at UBA will arrive the same day.
    \item The garage in Douala guarantees the repair for one year.
    \item In all four situations, the client receives a reference number.
\end{enumerate}
\end{exercice}

\begin{exercice}[Exercice 2 : QCM — Choisissez la bonne réponse]
\begin{enumerate}
    \item What is the monthly cost of the MTN Pulse plan (5 GB)?
    \begin{itemize}
        \item[A)] Free of charge
        \item[B)] 2,500 FCFA
        \item[C)] 5,000 FCFA
        \item[D)] 22,500 FCFA
    \end{itemize}
    \item What does "B\&B" stand for in the hotel context?
    \begin{itemize}
        \item[A)] Breakfast and Bed
        \item[B)] Bed and Breakfast
        \item[C)] Board and Breakfast
        \item[D)] Breakfast and Booking
    \end{itemize}
    \item How long does the SWIFT transfer take according to the UBA teller?
    \begin{itemize}
        \item[A)] Instant
        \item[B)] 24 hours
        \item[C)] 2 to 3 working days
        \item[D)] One week
    \end{itemize}
    \item What is the total estimated cost for the brake repair at Auto Rapide?
    \begin{itemize}
        \item[A)] 35,000 FCFA
        \item[B)] 15,000 FCFA
        \item[C)] 50,000 FCFA
        \item[D)] 65,000 FCFA
    \end{itemize}
\end{enumerate}
\end{exercice}

\begin{exercice}[Exercice 3 : Complétion de fiche — Informations chiffrées]
Complétez le tableau ci-dessous avec les chiffres exacts entendus (prix, délais, numéros).
\begin{center}
\begin{tabularx}{\linewidth}{|X|c|c|c|}\hline
\rowcolor{popblueL} \textbf{Service} & \textbf{Prix / Coût} & \textbf{Délai / Durée} & \textbf{N° de référence} \\ \hline
MTN Migration + Pulse & & Activation: \dots & MTN-BDA-\dots \\ \hline
Hôtel La Plage (2 nuits) & Total: \dots FCFA & Séjour: \dots nuits & HOT-KRI-\dots \\ \hline
UBA Virement SWIFT & Frais: \dots FCFA & \dots jours ouvrés & SWF-YDE-\dots \\ \hline
Garage Auto Rapide & Total: \dots FCFA & Réparation: \dots heures & (Non communiqué) \\ \hline
\end{tabularx}
\end{center}
\end{exercice}

\begin{exercice}[Exercice 4 : Réponses courtes — Questions indirectes]
Répondez en anglais par une phrase complète en reprenant la structure de la question indirecte.
\begin{enumerate}
    \item What does the client ask the MTN agent about the cost? (\eng{The client asks...})
    \item What does the client want to know from the hotel receptionist about cancellation?
    \item What information does the client request from the UBA teller regarding time?
    \item What two things does the client ask the mechanic?
\end{enumerate}
\end{exercice}

\begin{corrige}[Exercices 1 à 4]
\textbf{Exercice 1 :}
\begin{enumerate}
    \item \textbf{False.} \eng{« The migration to 4G is free of charge. »} (The 2,500 FCFA is for the monthly Pulse plan).
    \item \textbf{False.} \eng{« We require a 50\% deposit to secure the booking. »} (Not full payment).
    \item \textbf{False.} \eng{« The standard processing time is 2 to 3 working days. »}
    \item \textbf{False.} \eng{« We offer a 6-month warranty on parts and labour. »}
    \item \textbf{True.} MTN: \eng{MTN-BDA-8842}; Hotel: \eng{HOT-KRI-559}; UBA: \eng{SWF-YDE-1102}. (Garage: no ref number given).
\end{enumerate}

\textbf{Exercice 2 :}
1. B \quad 2. B \quad 3. C \quad 4. C

\textbf{Exercice 3 :}
\begin{center}
\begin{tabularx}{\linewidth}{|X|c|c|c|}\hline
\rowcolor{popblueL} \textbf{Service} & \textbf{Prix / Coût} & \textbf{Délai / Durée} & \textbf{N° de référence} \\ \hline
MTN Migration + Pulse & 2,500 FCFA/mois & Instant / Immédiat & MTN-BDA-8842 \\ \hline
Hôtel La Plage (2 nuits) & 45,000 FCFA (total) & 2 nuits (14 au 16 juil.) & HOT-KRI-559 \\ \hline
UBA Virement SWIFT & 12,500 FCFA (+ frais corresp.) & 2 à 3 jours ouvrés & SWF-YDE-1102 \\ \hline
Garage Auto Rapide & 50,000 FCFA (estime) & 3 heures (prêt 14h) & -- \\ \hline
\end{tabularx}
\end{center}

\textbf{Exercice 4 (Modèles de réponses) :}
\begin{enumerate}
    \item The client asks \textbf{could you please tell me how much the migration costs} (and when the service will be active).
    \item The client wants to know \textbf{what the cancellation policy is}.
    \item The client requests \textbf{the processing time} (or \textbf{how long the transfer will take}).
    \item The client asks \textbf{for a quote for front brake pads and discs} and \textbf{how long the repair will take}.
\end{enumerate}
\end{corrige}

\courssub{Lexique thématique : Requesting Services}

\begin{definition}[Vocabulaire essentiel par secteur d'activité (Mots-clés)]
\begin{center}
\begin{tabularx}{\linewidth}{|l|X|}\hline
\rowcolor{popblueL} \textbf{Secteur / Scénario} & \textbf{Vocabulaire (English — Français)} \\ \hline
\textbf{Télécommunications} & \eng{to migrate (to 4G/5G)} (migrer), \eng{a subscription / a plan / a bundle} (un forfait), \eng{data allowance} (volume data), \eng{activation fee} (frais d'activation), \eng{SIM card} (carte SIM), \eng{network coverage} (couverture réseau), \eng{reference number / transaction ID} (numéro de référence), \eng{to top up / to recharge} (recharger). \\ \hline
\textbf{Hôtellerie / Hébergement} & \eng{to book / to make a reservation} (réserver), \eng{a double / twin / single room} (chambre double / lits jumeaux / simple), \eng{sea view / garden view} (vue mer / jardin), \eng{full board / half board / B\&B (Bed \& Breakfast)} (pension complète / demi-pension / nuit+pdj), \eng{check-in / check-out} (arrivée / départ), \eng{cancellation policy} (conditions d'annulation), \eng{deposit / down payment} (acompte / arrhes), \eng{amenities} (équipements/confort). \\ \hline
\textbf{Banque / Finance} & \eng{a wire transfer / SWIFT transfer} (virement international), \eng{beneficiary} (bénéficiaire), \eng{transfer fees / charges / commission} (frais de virement / commission), \eng{exchange rate} (taux de change), \eng{processing time / lead time} (délai de traitement), \eng{receipt / proof of transfer / statement} (reçu / preuve / relevé), \eng{account number / IBAN / SWIFT code (BIC)} (numéro de compte / IBAN / code BIC), \eng{correspondent bank} (banque correspondante). \\ \hline
\textbf{Garage / Réparation auto} & \eng{a quote / an estimate} (un devis), \eng{brake pads / brake discs (rotors US)} (plaquettes / disques de frein), \eng{labour cost / labour charge} (coût de la main-d'œuvre), \eng{spare parts / genuine parts} (pièces détachées / pièces d'origine), \eng{breakdown / fault} (panne / défaut), \eng{to service / to carry out a service} (faire la révision), \eng{warranty / guarantee} (garantie), \eng{roadworthy} (en état de circuler). \\ \hline
\end{tabularx}
\end{center}
\end{definition}

\begin{exemplebox}[Verbes et collocations utiles (Action verbs)]
\begin{itemize}
    \item \eng{To request a service} (demander un service) — plus formel que \eng{ask for}.
    \item \eng{To inquire about} (se renseigner sur) — \eng{I am writing to inquire about your rates.}
    \item \eng{To subscribe to} (s'abonner à) — \eng{subscribe to a plan}.
    \item \eng{To transfer money to sb} (virer de l'argent à qqn).
    \item \eng{To book sth for sb} (réserver qqch pour qqn).
    \item \eng{To inspect / To diagnose} (inspecter / diagnostiquer).
    \item \eng{To confirm a booking / an appointment} (confirmer une réservation / un RDV).
    \item \eng{To process a request / a payment} (traiter une demande / un paiement).
\end{itemize}
\end{exemplebox}

\courssub{Point de grammaire : Politesse et questions indirectes}

\begin{propriete}[Les modaux de politesse : Degrés de formalité]
\begin{center}
\begin{tabularx}{\linewidth}{|l|X|X|}\hline
\rowcolor{popblueL} \textbf{Modal} & \textbf{Niveau / Fonction} & \textbf{Exemple traduit} \\ \hline
\eng{Can I / Can you} & Standard, neutre, courant & \eng{Can I have the receipt?} « Puis-je avoir le reçu ? » \\ \hline
\eng{Could I / Could you} & \cle{Plus poli}, recommandé guichet/téléphone & \eng{Could you check the price, please?} « Pourriez-vous vérifier le prix ? » \\ \hline
\eng{May I} & Très formel, demande de permission & \eng{May I see your ID, please?} « Puis-je voir votre pièce d'identité ? » \\ \hline
\eng{Would you like} & Proposition (Employé $\to$ Client) & \eng{Would you like a receipt?} « Voulez-vous un reçu ? » \\ \hline
\eng{I would like (I'd like)} & Expression du désir/besoin (Client) & \eng{I'd like to book a room.} « Je voudrais réserver une chambre. » \\ \hline
\eng{I would appreciate it if you could} & Très formel, écrit (email/lettre) & \eng{I would appreciate it if you could confirm.} « Je vous saurais gré de confirmer. » \\ \hline
\end{tabularx}
\end{center}

\textbf{Règles de forme importantes :}
\begin{itemize}
    \item \eng{Can / Could / May / Would} + \textbf{base verbale (sans \eng{to})}. \eng{*Can you to help?} $\rightarrow$ \eng{Can you help?}
    \item \eng{Would like} + \textbf{\eng{to} + base verbale} (infinitif). \eng{*I would like have} $\rightarrow$ \eng{I would like \textbf{to} have}.
    \item \eng{Would like} s'utilise pour les \textbf{désirs ponctuels} (pas pour les goûts permanents : \eng{I like coffee} $\neq$ \eng{I'd like a coffee}).
\end{itemize}
\end{propriete}

\begin{propriete}[Les questions indirectes (\eng{Indirect Questions}) — Structure clé]
Pour être plus poli, on « emboîte » la question dans une phrase introductive. \textbf{L'ordre des mots redevient l'ordre affirmatif (Sujet + Verbe). L'auxiliaire \eng{do/does/did} disparaît.}
\begin{center}
\begin{tabularx}{\linewidth}{|X|X|}\hline
\rowcolor{popblueL} \textbf{Question directe (Direct)} & \textbf{Question indirecte (Indirect)} \\ \hline
\eng{How much does it cost?} & \eng{Could you tell me \textbf{how much it costs}?} \\ \hline
\eng{When will it be ready?} & \eng{I'd like to know \textbf{when it will be ready}.} \\ \hline
\eng{What is the exchange rate?} & \eng{Can you tell me \textbf{what the exchange rate is}?} \\ \hline
\eng{Is the migration free?} & \eng{Do you know \textbf{if / whether the migration is free}?} \\ \hline
\end{tabularx}
\end{center}

\end{propriete}

\begin{attention}[Pièges classiques des francophones — Questions indirectes]
\begin{itemize}
    \item \textbf{Erreur :} \eng{*Could you tell me how much \textbf{does it cost}?} $\rightarrow$ \textbf{Correct :} \eng{how much \textbf{it costs}} (Pas d'inversion, pas de \eng{do}).
    \item \textbf{Erreur :} \eng{*I'd like to know when \textbf{will it be} ready?} $\rightarrow$ \textbf{Correct :} \eng{when \textbf{it will be} ready}.
    \item \textbf{Erreur :} \eng{*Can you tell me where \textbf{is} the garage?} $\rightarrow$ \textbf{Correct :} \eng{where \textbf{the garage is}}.
    \item \textbf{Point d'interrogation :} Si la phrase introductive est une question (\eng{Could you tell me...?}), on met un \textbf{?} à la fin. Si c'est une affirmation (\eng{I'd like to know...}), on met un \textbf{.}.
\end{itemize}
\end{attention}

\begin{propriete}[Accepter, Refuser, Informer (Côté Employé / Fournisseur)]
\begin{center}
\begin{tabularx}{\linewidth}{|l|X|}\hline
\rowcolor{popblueL} \textbf{Fonction} & \textbf{Formules types (English)} \\ \hline
\textbf{Accueillir / Proposer aide} & \eng{Good morning. How can I help you? / What can I do for you? / How may I assist you?} \\ \hline
\textbf{Accepter / Confirmer} & \eng{Certainly. / Of course. / Right away. / Absolutely. / I'll check that for you. / No problem.} \\ \hline
\textbf{Annoncer un chiffre / info} & \eng{The total is... / It costs... / The fee is... / It takes [duration]. / Your reference is...} \\ \hline
\textbf{Refuser / Problème (Politesse)} & \eng{I'm afraid (that)... / Unfortunately,... / I'm sorry, but... / That's not possible today. / The system is down.} \\ \hline
\textbf{Clôturer} & \eng{Thank you for choosing [Company]. / Have a nice day. / You're welcome. / Is there anything else?} \\ \hline
\end{tabularx}
\end{center}
\end{propriete}

\courssub{Production guidée : Je demande un service}

\begin{experience}[Jeu de rôle : « Guichets de services au Cameroun » — 45 min]
\textbf{Objectif :} Interagir oralement en situation authentique (Client $\leftrightarrow$ Employé) pour demander un service, obtenir des informations chiffrées (prix, délai, référence) et les noter.

\textbf{Matériel :} 4 \cle{Cartes-situations} (voir ci-dessous), 1 \cle{Fiche Client} par élève auditeur.

\end{experience}

\begin{methode}[Déroulement]
\textbf{Étape 1 : Préparation individuelle (5 min)}
\begin{enumerate}
    \item Piochez une \cle{carte-situation} au hasard. Lisez-la attentivement.
    \item Selon votre rôle attribué (A ou B), préparez vos répliques sur brouillon.
    \item \textbf{Rôle A (Client) :} Rédigez votre demande initiale avec \textbf{au moins deux formules de politesse différentes} (ex. \eng{I would like to...}, \eng{Could you please...}). Préparez \textbf{deux questions précises} : une sur le prix (\eng{How much...?}), une sur le délai (\eng{When...?} / \eng{How long...?}).
    \item \textbf{Rôle B (Employé) :} Préparez une réponse type : accueil, proposition de solution, annonce d'un \textbf{chiffre précis} (prix en FCFA, date, numéro de transaction, durée en heures/jours), formule de clôture.
\end{enumerate}

\textbf{Étape 2 : Jeu de rôle en binôme — Manche 1 (10 min)}
\begin{enumerate}
    \item Installez-vous face à face (dos à la classe si possible pour simuler un guichet / téléphone).
    \item Jouez la scène \textbf{sans lire vos notes}, de manière naturelle. L'enseignant chronomètre (2 min max par binôme).
    \item Les autres élèves (auditeurs) remplissent la \cle{Fiche Client} pendant l'échange.
\end{enumerate}

\textbf{Étape 3 : Échange de rôles et nouvelle carte — Manche 2 (10 min)}
\begin{enumerate}
    \item Changez de carte-situation (tirez-en une nouvelle).
    \item Inversez les rôles : A devient employé, B devient client.
    \item Rejouez la scène. Les auditeurs complètent une nouvelle fiche.
\end{enumerate}

\textbf{Étape 4 : Restitution et correction collective (15 min)}
\begin{enumerate}
    \item L'enseignant désigne 2 ou 3 binômes pour rejouer leur scène devant la classe.
    \item Les auditeurs donnent les réponses de leur fiche (prix, délai, numéro). L'enseignant note au tableau les bonnes formulations entendues et corrige les erreurs majeures (ordre des mots, modaux, vocabulaire).
    \item Auto-évaluation : chaque élève coche les items atteints dans le \cle{Bilan} final.
\end{enumerate}
\end{methode}


\begin{textebox}[Cartes-situations — À découper / Afficher]
\textbf{Carte 1 — Agence MTN Bamenda (Quarter Mile 3)} \\
\eng{Client:} You want to migrate your number to 4G and subscribe to the \eng{MTN Pulse} plan (5 GB/month). Ask: migration cost? activation date? \\
\eng{Agent:} Migration is free. Pulse 5 GB = 2,500 FCFA/month. Activation instant. Ref: MTN-BDA-8842.

\textbf{Carte 2 — Hôtel La Plage Kribi (Bord de mer)} \\
\eng{Client:} Book a double room, sea view, July 14–16 (2 nights). Ask: total price (breakfast included)? cancellation policy? \\
\eng{Receptionist:} 22,500 FCFA/night (B\&B). Total 45,000 FCFA. 50\% deposit. Free cancellation 48h before. Ref: HOT-KRI-559.

\textbf{Carte 3 — Agence UBA Yaoundé (Carrefour Warda)} \\
\eng{Client:} Urgent SWIFT transfer 250,000 FCFA to France (tuition). Ask: transfer fees? processing time? exchange rate? \\
\eng{Teller:} Fee 12,500 FCFA + correspondent charges. Time: 2–3 working days. Rate: 1 EUR = 655.957 FCFA. Ref: SWF-YDE-1102.

\textbf{Carte 4 — Garage Auto Rapide Douala (Akwa)} \\
\eng{Client:} Toyota Corolla 2015, grinding noise when braking. Ask: quote for front brake pads + discs? repair duration? \\
\eng{Mechanic:} Parts 35,000 FCFA. Labour 15,000 FCFA. Total 50,000 FCFA. Duration 3 hours (ready 2 PM). 6-month warranty.
\end{textebox}

\begin{textebox}[Fiche Client — Grille d'écoute pour les auditeurs (à recopier)]
\begin{center}
\begin{tabularx}{\linewidth}{|X|c|X|X|}\hline
\rowcolor{popblueL} \textbf{Scénario (Lieu / Ville)} & \textbf{Rôle A (Client)} & \textbf{Service demandé} & \textbf{Infos chiffrées (Prix / Délai / N°)} \\ \hline
1. MTN Bamenda & & & \\ \hline
2. Hôtel Kribi & & & \\ \hline
3. UBA Yaoundé & & & \\ \hline
4. Garage Douala & & & \\ \hline
\end{tabularx}
\end{center}
\emph{Consigne :} Remplissez en anglais ou en français. Notez le chiffre exact entendu (ex. \eng{25,000 FCFA}, \eng{48 hours}, \eng{Ref: TXN7789}).
\end{textebox}

\begin{exoresolu}[Modèles de dialogues attendus — Corrigé commenté pour l'enseignant]
\textbf{Scénario 1 : MTN Bamenda}
\begin{textebox}[Dialogue modèle]
\eng{Agent:} Good morning. Welcome to MTN. How can I help you?
\eng{Client:} Good morning. \textbf{I would like to migrate my number to 4G} and \textbf{subscribe to} the MTN Pulse plan. \textbf{Could you please tell me how much the migration costs?}
\eng{Agent:} Certainly. The migration is \textbf{free of charge}. The Pulse plan with 5 GB costs \textbf{2,500 FCFA} per month. Activation is \textbf{instant}. \textbf{Your reference number is MTN-BDA-8842.}
\eng{Client:} \textbf{Thank you very much.} Have a nice day.
\eng{Agent:} You're welcome. Enjoy your 4G connection!
\end{textebox}
\textbf{Commentaire :} \eng{migrate to} / \eng{subscribe to} (prépositions). \eng{Free of charge} (gratuit). \eng{Instant} (immédiat). Structure cible : \eng{I would like to} + BV ; \eng{Could you tell me} + ordre affirmatif.

\textbf{Scénario 2 : Hôtel Kribi}
\begin{textebox}[Dialogue modèle]
\eng{Receptionist:} Good afternoon. Hotel La Plage. How may I assist you?
\eng{Client:} Good afternoon. \textbf{I'd like to book a double room with sea view} for July 14th to 16th. \textbf{Could you tell me the total price including breakfast?} Also, \textbf{what is your cancellation policy?}
\eng{Receptionist:} Of course. It's \textbf{22,500 FCFA per night}, breakfast included. Total \textbf{45,000 FCFA} for two nights. \textbf{50\% deposit} required. Free cancellation \textbf{up to 48 hours before check-in}. Your booking reference: \textbf{HOT-KRI-559}.
\eng{Client:} Perfect. I'll send the deposit. Thank you.
\eng{Receptionist:} You're welcome. We look forward to welcoming you.
\end{textebox}
\textbf{Commentaire :} \eng{Including breakfast} (pdj inclus). \eng{Up to 48 hours before} (jusqu'à 48h avant). \eng{Look forward to + -ing} (avoir hâte de).

\textbf{Scénario 3 : UBA Yaoundé}
\begin{textebox}[Dialogue modèle]
\eng{Teller:} Good morning. Welcome to UBA. What can I do for you?
\eng{Client:} Good morning. \textbf{I need to make a SWIFT transfer of 250,000 FCFA to France.} \textbf{Could you tell me the transfer fees and how long it will take?}
\eng{Teller:} Right away. The fee is \textbf{12,500 FCFA}. Correspondent bank charges may apply. Processing time: \textbf{2 to 3 working days}. Exchange rate today: \textbf{655.957 FCFA} for 1 Euro. Reference: \textbf{SWF-YDE-1102}.
\eng{Client:} Thank you. Here is the beneficiary's IBAN.
\eng{Teller:} Thank you. Here is your receipt.
\end{textebox}
\textbf{Commentaire :} \eng{Working days} (jours ouvrés). \eng{May apply} (peuvent s'appliquer - probabilité). \eng{Beneficiary} (bénéficiaire).

\textbf{Scénario 4 : Garage Douala}
\begin{textebox}[Dialogue modèle]
\eng{Mechanic:} Good morning. Auto Rapide. How can I help?
\eng{Client:} Good morning. My Corolla \textbf{makes a grinding noise when braking}. \textbf{I'd like a quote for front brake pads and discs.} \textbf{How long will the repair take?}
\eng{Mechanic:} Parts: \textbf{35,000 FCFA}. Labour: \textbf{15,000 FCFA}. Total estimate: \textbf{50,000 FCFA}. It will take \textbf{3 hours}. Ready by \textbf{2 PM}. \textbf{6-month warranty} included.
\eng{Client:} That works. I'll bring it right away. Thanks.
\eng{Mechanic:} You're welcome. See you soon.
\end{textebox}
\textbf{Commentaire :} \eng{Grinding noise} (bruit de grincement/ferraille). \eng{Quote / Estimate} (devis). \eng{Labour} (main-d'œuvre). \eng{Ready by + heure} (prêt pour/à).
\end{exoresolu}

\courssub{Bilan et auto-évaluation}

\begin{aretenir}[Auto-évaluation — Cochez votre niveau]
\begin{center}
\begin{tabularx}{\linewidth}{|X|c|c|c|}\hline
\rowcolor{popblueL} \textbf{Compétence / Savoir-faire} & \textbf{\eng{Not yet}} & \textbf{\eng{With help}} & \textbf{\eng{On my own}} \\ \hline
Je formule une demande de service avec \eng{I would like to} / \eng{Could you}. & & & \\ \hline
J'utilise au moins deux formules de politesse différentes. & & & \\ \hline
Je pose une question sur le prix (\eng{How much...}) et une sur le délai (\eng{When/How long}). & & & \\ \hline
Je comprends et note un prix, une date, un numéro de référence entendus. & & & \\ \hline
Je joue le rôle de l'employé : j'annonce un chiffre précis et je réponds poliment. & & & \\ \hline
Je corrige mes erreurs d'ordre des mots (questions indirectes) et de modaux. & & & \\ \hline
\end{tabularx}
\end{center}
\end{aretenir}

\courssub{Prolongement : Écrit et Recherche}

\begin{methode}[Devoir / Projet — À rendre la semaine prochaine]
\begin{enumerate}
    \item \textbf{Expression écrite (120–150 mots) :} Rédigez un \textbf{e-mail formel} à la direction de l'un des quatre services (MTN, Hôtel, UBA, Garage) pour \textbf{confirmer votre demande} ou faire une \textbf{réclamation} (ex. : délai non respecté, erreur de facturation, service non conforme).
    \begin{itemize}
        \item Structure imposée :
        \eng{Subject: Confirmation / Complaint — [Service Name] — Ref: [Number]}
        \eng{Dear Sir / Madam,}
        \eng{I am writing to... (confirm my request / complain about...)}
        \eng{On [date], I requested... / I was told that...}
        \eng{However, ... (problem) / I would appreciate it if you could... (action desired).}
        \eng{Yours faithfully,}
        \eng{[Your Name]}
    \end{itemize}
    \item \textbf{Recherche documentaire :} Comparez les tarifs réels actuels (sites web / appels) pour : MTN vs Orange (forfaits 4G/5G), Hôtels à Kribi (vue mer), Banques UBA vs SCB vs Afriland (frais SWIFT), Garages Douala/Yaoundé (plaquettes Corolla). Présentez un \textbf{tableau comparatif en anglais} (Service | Provider | Price | Delay | Pros/Cons) en classe prochaine.
\end{enumerate}
\end{methode}

\courssub{Le saviez-vous ? — Le secteur des services au Cameroun}

\begin{savaistu}[Contexte économique et linguistique]
Le secteur tertiaire (services) représente plus de \textbf{50 \% du PIB camerounais}. Les télécommunications (MTN, Orange, Camtel) et la finance mobile (\eng{Mobile Money} : \eng{MoMo}, \eng{Orange Money}) sont les moteurs de l'inclusion financière. En 2023, plus de 10 millions de Camerounais utilisaient le \eng{Mobile Money} pour payer factures, transferts et achats quotidiens.

Maîtriser l'anglais des services (\eng{requesting services, making complaints, negotiating prices}) est un atout majeur pour l'employabilité dans les villes bilingues (Yaoundé, Douala, Buea, Bamenda, Limbe) et pour travailler dans les entreprises internationales, ONGs, banques, hôtels et centres d'appels (\eng{call centers}) implantés au Cameroun.

\textbf{Termes camerounais courants en anglais :}
\begin{itemize}
    \item \eng{Bendskin} / \eng{Okada} (moto-taxi) — \eng{Commercial bike}.
    \item \eng{Call box} / \eng{Cabinet} (kiosque de transfert d'argent/téléphone).
    \item \eng{Buyam-sellam} (commerçant ambulant) — \eng{Street vendor / Hawker}.
    \item \eng{Ngoua} / \eng{Eru} / \eng{Koki} (plats locaux) — \eng{Local dishes}.
\end{itemize}
\end{savaistu}

\begin{attention}[Rappel orthographe britannique (UK) — Examen / Bac]
\begin{itemize}
    \item \eng{Programme} (pas \eng{program} — sauf informatique), \eng{Centre} (pas \eng{center}), \eng{Metre} (unité), \eng{Litre}, \eng{Fibre}.
    \item \eng{Cheque} (chèque bancaire), \eng{Analyse}, \eng{Catalogue}, \eng{Dialogue}.
    \item \eng{Organisation}, \eng{Realisation}, \eng{Specialise} (terminaison \eng{-ise} / \eng{-isation} préférée au \eng{-ize} / \eng{-ization} au Royaume-Uni, bien que \eng{-ize} soit accepté par Oxford).
    \item \eng{Travelling} (double \eng{l}), \eng{Cancelled}, \eng{Modelling}, \eng{Labelled}.
    \item \eng{Practise} (verbe) vs \eng{Practice} (nom). \eng{Advise} (verbe) vs \eng{Advice} (nom). \eng{Device} (nom) vs \eng{Devise} (verbe).
    \item Dates : \eng{14 July 2025} ou \eng{14th July 2025} (pas \eng{July 14th} seul en UK formel). Heures : \eng{2.00 pm} / \eng{14:00}.
\end{itemize}
\end{attention}

\newpage

\coursec{Méthodes et exercices résolus}

\courssub{Méthodes et exercices résolus}

\begin{methode}[Identifier le thème et le contexte d'un échange de service]
\textbf{Objectif :} Comprendre la situation globale (qui parle, où, dans quel but) dès la première écoute.
\begin{enumerate}
    \item \textbf{Analyser le titre et les consignes :} Ils donnent le cadre (ex. \eng{At the service desk}, \eng{At the post office}).
    \item \textbf{Écouter les premiers échanges :} Les formules de politesse (\eng{Good morning}, \eng{How can I help you?}) et le ton (formel/informel) situent la relation (client/employé, usager/administration).
    \item \textbf{Repérer le \cle{champ lexical} dominant :} Mots-clés récurrents : \eng{parcel}, \eng{form}, \eng{queue}, \eng{receipt}, \eng{tracking number}, \eng{complaint}, \eng{refund}.
    \item \textbf{Poser l'hypothèse :} « Il s'agit d'un client qui récupère un colis / dépose une réclamation / demande un renseignement. »
    \item \textbf{Valider à la deuxième écoute :} Confirmer ou affiner l'hypothèse grâce aux détails.
\end{enumerate}
\end{methode}

\begin{textebox}[Script d'écoute : At the Post Office]
\eng{-- Good morning, Madam. How can I help you today?}
\eng{-- Good morning. I would like to send this parcel to Bamenda, please.}
\eng{-- Certainly. Let me weigh it... It's 2.5 kilos. Do you want it sent by Express or Standard mail?}
\eng{-- How long does Standard take?}
\eng{-- Usually three to four working days. Express is next day before noon.}
\eng{-- I'll take Express, please. It's urgent.}
\eng{-- Right. That will be 4,500 CFA francs. Do you have the recipient's address and phone number?}
\eng{-- Yes, here is the form filled out.}
\eng{-- Thank you. Here is your receipt and the tracking number. You can track it online.}
\eng{-- Thank you very much. Goodbye.}
\eng{-- Goodbye, Madam.}
\end{textebox}

\begin{exoresolu}[Application : Contexte global]
\textbf{Écoutez le dialogue ci-dessus (lu par le professeur ou audio) et répondez aux questions.}
\begin{enumerate}
    \item Where does the scene take place?
    \item Who are the speakers?
    \item What is the customer's main request?
    \item Which option does the customer choose and why?
\end{enumerate}
\tcblower
\textbf{Solution détaillée :}
\begin{enumerate}
    \item \textbf{Where does the scene take place?} \eng{At the post office.} (Indices : \eng{“send this parcel”}, \eng{“Express or Standard mail”}, \eng{“tracking number”}, \eng{“receipt”}).
    \item \textbf{Who are the speakers?} A \textbf{clerk / postal worker} (employee) and a \textbf{customer} (client). Le ton est formel (\eng{“Madam”}, \eng{“How can I help you?”}).
    \item \textbf{What is the customer's main request?} \eng{She would like to send a parcel to Bamenda.} (Structure : \eng{I would like to} + infinitif pour exprimer un souhait poli).
    \item \textbf{Which option does the customer choose and why?} \eng{She chooses Express mail because it is urgent.} (Justification : \eng{“It's urgent”} / \eng{“next day before noon”}).
\end{enumerate}
\textbf{Conclusion :} La première écoute suffit à situer la scène (bureau de poste), les acteurs (employé/client) et l'action principale (envoi colis). La deuxième écoute précise le choix (Express) et la raison (urgence).
\end{exoresolu}

\begin{methode}[Repérer les informations chiffrées, les noms propres et les mots-clés]
\textbf{Objectif :} Isoler les données factuelles précises (prix, poids, délais, noms, numéros) souvent ciblées par les items Vrai/Faux ou QCM.
\begin{enumerate}
    \item \textbf{Pré-lecture des questions :} Avant la 2ème écoute, lire les items. Souligner les \cle{mots-questions} : \eng{How much?}, \eng{How long?}, \eng{What time?}, \eng{What name?}, \eng{Which number?}.
    \item \textbf{Écoute active ciblée :} Ne pas tout réécouter de la même façon. Tendre l'oreille sur les \cle{chiffres} (souvent prononcés distinctement), les \cle{noms propres} (orthographiés parfois : \eng{B-A-M-E-N-D-A}) et les \cle{mots de liaison} logiques (\eng{because}, \eng{but}, \eng{however}, \eng{so}).
    \item \textbf{Prise de notes abrégée :} Noter en chiffres (\eng{2.5 kg}, \eng{4500 CFA francs}, \eng{24h}) et abréviations (\eng{Expr.} pour Express, \eng{Std} pour Standard). Ne pas écrire des phrases complètes.
    \item \textbf{Vérification croisée :} Confronter ses notes aux items. Attention aux \textbf{pièges} : l'employé donne souvent deux options (prix Express ET prix Standard), la question ne porte que sur celle choisie.
\end{enumerate}
\end{methode}

\begin{exoresolu}[Application : Détails factuels (Vrai / Faux / Non mentionné)]
\textbf{D'après le dialogue précédent, dites si les affirmations sont True (T), False (F) ou Not Mentioned (NM). Justifiez par une citation du texte.}
\begin{enumerate}
    \item The parcel weighs 3 kilos.
    \item Standard mail takes three to four working days.
    \item The customer pays 5,000 CFA francs.
    \item The recipient's name is Mr. Nfor.
    \item The tracking number allows online tracking.
\end{enumerate}
\tcblower
\textbf{Solution détaillée :}
\begin{enumerate}
    \item \textbf{F (False).} Texte : \eng{“It's 2.5 kilos.”} $\neq$ 3 kilos.
    \item \textbf{T (True).} Texte : \eng{“Usually three to four working days.”} Correspondance exacte.
    \item \textbf{F (False).} Texte : \eng{“That will be 4,500 CFA francs.”} $\neq$ 5,000. \textbf{Piège :} Ne pas confondre avec un arrondi ou une autre devise.
    \item \textbf{NM (Not Mentioned).} Le texte donne l'adresse et le numéro de téléphone (\eng{“recipient's address and phone number”}), mais \textbf{jamais le nom} du destinataire.
    \item \textbf{T (True).} Texte : \eng{“You can track it online.”} (Référence au \eng{tracking number}).
\end{enumerate}
\textbf{Astuce Bac :} Pour \textbf{NM}, ne pas déduire. Si l'info n'est pas prononcée $\rightarrow$ NM. Ici, on ne sait pas si le destinataire s'appelle Nfor ou pas.
\end{exoresolu}

\begin{methode}[Maîtriser le lexique et les formules de politesse pour demander un service]
\textbf{Objectif :} Utiliser les structures standards pour formuler une demande, accepter, refuser ou préciser un besoin dans un contexte formel (administration, commerce, services publics).
\begin{enumerate}
    \item \textbf{Classer le lexique par fonction :}
    \begin{itemize}
        \item \textbf{Demander :} \eng{I would like to...}, \eng{Could you please...?}, \eng{I need to...}, \eng{May I...?}
        \item \textbf{Préciser :} \eng{Urgent}, \eng{Standard}, \eng{Express}, \eng{Fragile}, \eng{Registered}, \eng{Tracking number}, \eng{Receipt}, \eng{Proof of identity / ID}.
        \item \textbf{Répondre (Employé) :} \eng{Certainly.}, \eng{Of course.}, \eng{Would you like...?}, \eng{Do you have...?}, \eng{Here is your...}, \eng{The fee is...}
    \end{itemize}
    \item \textbf{Respecter la structure \cle{Modale + Sujet + Base Verbale} :} \eng{Could you \textbf{send}...?} (pas \eng{to send}), \eng{Would you \textbf{like}...?}.
    \item \textbf{Distinguier \eng{I would like} (souhait poli) et \eng{I like} (goût général) :} \eng{I would like to send} = \eng{Je voudrais envoyer}. \eng{I like sending} = \eng{J'aime envoyer}.
    \item \textbf{S'entraîner à la substitution :} Remplacer l'objet (\eng{parcel} $\rightarrow$ \eng{letter}, \eng{money}, \eng{form}) et la destination (\eng{Bamenda} $\rightarrow$ \eng{Yaoundé}, \eng{London}).
\end{enumerate}
\end{methode}

\begin{exoresolu}[Application : Vocabulaire et transformation]
\begin{enumerate}
    \item \textbf{Complétez avec le mot approprié :} \eng{receipt}, \eng{tracking number}, \eng{counter}, \eng{parcel}, \eng{form}.
    \begin{enumerate}
        \item Please fill out this \_\_\_\_\_\_ with your address.
        \item I want to send this \_\_\_\_\_\_ to Douala.
        \item Keep the \_\_\_\_\_\_ as proof of payment.
        \item You can check the delivery status with the \_\_\_\_\_\_.
        \item Go to \_\_\_\_\_\_ number 3 for registered mail.
    \end{enumerate}
    \item \textbf{Transformez ces phrases directes en demandes polies (formelles) :}
    \begin{enumerate}
        \item \eng{Send this letter to Buea.}
        \item \eng{Give me a form.}
        \item \eng{I want to pay by card.}
    \end{enumerate}
\end{enumerate}
\tcblower
\textbf{Solution détaillée :}
\textbf{1. Vocabulaire contextuel :}
\begin{enumerate}
    \item \eng{form} (collocation : \eng{fill out a form}).
    \item \eng{parcel} (colis, objet physique à envoyer).
    \item \eng{receipt} (reçu, preuve de paiement).
    \item \eng{tracking number} (numéro de suivi, pour \eng{check status}).
    \item \eng{counter} (guichet, lieu physique numéroté).
\end{enumerate}
\textbf{2. Transformation en demandes polies (Register formel) :}
\begin{enumerate}
    \item \eng{Could you please send this letter to Buea?} (ou \eng{Would you mind sending...?})
    \item \eng{Could I have a form, please?} / \eng{May I have a form?} (Demander \textbf{pour soi} $\rightarrow$ \eng{Can/Could/May I...?})
    \item \eng{I would like to pay by card, please.} / \eng{Could I pay by card?} (\eng{I want} $\rightarrow$ trop direct/impoli $\rightarrow$ \eng{I would like}).
\end{enumerate}
\textbf{Rappel grammatical :} \eng{Could you + BV} (demande à autrui) vs \eng{Could I + BV} (demande pour soi / permission).
\end{exoresolu}

\begin{methode}[Structurer une réponse courte et une production orale guidée (Role-play)]
\textbf{Objectif :} Réinvestir le lexique et les structures pour simuler un échange au guichet (Speaking) ou rédiger un dialogue court (Writing).
\begin{enumerate}
    \item \textbf{Identifier les \cle{étapes obligatoires} de l'échange type :}
    \begin{enumerate}
        \item Salutation polie (\eng{Good morning/afternoon}).
        \item Offre d'aide (\eng{How can I help?}).
        \item Formulation de la demande (\eng{I'd like to... / Could you...?}).
        \item Questions de précision (Poids, destination, mode, paiement).
        \item Transaction (Prix, paiement, remise reçu).
        \item Clôture polie (\eng{Thank you / Goodbye}).
    \end{enumerate}
    \item \textbf{Utiliser les \cle{connecteurs logiques} :} \eng{First}, \eng{Then}, \eng{Also}, \eng{Finally} (pour structurer l'oral).
    \item \textbf{Gérer l'imprévu (stratégie) :} Si on ne comprend pas $\rightarrow$ \eng{“Sorry, could you repeat please?”}, \eng{“Could you speak more slowly?”}, \eng{“How do you spell that?”}.
    \item \textbf{S'entraîner en binôme :} A = Employé (a la fiche prix/délais), B = Client (a une carte de situation : \eng{Send 5kg parcel to London, Express, pay cash}). Échanger les rôles.
\end{enumerate}
\end{methode}

\begin{textebox}[Dialogue : Opening a Savings Account]
\eng{\textbf{Clerk:} Good morning, Sir. How can I help you today?}
\eng{\textbf{Client:} Good morning. I would like to open a savings account, please.}
\eng{\textbf{Clerk:} Certainly. Do you have your ID card and a passport photo?}
\eng{\textbf{Client:} Yes, here is my ID card and here are two photos.}
\eng{\textbf{Clerk:} Thank you. We also need an initial deposit of 10,000 CFA francs.}
\eng{\textbf{Client:} No problem. Here is the money.}
\eng{\textbf{Clerk:} Perfect. You will receive a free debit card and access to our mobile banking app.}
\eng{\textbf{Client:} That sounds great. How do I activate the app?}
\eng{\textbf{Clerk:} You will receive an SMS with a link. Just follow the instructions.}
\eng{\textbf{Client:} Thank you very much for your help.}
\eng{\textbf{Clerk:} You're welcome, Sir. Have a nice day.}
\eng{\textbf{Client:} Goodbye.}
\end{textebox}

\begin{exoresolu}[Production guidée : Écriture d'un dialogue]
\textbf{Consigne :} Rédigez un dialogue complet (10-12 répliques) entre un employé de banque et un client. Le client veut ouvrir un compte d'épargne (\eng{savings account}). Il doit fournir une pièce d'identité (\eng{ID card}), une photo d'identité (\eng{passport photo}) et un dépôt initial de 10 000 CFA francs. L'employé explique les conditions (carte bancaire gratuite, application mobile).
\textbf{Contraintes :} Utiliser \eng{I would like to}, \eng{Could you please}, \eng{Here is}, \eng{You will receive}.
\tcblower
\textbf{Proposition de solution (Modèle corrigé) :}
\eng{(Voir le dialogue ci-dessus)}
\textbf{Analyse des structures utilisées :}
\begin{itemize}
    \item \eng{I would like to open} (Demande polie).
    \item \eng{Do you have...?} / \eng{Here is/are...} (Échange d'objets).
    \item \eng{You will receive} (Futur simple pour conséquence certaine / information future).
    \item \eng{Just follow the instructions} (Impératif pour mode d'emploi).
\end{itemize}
\end{exoresolu}

\begin{methode}[Répondre aux questions de compréhension détaillée (Inférence et Reformulation)]
\textbf{Objectif :} Aller au-delà du littéral : comprendre une implication, une cause, une conséquence ou reformuler une idée avec ses propres mots (compétence Bac : \eng{Answer in your own words}).
\begin{enumerate}
    \item \textbf{Identifier le type de question :}
    \begin{itemize}
        \item \eng{Why...?} $\rightarrow$ Cause/Raison (\eng{Because}, \eng{Since}, \eng{As}).
        \item \eng{How...?} $\rightarrow$ Manière/Moyen/Durée/Prix.
        \item \eng{What does X mean?} $\rightarrow$ Définition contextuelle (synonyme).
        \item \eng{According to the text...} $\rightarrow$ Preuve textuelle obligatoire.
    \end{itemize}
    \item \textbf{Technique de la \cle{reformulation} (Paraphrase) :}
    \begin{itemize}
        \item Changer la structure : \eng{Active $\leftrightarrow$ Passive}, \eng{Direct $\leftrightarrow$ Indirect}.
        \item Remplacer le vocabulaire par des synonymes : \eng{parcel $\rightarrow$ package}, \eng{urgent $\rightarrow$ pressing}, \eng{fee $\rightarrow$ cost/price}, \eng{clerk $\rightarrow$ assistant}.
        \item Garder le sens exact, ne pas ajouter d'opinion personnelle.
    \end{itemize}
    \item \textbf{Citer la preuve :} Pour les questions \eng{Justify with a quote}, copier \textbf{exactement} le segment de phrase (guillemets anglais “ ”), sans le déformer.
\end{enumerate}
\end{methode}

\begin{textebox}[Listening Script : At the Internet Café]
\eng{Manager: Welcome to Connect World. Do you have a membership card?}
\eng{Student: No, I don't. I just need to print three pages urgently.}
\eng{Manager: Okay. Printing is 100 francs per page. But the printer has a paper jam right now. The technician is coming in twenty minutes.}
\eng{Student: Twenty minutes? I have an exam in an hour. Can I use a computer to type my CV while waiting?}
\eng{Manager: Sure. Internet access is 200 francs per hour. You can use booth number 4.}
\eng{Student: Fine. I'll pay for the hour and the printing later.}
\end{textebox}

\begin{exoresolu}[Compréhension détaillée et Reformulation]
\textbf{Texte support (lu deux fois) : voir le script ci-dessus.}
\textbf{Questions :}
\begin{enumerate}
    \item Why can't the student print immediately? (Answer in a full sentence).
    \item How much will the student pay in total if he uses the computer for one hour and prints 3 pages? (Calculate).
    \item Find in the text a synonym for "booth" (line 6) in this context.
    \item Reformulate: "The printer has a paper jam." (Start with: Printing is impossible...).
\end{enumerate}
\tcblower
\textbf{Solution détaillée :}
\begin{enumerate}
    \item \textbf{Why?} \eng{The student cannot print immediately because the printer has a paper jam (and the technician is only coming in twenty minutes).}
    \item \textbf{Calcul :} Internet 1h = 200 FCFA. Impression 3 pages = 3 $\times$ 100 = 300 FCFA. \textbf{Total = 500 FCFA.}
    \item \textbf{Synonyme de "booth" :} \eng{Cubicle} / \eng{Workstation} / \eng{Terminal} / \eng{Station}. (Contexte : place individuelle avec ordinateur).
    \item \textbf{Reformulation :} \eng{Printing is impossible at the moment because the printer has a paper jam.} (Passif / Structure causale). Autre possibilité : \eng{Printing cannot be done right now due to a paper jam.}
\end{enumerate}
\textbf{Point méthode :} Pour le calcul, noter les chiffres \textbf{au fur et à mesure} (200, 100, 3). Pour la reformulation, ne pas changer le temps (présent $\rightarrow$ présent) ni le sens.
\end{exoresolu}

\begin{methode}[Synthèse : La fiche "Survie" pour l'épreuve de Listening (Bac)]
\textbf{Objectif :} Avoir une check-list mentale le jour de l'examen.
\begin{enumerate}
    \item \textbf{Avant l'écoute (1 min) :} Lire le titre, les questions, souligner les \cle{mots-clés} (Who, Where, How much, Why, True/False).
    \item \textbf{1ère écoute (Global) :} Noter : Lieu, Acteurs, Sujet principal, Ton. \textbf{Ne pas écrire de phrases}, juste des mots : \eng{Post office / Clerk-Client / Send parcel / Express / 4500}.
    \item \textbf{Pause (30s) :} Relire les questions. Anticiper les réponses. Préparer le stylo pour les chiffres/noms.
    \item \textbf{2ème écoute (Détails) :} Cibler \textbf{uniquement} les infos manquantes. Chiffres $\rightarrow$ écrire en chiffres. Noms propres $\rightarrow$ phonétique approximative acceptée si pas d'orthographe demandée.
    \item \textbf{Après l'écoute (Relecture) :} Vérifier la cohérence (le prix correspond-il au choix Express ?). Remplir les items Vrai/Faux/Justifier. \textbf{Ne laisser aucune case vide.}
\end{enumerate}
\end{methode}

\begin{textebox}[Script : At the Travel Agency]
\eng{Agent: Good afternoon. Buea Travels. How may I assist you?}
\eng{Customer: Hello. I need a return ticket to Douala for tomorrow morning.}
\eng{Agent: We have a bus leaving at 6:30 AM and another at 8:00 AM. The 6:30 is a VIP coach, 5,000 francs. The 8:00 is standard, 3,500 francs.}
\eng{Customer: I'll take the VIP one. I need to be in Douala before 10 AM for a meeting.}
\eng{Agent: The VIP arrives at 9:45 AM. Do you want to book now or pay at the agency?}
\eng{Customer: I'll book now. Can I pay by Mobile Money?}
\eng{Agent: Yes, send to 6XX XXX XXX. Reference: VIP-DLA-001.}
\eng{Customer: Done. I've sent 5,000 francs.}
\eng{Agent: Received. Your e-ticket is sent to your WhatsApp. Safe journey.}
\eng{Customer: Thank you. Bye.}
\end{textebox}

\begin{exoresolu}[Entraînement final : Mini-Bac Listening]
\textbf{Script (Lu 2 fois par l'enseignant) : voir le script ci-dessus.}
\textbf{Questions d'examen (Type Bac) :}
\begin{enumerate}
    \item \textbf{MCQ :} The customer chooses the 6:30 AM bus because: A) It is cheaper. B) He has a meeting before 10 AM. C) It is a standard coach.
    \item \textbf{Short Answer :} What is the arrival time of the VIP bus in Douala?
    \item \textbf{True/False + Justify :} The customer pays cash at the agency. (T/F + Quote).
    \item \textbf{Vocabulary :} Find a word in the text meaning "reserve a seat".
    \item \textbf{Grammar in context :} Rewrite: "I need a ticket" $\rightarrow$ Polite request with "Would like".
\end{enumerate}
\tcblower
\textbf{Corrigé type Bac :}
\begin{enumerate}
    \item \textbf{B.} Justification texte : \eng{“I need to be in Douala before 10 AM for a meeting.”} (A est faux : 5000 > 3500. C est faux : c'est VIP).
    \item \textbf{9:45 AM.} (Info chiffrée précise, notée à la 2ème écoute).
    \item \textbf{False.} Justification : \eng{“I'll book now. Can I pay by Mobile Money?”} / \eng{“I've sent 5,000 francs.”} (Paiement à distance, pas à l'agence).
    \item \textbf{Book} (verbe) ou \textbf{Booking} (nom). Texte : \eng{“Do you want to book now?”}
    \item \eng{I would like a ticket, please.} / \eng{I would like to book a ticket.}
\end{enumerate}
\textbf{Bilan :} Maîtriser les chiffres (heures, prix), les connecteurs logiques (\eng{because}), les synonymes (\eng{book/reserve}) et la politesse (\eng{would like}).
\end{exoresolu}

\begin{attention}[Les 5 erreurs qui coûtent des points au Bac]
\begin{enumerate}
    \item \textbf{Confusion \eng{I want} / \eng{I would like} :} \eng{“I want a ticket”} = Impoli / Familier. \textbf{Attendu :} \eng{I would like a ticket} ou \eng{Could I have a ticket?} (Register formel obligatoire en situation de service).
    \item \textbf{Mauvaise formation de la question polie :} \eng{*“Could you to send?”} (Intrus \eng{to}). \textbf{Règle :} \eng{Modal + Sujet + Base Verbale (sans to)}. \eng{Could you send?}, \eng{Would you mind sending?} (Gérondif après \eng{mind}).
    \item \textbf{Oubli du \eng{please} / \eng{thank you} :} En anglais, l'absence de marqueur de politesse dans une demande change le sens en ordre/impolitesse. Toujours ajouter \eng{please} en fin de phrase ou \eng{could you please}.
    \item \textbf{Erreurs de chiffres à l'oral/écrit :} Confusion \eng{15 / 50} (fifteen/fifty), \eng{13 / 30} (thirteen/thirty). \textbf{Parade :} Écrire en chiffres arabes (15, 50) lors de la prise de notes. Pour l'oral : bien articuler le « n » de \eng{thirteen} (son « ène ») vs le « t » de \eng{thirty} (son « ète »).
    \item \textbf{Prépositions incorrectes avec le vocabulaire du service :}
    \begin{itemize}
        \item \eng{*“I go to the post office \textbf{for} send a parcel”} $\rightarrow$ \eng{\textbf{to send}} (but finalité).
        \item \eng{*“I pay \textbf{by} cash”} $\rightarrow$ \eng{\textbf{in} cash} (mais \eng{by card}, \eng{by Mobile Money}, \eng{by cheque}).
        \item \eng{*“The parcel is \textbf{to} Bamenda”} $\rightarrow$ \eng{\textbf{for} Bamenda} (destination colis) ou \eng{\textbf{to} Bamenda} (verbe mouvement \eng{send to}). \eng{Send \textbf{to} [person/place]}, \eng{A parcel \textbf{for} [destination]}.
    \end{itemize}
\end{enumerate}
\end{attention}

\newpage

\coursec{Exercices}

\courssub{Je vérifie mes connaissances}

\begin{exercice}[QCM : Compréhension globale d'un dialogue]
\textbf{Écoutez le dialogue suivant (ou lisez-le) et choisissez la bonne réponse.}
\begin{textebox}[Script : At the Post Office in Buea]
Clerk: Good morning, sir. How can I help you today?
Customer: Good morning. I'd like to send this parcel to Yaoundé, please.
Clerk: Certainly. Let me weigh it. It's 2.5 kg. That will be 3,500 CFA francs for express delivery.
Customer: Express? How long does it take?
Clerk: It arrives tomorrow before noon. Standard delivery takes three days and costs 2,000 francs.
Customer: I'll take the express option. Here is the money.
Clerk: Thank you. Here is your receipt and the tracking number. Have a nice day.
Customer: Thank you. Goodbye.
\end{textebox}
\begin{enumerate}
\item Where does the conversation take place? 
\begin{itemize}
\item[A.] At a bank. \item[B.] At a post office. \item[C.] At a travel agency. \item[D.] At a supermarket.
\end{itemize}
\item What does the customer want to do? 
\begin{itemize}
\item[A.] Buy stamps. \item[B.] Send a parcel. \item[C.] Receive a parcel. \item[D.] Open a bank account.
\end{itemize}
\item How much does the express delivery cost? 
\begin{itemize}
\item[A.] 2,000 CFA francs. \item[B.] 3,000 CFA francs. \item[C.] 3,500 CFA francs. \item[D.] 5,000 CFA francs.
\end{itemize}
\item When will the parcel arrive if sent by express? 
\begin{itemize}
\item[A.] Today. \item[B.] Tomorrow before noon. \item[C.] In three days. \item[D.] Next week.
\end{itemize}
\item What does the clerk give the customer at the end? 
\begin{itemize}
\item[A.] A stamp and a pen. \item[B.] A box and tape. \item[C.] A receipt and a tracking number. \item[D.] A letter and an envelope.
\end{itemize}
\end{enumerate}
\end{exercice}

\begin{exercice}[Vrai / Faux justifié : Repérage de détails]
\textbf{Lisez le script ci-dessous. Pour chaque affirmation, cochez V (Vrai) ou F (Faux) et justifiez en citant un passage du texte.}
\begin{textebox}[Script : At the Bank Counter]
Customer: Excuse me, I would like to open a savings account.
Bank Officer: Of course. Do you have your national identity card and a proof of address?
Customer: Yes, here is my ID card. For the proof of address, I have an electricity bill from last month.
Bank Officer: That's perfect. The minimum opening deposit is 10,000 CFA francs. There are no monthly fees for students under 25.
Customer: I am 22 and I study at the University of Douala. Here is 15,000 francs.
Bank Officer: Thank you. Please sign this form. Your account is now active. You will receive your ATM card by post within five working days.
Customer: Thank you very much. Goodbye.
\end{textebox}
\begin{enumerate}
\item The customer wants to open a current account (compte courant).
\item The customer provides a water bill as proof of address.
\item The minimum deposit required is 10,000 CFA francs.
\item The customer is a student at the University of Yaoundé.
\item The ATM card will be delivered immediately at the counter.
\end{enumerate}
\end{exercice}

\begin{exercice}[Vocabulaire : Les formules de demande de service]
\textbf{Complétez les phrases suivantes avec le mot ou l'expression approprié(e) parmi la liste : \eng{receipt}, \eng{tracking number}, \eng{deposit}, \eng{proof of address}, \eng{counter}, \eng{parcel}, \eng{express delivery}, \eng{form}.}
\begin{enumerate}
\item I need to send this \_\_\_\_\_\_\_\_ to my brother in Bamenda.
\item Please go to \_\_\_\_\_\_\_\_ number 3 to pay your bill.
\item The \_\_\_\_\_\_\_\_ allows you to follow your package online.
\item Keep the \_\_\_\_\_\_\_\_ as proof of payment.
\item For \_\_\_\_\_\_\_\_, the package arrives in 24 hours.
\item You must fill in this \_\_\_\_\_\_\_\_ with your personal details.
\item A recent electricity bill serves as \_\_\_\_\_\_\_\_.
\item The initial \_\_\_\_\_\_\_\_ for this account is 5,000 francs.
\end{enumerate}
\end{exercice}

\begin{exercice}[Grammaire : Modaux et formes de politesse]
\textbf{Mettez les verbes entre parenthèses à la forme correcte (modaux de politesse : \eng{could}, \eng{would}, \eng{can}, \eng{may}, \eng{shall}) pour formuler des demandes de service.}
\begin{enumerate}
\item \_\_\_\_\_\_\_\_ you please tell me where the nearest post office is?
\item I \_\_\_\_\_\_\_\_ like to book a table for two for tonight.
\item \_\_\_\_\_\_\_\_ I help you with your luggage, madam?
\item \_\_\_\_\_\_\_\_ you check if this book is available in the library?
\item We \_\_\_\_\_\_\_\_ be grateful if you could send the confirmation by email.
\item \_\_\_\_\_\_\_\_ you mind opening the window? It is very hot in here.
\end{enumerate}
\end{exercice}

\courssub{Je m'entraîne}

\begin{exercice}[Traduction : Du français vers l'anglais]
\textbf{Traduisez les phrases suivantes en anglais en utilisant le vocabulaire et les structures vus (formules de politesse, lexique des services).}
\begin{enumerate}
\item Bonjour, je voudrais envoyer un colis express pour Douala.
\item Combien coûte l'envoi standard pour ce paquet de 3 kg ?
\item Pourriez-vous me donner un reçu, s'il vous plaît ?
\item Je n'ai pas de justificatif de domicile sur moi.
\item Le guichetier m'a dit que la carte arriverait dans cinq jours.
\item Est-ce que je peux payer par carte Mobile Money ?
\end{enumerate}
\end{exercice}

\begin{exercice}[Transformation : Registre de langue / Politesse]
\textbf{Transformez les phrases directes et familières en demandes polies et formelles adaptées à un guichet (banque, poste, mairie).}
\begin{enumerate}
\item \eng{Give me a form.} $\rightarrow$ \dots
\item \eng{I want to send this letter.} $\rightarrow$ \dots
\item \eng{How much is it?} $\rightarrow$ \dots
\item \eng{Wait a minute.} $\rightarrow$ \dots
\item \eng{Where do I sign?} $\rightarrow$ \dots
\item \eng{Call me when it's ready.} $\rightarrow$ \dots
\end{enumerate}
\end{exercice}

\begin{exercice}[Texte à trous : Dialogue à la banque]
\textbf{Complétez le dialogue entre un client et un employé de banque avec les expressions proposées : \eng{Certainly}, \eng{I'd like to}, \eng{Could you}, \eng{I'm afraid}, \eng{Here is}, \eng{Would you like}, \eng{Please sign}, \eng{Thank you}.}
\begin{textebox}[At the Bank]
Customer: Good morning. \_\_\_\_\_\_\_\_ withdraw some money from my savings account.
Clerk: \_\_\_\_\_\_\_\_. \_\_\_\_\_\_\_\_ please show me your ID card and your passbook?
Customer: Yes, of course. \_\_\_\_\_\_\_\_ they are.
Clerk: \_\_\_\_\_\_\_\_. How much would you like to withdraw?
Customer: 50,000 CFA francs, please.
Clerk: \_\_\_\_\_\_\_\_ this withdrawal slip. \_\_\_\_\_\_\_\_ you like the money in large or small bills?
Customer: Large bills, please.
Clerk: \_\_\_\_\_\_\_\_ your money. Have a good day.
Customer: \_\_\_\_\_\_\_\_. Goodbye.
\end{textebox}
\end{exercice}

\begin{exercice}[Appariements : Situations et réponses]
\textbf{Associez chaque situation (1--6) à la réplique la plus appropriée du personnel de service (A--F).}
\begin{center}
\begin{tabularx}{\linewidth}{|X|X|}
\hline
\textbf{Situation} & \textbf{Réplique du personnel} \\
\hline
1. Un client arrive à la réception d'un hôtel sans réservation. & A. \eng{Certainly, sir. Would you like a single or a double room?} \\
\hline
2. Un client signale que sa carte SIM ne fonctionne plus. & B. \eng{I'm afraid we are fully booked for tonight.} \\
\hline
3. Un client veut payer sa facture d'électricité en retard. & C. \eng{Please fill in this claim form and we will replace it.} \\
\hline
4. Un client dépose une valise abîmée au service bagages. & D. \eng{We can process the payment here. Do you have the bill reference?} \\
\hline
5. Un client demande un renseignement sur les horaires de bus. & E. \eng{The next bus to Yaoundé leaves at 7:00 AM from platform 3.} \\
\hline
6. Un client veut ouvrir un compte étudiant. & F. \eng{You need your student card and a proof of address.} \\
\hline
\end{tabularx}
\end{center}
\end{exercice}

\courssub{Je me prépare au Bac}

\begin{exercice}[Compréhension de l'oral / Écrit type Bac : Texte long + Questions]
\textbf{Lisez attentivement le texte suivant puis répondez aux questions qui suivent.}
\begin{textebox}[A Visit to the Telecom Agency]
Last Monday, Mr. Nfor, a teacher at Government High School Bamenda, went to the MTN agency at Commercial Avenue to solve a problem with his mobile line. He had been unable to make calls for two days, although he had credit on his account. When he arrived at the agency around 10:00 AM, there was a long queue. He took a ticket from the machine: number 412. The screen displayed number 398.

While waiting, he observed the agents. They were busy attending to customers with various requests: SIM card registration, mobile money deposits, data bundle subscriptions, and complaints about network coverage. After about forty minutes, his number appeared on the screen. He walked to counter 3.

"Good morning, sir. How can I assist you?" the agent asked politely.
"Good morning. I have a problem with my line. I cannot make outgoing calls," Mr. Nfor explained.
The agent asked for his phone number and his national identity card. He typed on his keyboard for a few minutes. "I see the issue, sir. Your SIM card is damaged. We need to do a SIM swap. It costs 1,000 CFA francs and you will keep the same number. Do you agree?"
"Yes, I agree. How long will it take?"
"About fifteen minutes. Please have a seat."
Fifteen minutes later, the agent handed him a new SIM card. "Here is your new SIM. It is already active. Please insert it into your phone and dial *123\# to check your balance."
Mr. Nfor did as instructed. His phone immediately connected to the network. "It works! Thank you very much for your help."
"You are welcome, sir. Enjoy your day."
Mr. Nfor left the agency satisfied. He could finally call the school principal to confirm his attendance for the week.
\end{textebox}

\textbf{Questions de compréhension (répondez en anglais, phrases complètes) :}
\begin{enumerate}
\item Why did Mr. Nfor go to the MTN agency?
\item What was the number on his ticket and what number was being served when he arrived?
\item List three different services the agents were providing to other customers.
\item What was the diagnosis for Mr. Nfor's phone problem?
\item How much did the solution cost and how long did it take?
\item What did the agent ask Mr. Nfor to do to check if the new SIM worked?
\end{enumerate}

\textbf{Questions de langue :}
\begin{enumerate}
\setcounter{enumi}{6}
\item Find in the text the synonyms for: (a) "queue" (line 4), (b) "fix/solve" (line 14), (c) "broken" (line 17).
\item Put the following verbs from the text in the Past Simple: (a) \eng{go}, (b) \eng{have}, (c) \eng{take}, (d) \eng{give}, (e) \eng{leave}.
\item Rewrite the sentence in reported speech: The agent asked: "Do you agree?" $\rightarrow$ The agent asked him \dots
\item Turn into the passive voice: "They handed him a new SIM card."
\end{enumerate}
\end{exercice}

\begin{exercice}[Traduction : Extrait du texte vers le français]
\textbf{Traduisez en français le passage suivant (lignes 14 à 20 du texte ci-dessus) :}
\begin{quote}
"The agent asked for his phone number and his national identity card. He typed on his keyboard for a few minutes. 'I see the issue, sir. Your SIM card is damaged. We need to do a SIM swap. It costs 1,000 CFA francs and you will keep the same number. Do you agree?'"
\end{quote}
\end{exercice}

\begin{exercice}[Expression écrite : Sujet de composition type Bac]
\textbf{Choisissez \textbf{UN} des deux sujets suivants. Votre composition comportera environ \textbf{150 à 180 mots}. Vous respecterez le format demandé (lettre formelle ou article).}

\textbf{Sujet A (Lettre formelle) :}
Vous êtes \eng{Ngono Essomba}, élève en classe de Première A au Lycée de Nkolbisson (Yaoundé). Vous avez acheté un smartphone neuf chez \eng{"Techno World"} au quartier Mvog-Ada. Après trois jours, l'écran tactile ne répond plus. Écrivez une lettre de réclamation au gérant de la boutique pour demander un échange ou un remboursement. Mentionnez : la date d'achat, le numéro de facture, le défaut constaté, vos démarches précédentes (appel téléphonique sans suite) et votre demande précise.

\textbf{Sujet B (Article pour le journal scolaire) :}
Vous êtes rédacteur au journal de votre établissement. Écrivez un article intitulé : \eng{"How to Request a Service Effectively in Public Places"}. Donnez au moins 4 conseils pratiques (politesse, documents nécessaires, heures creuses, note-taking, suivi) et illustrez avec un exemple vécu (poste, banque, mairie, agence de téléphonie).
\end{exercice}

\newpage

\coursec{Corrigés des exercices}

\begin{corrige}[Exercice 1 — QCM : Compréhension globale d'un dialogue]
\begin{enumerate}
\item \textbf{B.} At a post office. — Le titre du script et le vocabulaire (\eng{send this parcel}, \eng{express delivery}, \eng{tracking number}) indiquent clairement une poste.
\item \textbf{B.} Send a parcel. — \eng{“I'd like to send this parcel to Yaoundé, please.”}
\item \textbf{C.} 3,500 CFA francs. — \eng{“That will be 3,500 CFA francs for express delivery.”} Attention à ne pas confondre avec les 2,000 francs de la livraison standard.
\item \textbf{B.} Tomorrow before noon. — \eng{“It arrives tomorrow before noon.”} Les trois jours correspondent à la livraison standard.
\item \textbf{C.} A receipt and a tracking number. — \eng{“Here is your receipt and the tracking number.”}
\end{enumerate}
\textbf{Méthode :} avant d'écouter, lisez les questions pour savoir quoi repérer (lieu, intention, chiffres, objets remis). Pendant l'écoute, notez les chiffres et les noms propres dès qu'ils sont prononcés.
\end{corrige}

\begin{corrige}[Exercice 2 — Vrai / Faux justifié : Repérage de détails]
\begin{enumerate}
\item \textbf{F (Faux).} — \eng{“I would like to open a savings account.”} Le client veut un compte d'épargne, pas un compte courant.
\item \textbf{F (Faux).} — \eng{“I have an electricity bill from last month.”} Il s'agit d'une facture d'électricité, pas d'eau.
\item \textbf{V (Vrai).} — \eng{“The minimum opening deposit is 10,000 CFA francs.”}
\item \textbf{F (Faux).} — \eng{“I study at the University of Douala.”} Le client étudie à Douala, pas à Yaoundé.
\item \textbf{F (Faux).} — \eng{“You will receive your ATM card by post within five working days.”} La carte arrive par courrier sous cinq jours ouvrés, pas immédiatement au guichet.
\end{enumerate}
\textbf{Point de vigilance :} en compréhension, chaque détail compte (type de compte, type de facture, ville, délai). Les distracteurs jouent toujours sur des mots proches présents dans le texte.
\end{corrige}

\begin{corrige}[Exercice 3 — Vocabulaire : Les mots-clés de la demande de service]
\begin{enumerate}
\item \eng{parcel} — I need to send this \textbf{parcel} to my brother in Bamenda. (colis)
\item \eng{counter} — Please go to \textbf{counter} number 3 to pay your bill. (guichet)
\item \eng{tracking number} — The \textbf{tracking number} allows you to follow your package online. (numéro de suivi)
\item \eng{receipt} — Keep the \textbf{receipt} as proof of payment. (reçu)
\item \eng{express delivery} — For \textbf{express delivery}, the package arrives in 24 hours. (livraison express)
\item \eng{form} — You must fill in this \textbf{form} with your personal details. (formulaire)
\item \eng{proof of address} — A recent electricity bill serves as \textbf{proof of address}. (justificatif de domicile)
\item \eng{deposit} — The initial \textbf{deposit} for this account is 5,000 francs. (dépôt initial)
\end{enumerate}
\textbf{Astuce :} repérez les indices grammaticaux : \eng{this} + nom singulier, \eng{the} + nom, \eng{as} + nom de fonction.
\end{corrige}

\begin{corrige}[Exercice 4 — Grammaire : Modaux et formes de politesse]
\begin{enumerate}
\item \eng{Could} you please tell me where the nearest post office is? — \eng{Could} + pronom + verbe : demande polie d'information. (Pourriez-vous me dire où est la poste la plus proche ?)
\item I \eng{would} like to book a table for two for tonight. — \eng{would like} = « je voudrais », formule de politesse standard. (Je voudrais réserver une table pour deux ce soir.)
\item \eng{May} I help you with your luggage, madam? — \eng{May I...?} : proposition d'aide très polie. (Puis-je vous aider avec vos bagages, madame ?)
\item \eng{Could} you check if this book is available in the library? — demande polie d'un service. (Pourriez-vous vérifier si ce livre est disponible ?)
\item We \eng{would} be grateful if you could send the confirmation by email. — \eng{We would be grateful if you could...} : formule très formelle de correspondance. (Nous vous serions reconnaissants de bien vouloir envoyer la confirmation par email.)
\item \eng{Would} you mind opening the window? — \eng{Would you mind} + V-ing : demande très polie. (Cela vous dérangerait-il d'ouvrir la fenêtre ?)
\end{enumerate}
\textbf{Rappel :} après un modal, le verbe est toujours à la base verbale ; seul \eng{would you mind} est suivi du gérondif.
\end{corrige}

\begin{corrige}[Exercice 5 — Traduction : Du français vers l'anglais]
\begin{enumerate}
\item \eng{Good morning, I would like to send an express parcel to Douala.}
\item \eng{How much does standard delivery cost for this 3 kg package?}
\item \eng{Could you give me a receipt, please?}
\item \eng{I do not have any proof of address with me.}
\item \eng{The clerk told me that the card would arrive within five days.}
\item \eng{Can I pay by Mobile Money?}
\end{enumerate}
\textbf{Points de vigilance :}
\begin{itemize}
\item « Je voudrais » se traduit par \eng{I would like}, jamais \eng{I want} (trop direct).
\item « Combien coûte... ? » : \eng{How much does ... cost?} — ne pas oublier l'auxiliaire \eng{does}.
\item Phrase 5 : discours rapporté, recul du temps : \eng{will} $\rightarrow$ \eng{would}.
\item \eng{proof of address} est invariable dans cette expression ; avec \eng{any}, on ne met pas d'article devant \eng{proof}.
\end{itemize}
\end{corrige}

\begin{corrige}[Exercice 6 — Transformation : Registre de langue / Politesse]
Réponses proposées (d'autres formulations polies sont acceptables) :
\begin{enumerate}
\item \eng{Give me a form.} $\rightarrow$ \eng{Could I have a form, please?}
\item \eng{I want to send this letter.} $\rightarrow$ \eng{I would like to send this letter, please.}
\item \eng{How much is it?} $\rightarrow$ \eng{Could you tell me how much it costs, please?}
\item \eng{Wait a minute.} $\rightarrow$ \eng{Would you mind waiting a moment, please?}
\item \eng{Where do I sign?} $\rightarrow$ \eng{Could you show me where I should sign, please?}
\item \eng{Call me when it's ready.} $\rightarrow$ \eng{Would you mind calling me when it is ready, please?}
\end{enumerate}
\textbf{Méthode :} pour passer au registre formel, on (1) remplace l'impératif par \eng{Could you...?} ou \eng{Would you mind...?}, (2) remplace \eng{want} par \eng{would like}, (3) ajoute \eng{please}, (4) évite les contractions familières.
\end{corrige}

\begin{corrige}[Exercice 7 — Texte à trous : Dialogue à la banque]
\textbf{Corrigé (ordre des cases) :}
\begin{center}
\begin{tabular}{|l|l|}
\hline
\textbf{Case} & \textbf{Réponse} \\
\hline
1 (client) & \eng{I'd like to} \\
\hline
2 (guichetier) & \eng{Certainly} \\
\hline
3 (guichetier) & \eng{Could you} \\
\hline
4 (client) & \eng{Here they are} \\
\hline
5 (guichetier) & \eng{Thank you} \\
\hline
6 (guichetier) & \eng{Please sign} \\
\hline
7 (guichetier) & \eng{Would you} \\
\hline
8 (guichetier) & \eng{Here is} \\
\hline
\end{tabular}
\end{center}
\textbf{Dialogue complet reconstitué :}
\begin{itemize}
\item Customer: Good morning. \eng{I'd like to} withdraw some money from my savings account.
\item Clerk: \eng{Certainly}. \eng{Could you} please show me your ID card and your passbook?
\item Customer: Yes, of course. \eng{Here they are}.
\item Clerk: \eng{Thank you}. How much would you like to withdraw?
\item Customer: 50,000 CFA francs, please.
\item Clerk: \eng{Please sign} this withdrawal slip. \eng{Would you} like the money in large or small bills?
\item Customer: Large bills, please.
\item Clerk: \eng{Here is} your money. Have a nice day!
\end{itemize}
\textbf{Note pédagogique :} \eng{I'm afraid} ne convient à aucune case ici : il sert à annoncer une mauvaise nouvelle (\eng{I'm afraid we are closed.}) ; c'est un distracteur. Attention à l'accord : \eng{Here is} + singulier (\eng{Here is your money.}), \eng{Here are} / \eng{Here they are} + pluriel.
\end{corrige}

\begin{corrige}[Exercice 8 — Appariements : Situations et réponses]
\begin{center}
\begin{tabular}{|c|c|l|}
\hline
\textbf{Situation} & \textbf{Réponse} & \textbf{Indice-clé} \\
\hline
1 & B & \eng{fully booked} = complet (hôtel sans réservation) \\
\hline
2 & C & \eng{claim form} + \eng{replace it} = carte SIM défectueuse \\
\hline
3 & D & \eng{process the payment} + \eng{bill reference} = facture \\
\hline
4 & A & \eng{single or double room} = enregistrement à l'hôtel \\
\hline
5 & E & \eng{bus} + \eng{platform 3} = horaires de bus \\
\hline
6 & F & \eng{student card} + \eng{proof of address} = compte étudiant \\
\hline
\end{tabular}
\end{center}
\textbf{Correction détaillée :}
\begin{itemize}
\item Situation 1 (client à l'hôtel \emph{sans réservation}, l'hôtel est complet) $\rightarrow$ \textbf{B} : \eng{I'm afraid we are fully booked for tonight.} L'employé annonce une mauvaise nouvelle, d'où \eng{I'm afraid}.
\item Situation 2 (carte SIM qui ne fonctionne plus) $\rightarrow$ \textbf{C} : \eng{Please fill in this claim form and we will replace it.} Le formulaire de réclamation et le remplacement correspondent exactement au problème.
\item Situation 3 (paiement d'une facture d'électricité) $\rightarrow$ \textbf{D} : \eng{I can process the payment if you give me your bill reference.}
\item Situation 4 (client qui arrive à la réception d'un hôtel \emph{avec} réservation) $\rightarrow$ \textbf{A} : \eng{Would you like a single or a double room?} L'agent propose le type de chambre.
\item Situation 5 (renseignement sur les départs de bus) $\rightarrow$ \textbf{E} : \eng{The next bus leaves from platform 3 at 2.30 p.m.}
\item Situation 6 (ouverture d'un compte étudiant) $\rightarrow$ \textbf{F} : \eng{You will need your student card and a proof of address.}
\end{itemize}
\textbf{Piège classique :} ne pas confondre les situations 1 et 4, qui se passent toutes les deux à l'hôtel. C'est le contexte qui tranche : sans réservation et hôtel complet $\rightarrow$ \eng{fully booked} ; avec réservation $\rightarrow$ proposition d'une chambre simple ou double.
\end{corrige}

\begin{corrige}[Exercice 9 — Compréhension type Bac : Texte long + Questions]
\textbf{Questions de compréhension (réponses modèles) :}
\begin{enumerate}
\item \eng{Mr Nfor went to the MTN agency because he had been unable to make calls for two days, although he had credit on his account.}
\item \eng{His ticket number was 412 and the screen was displaying number 398 when he arrived.}
\item \eng{The agents were providing SIM card registration, mobile money deposits, data bundle subscriptions, and handling complaints about network coverage.} (trois de ces services suffisent)
\item \eng{The agent diagnosed that his SIM card was damaged and that he needed a SIM swap.}
\item \eng{The SIM swap cost 1,000 CFA francs and it took about fifteen minutes.}
\item \eng{The agent asked him to insert the new SIM into his phone and to dial *123\# to check his balance.}
\end{enumerate}
\textbf{Questions de langue :}
\begin{enumerate}
\setcounter{enumi}{6}
\item (a) queue = \eng{line} ; (b) fix = \eng{solve} ; (c) broken = \eng{damaged}.
\item (a) \eng{go} $\rightarrow$ \eng{went} ; (b) \eng{have} $\rightarrow$ \eng{had} ; (c) \eng{take} $\rightarrow$ \eng{took} ; (d) \eng{give} $\rightarrow$ \eng{gave} ; (e) \eng{leave} $\rightarrow$ \eng{left}.
\item \eng{The agent asked him if he agreed.} — discours indirect : question fermée introduite par \eng{if}, suppression du point d'interrogation, recul du présent vers le prétérit (\eng{do you agree} $\rightarrow$ \eng{he agreed}).
\item \eng{He was handed a new SIM card.} ou \eng{A new SIM card was handed to him.} — voix passive : auxiliaire \eng{be} au prétérit + participe passé \eng{handed}.
\end{enumerate}
\textbf{Barème indicatif (sur 20) :} compréhension 6 $\times$ 1,5 = 9 pts ; synonymes 3 pts ; verbes irréguliers 2,5 pts ; discours rapporté 2,5 pts ; voix passive 3 pts. La qualité de la langue est prise en compte dans chaque réponse.
\textbf{Points de vigilance :} répondre par des phrases complètes ; ne pas recopier tout le texte ; au discours rapporté, ne jamais garder le point d'interrogation ; à la voix passive, ne pas oublier l'accord de \eng{be} (\eng{was}, singulier).
\end{corrige}

\begin{corrige}[Exercice 10 — Traduction : Extrait du texte vers le français]
« L'agent lui demanda son numéro de téléphone et sa carte nationale d'identité. Il tapa sur son clavier pendant quelques minutes. “Je vois le problème, monsieur. Votre carte SIM est endommagée. Nous devons procéder à un échange de SIM (un \emph{SIM swap}). Cela coûte 1 000 francs CFA et vous conserverez le même numéro. Êtes-vous d'accord ?” »
\textbf{Points de vigilance :}
\begin{itemize}
\item \eng{issue} est un faux ami : ici « problème », pas « question » ou « sortie ».
\item \eng{damaged} = « endommagée » (et non « damnée » !).
\item \eng{you will keep the same number} = « vous conserverez le même numéro » (futur de garantie).
\item \eng{Do you agree?} se traduit naturellement par « Êtes-vous d'accord ? » ou « Acceptez-vous ? ».
\end{itemize}
\end{corrige}

\begin{corrige}[Exercice 11 — Expression écrite : Sujet de composition type Bac]
\textbf{Barème indicatif (sur 20) :} respect du format et de la consigne 4 pts ; contenu et pertinence des informations 6 pts ; organisation et cohérence 4 pts ; correction grammaticale et orthographique 4 pts ; richesse du vocabulaire 2 pts.

\textbf{Proposition de rédaction modèle — Sujet A (lettre formelle) :}
\begin{textebox}[Model answer — Formal letter of complaint]
Ngono Essomba\\
Lycée de Nkolbisson\\
P.O. Box 1234, Yaoundé\\
15th May 2024

The Manager\\
Techno World\\
Mvog-Ada, Yaoundé

Dear Sir or Madam,

I am writing to complain about a smartphone I bought in your shop on 10th May 2024 (invoice number TW-0457). After only three days of normal use, the touch screen stopped responding completely.

I called your customer service twice last week, but nobody has contacted me since then. As the phone is still under warranty, I would be grateful if you could either exchange it for a new one or refund my money.

I enclose a copy of the receipt. I look forward to hearing from you within a week.

Yours faithfully,\\
Ngono Essomba
\end{textebox}
(Environ 115 mots pour le corps ; avec les adresses, on atteint la fourchette demandée.)

\textbf{Proposition de rédaction modèle — Sujet B (article) :}
\begin{textebox}[Model answer — School magazine article]
How to Request a Service Effectively in Public Places

Have you ever spent a whole morning at a counter without getting what you wanted? Here are four practical tips to avoid that frustration.

First, always be polite: start with “Good morning” and use “Could you...?” instead of giving orders. Second, prepare all your documents before leaving home: ID card, bills, receipts or reference numbers. Third, avoid rush hours; go early in the morning or in the middle of the week. Finally, take notes: write down names, dates and reference numbers so that you can follow up your request later.

Last month, I went to the post office in Yaoundé to send a parcel. Because I had my receipt and my ID ready, the whole operation took ten minutes. My neighbour, who had forgotten his documents, had to come back the next day!

Good preparation and politeness really open doors.
\end{textebox}
(Environ 150 mots.)

\textbf{Points de vigilance :}
\begin{itemize}
\item Lettre formelle : adresses en haut, date, formule d'appel \eng{Dear Sir or Madam}, formule de politesse finale \eng{Yours faithfully} (destinataire inconnu).
\item Vocabulaire de la réclamation : \eng{complain about}, \eng{under warranty}, \eng{exchange}, \eng{refund}, \eng{I look forward to hearing from you} (+ V-ing !).
\item Article : titre, paragraphes courts, impératifs et conseils, exemple personnel au prétérit.
\item Éviter les faux amis : \eng{actually} = « en fait » (pas « actuellement ») ; \eng{to assist} = « aider » (pas « assister à »).
\end{itemize}
\end{corrige}

\newpage

\coursec{Fiche bilan}

\begin{popfigure}
\begin{center}\tcbox[colback=poporange,colframe=poporange,coltext=white,arc=9pt,boxsep=4pt,left=6pt,right=6pt]{\bfseries\small Requesting Services}\end{center}
\vspace{-2pt}
\begin{tcbraster}[raster columns=3,raster equal height=rows,raster column skip=6pt,raster row skip=6pt,enhanced,blankest]
\begin{tcolorbox}[enhanced,colback=white,colframe=poporange,boxrule=0.9pt,arc=3pt,title={Listening Strategies},fonttitle=\bfseries\scriptsize,coltitle=white,colbacktitle=poporange,left=4pt,right=4pt,top=2pt,bottom=2pt,boxsep=1pt,toptitle=1.5pt,bottomtitle=1.5pt,halign=left]
\begin{itemize}[nosep,leftmargin=*,itemsep=1pt,font=\scriptsize]
\item Keywords
\item Numbers/Names
\item Main Idea
\end{itemize}
\end{tcolorbox}
\begin{tcolorbox}[enhanced,colback=white,colframe=popgreen,boxrule=0.9pt,arc=3pt,title={Polite Requests},fonttitle=\bfseries\scriptsize,coltitle=white,colbacktitle=popgreen,left=4pt,right=4pt,top=2pt,bottom=2pt,boxsep=1pt,toptitle=1.5pt,bottomtitle=1.5pt,halign=left]
\begin{itemize}[nosep,leftmargin=*,itemsep=1pt,font=\scriptsize]
\item Could you...?
\item Would you mind...?
\item I would like...
\end{itemize}
\end{tcolorbox}
\begin{tcolorbox}[enhanced,colback=white,colframe=poppurple,boxrule=0.9pt,arc=3pt,title={Service Vocabulary},fonttitle=\bfseries\scriptsize,coltitle=white,colbacktitle=poppurple,left=4pt,right=4pt,top=2pt,bottom=2pt,boxsep=1pt,toptitle=1.5pt,bottomtitle=1.5pt,halign=left]
\begin{itemize}[nosep,leftmargin=*,itemsep=1pt,font=\scriptsize]
\item Booking
\item Complaint
\item Information
\end{itemize}
\end{tcolorbox}
\begin{tcolorbox}[enhanced,colback=white,colframe=poppink,boxrule=0.9pt,arc=3pt,title={Grammar Focus},fonttitle=\bfseries\scriptsize,coltitle=white,colbacktitle=poppink,left=4pt,right=4pt,top=2pt,bottom=2pt,boxsep=1pt,toptitle=1.5pt,bottomtitle=1.5pt,halign=left]
\begin{itemize}[nosep,leftmargin=*,itemsep=1pt,font=\scriptsize]
\item Modals
\item Conditionals
\item Indirect Q.
\end{itemize}
\end{tcolorbox}
\begin{tcolorbox}[enhanced,colback=white,colframe=popgold,boxrule=0.9pt,arc=3pt,title={Contexts},fonttitle=\bfseries\scriptsize,coltitle=white,colbacktitle=popgold,left=4pt,right=4pt,top=2pt,bottom=2pt,boxsep=1pt,toptitle=1.5pt,bottomtitle=1.5pt,halign=left]
\begin{itemize}[nosep,leftmargin=*,itemsep=1pt,font=\scriptsize]
\item Hotel/Reception
\item Bank/Post Office
\item Admin/Camrail
\end{itemize}
\end{tcolorbox}
\begin{tcolorbox}[enhanced,colback=white,colframe=popblue,boxrule=0.9pt,arc=3pt,title={Production},fonttitle=\bfseries\scriptsize,coltitle=white,colbacktitle=popblue,left=4pt,right=4pt,top=2pt,bottom=2pt,boxsep=1pt,toptitle=1.5pt,bottomtitle=1.5pt,halign=left]
\begin{itemize}[nosep,leftmargin=*,itemsep=1pt,font=\scriptsize]
\item Role-play
\item Dialogue
\item Formal Letter
\end{itemize}
\end{tcolorbox}
\end{tcbraster}

\legende{Carte mentale : Lesson 16 — Requesting Services}
\end{popfigure}

\begin{bilan}

\textbf{Lexique clé — Key Vocabulary (English / Français)}

\begin{center}
\begin{tabularx}{\linewidth}{|X|X|X|}
\hline
\rowcolor{popblueL} \textbf{English} & \textbf{Français} & \textbf{Contexte / Collocation} \\
\hline
\cle{Reception desk} & Guichet d'accueil / Réception & \eng{Go to the reception desk.} \\
\cle{Booking / Reservation} & Réservation & \eng{Make a booking for a room.} \\
\cle{Check-in / Check-out} & Arrivée / Départ (hôtel) & \eng{Check-in is at 2 PM.} \\
\cle{Single / Double room} & Chambre simple / double & \eng{I'd like a single room.} \\
\cle{Available / Fully booked} & Disponible / Complet & \eng{We are fully booked tonight.} \\
\cle{Complaint} & Réclamation & \eng{Make a complaint about the noise.} \\
\cle{Refund} & Remboursement & \eng{Ask for a refund.} \\
\cle{Receipt} & Reçu & \eng{Keep your receipt.} \\
\cle{Form / Application form} & Formulaire & \eng{Fill in this form, please.} \\
\cle{ID card / Passport} & Carte d'identité / Passeport & \eng{Show your ID card.} \\
\cle{Queue / Line} & File d'attente & \eng{Wait in the queue.} \\
\cle{Assistance / Help} & Assistance / Aide & \eng{Do you need assistance?} \\
\cle{Service charge} & Frais de service & \eng{Service charge included.} \\
\hline
\end{tabularx}
\end{center}

\vspace{0.5cm}

\textbf{Structures grammaticales et fonctionnelles — Grammar \& Functions}

\begin{center}
\begin{tabularx}{\linewidth}{|X|X|X|}
\hline
\rowcolor{poporangeL} \textbf{Fonction} & \textbf{Structure / Exemple} & \textbf{Niveau de politesse} \\
\hline
Demande directe & \eng{Can you help me?} & Standard \\
Demande polie & \eng{Could you help me, please?} & Poli \\
Demande très formelle & \eng{Would you mind helping me?} & Très poli / Formel \\
Exprimer un besoin & \eng{I would like to book a room.} & Standard (Would like) \\
Demande indirecte & \eng{I was wondering if I could book...} & Très poli (hésitation) \\
Demander la permission & \eng{May I see the manager?} & Formel \\
Répondre positivement & \eng{Certainly. / Of course. / Right away.} & -- \\
Répondre négativement & \eng{I'm afraid not. / I'm sorry, but...} & -- \\
\hline
\end{tabularx}
\end{center}

\vspace{0.5cm}

\textbf{Méthodes Express — Listening Strategies}

\end{bilan}

\begin{methode}[Réussir l'écoute « Requesting Services »]
\emph{Étape 1 — Avant l'écoute :} Lire les questions, souligner les \cle{mots-clés} (noms, verbes, chiffres), prédire le contexte (hôtel, banque, gare).
\emph{Étape 2 — Première écoute :} Identifier le \cle{thème général}, les \cle{interlocuteurs} (client/employé), le \cle{lieu} et le \cle{but} de la demande.
\emph{Étape 3 — Deuxième écoute :} Noter les \cle{détails précis} : dates, heures, numéros, noms propres, prix, numéros de chambre/vol.
\emph{Étape 4 — Vérification :} Revoir les questions Vrai/Faux ou QCM, attention aux \cle{pièges} (négation, changement d'avis, "but", "however").
\end{methode}

\vspace{0.5cm}

\textbf{Expressions utiles pour la production orale/écrite}

\begin{itemize}
\item \eng{Good morning, I have a reservation under the name of...} (Bonjour, j'ai une réservation au nom de...)
\item \eng{Could you tell me where the nearest bank is?} (Pourriez-vous me dire où est la banque la plus proche ?)
\item \eng{I would like to report a problem with...} (Je voudrais signaler un problème avec...)
\item \eng{Is it possible to change my ticket?} (Est-il possible de changer mon billet ?)
\item \eng{Thank you for your assistance.} (Merci pour votre aide.)
\end{itemize}



\begin{aretenir}[Checklist avant le Bac]
$\square$ Je sais identifier le lieu et les interlocuteurs dès la première écoute. \\
$\square$ Je reconnais les formules de politesse : \eng{Could you}, \eng{Would you mind}, \eng{I would like}. \\
$\square$ Je note correctement les chiffres : prix, dates, heures, numéros de téléphone/réservation. \\
$\square$ Je distingue une \cle{demande} d'une \cle{réclamation} (complaint) à l'oral. \\
$\square$ Je connais le vocabulaire de base : \eng{booking, check-in, receipt, refund, form, queue}. \\
$\square$ Je sais reformuler une demande directe en demande indirecte (plus polie). \\
$\square$ Je repère les marqueurs de contraste : \eng{but, however, unfortunately, I'm afraid}. \\
$\square$ Je suis capable de jouer un dialogue « Guichet / Client » en 5-6 répliques. \\
$\square$ Je sais rédiger une courte réclamation écrite formelle (structure : faits $\to$ problème $\to$ solution attendue). \\
$\square$ Je ne confonds pas \eng{receipt} (reçu) et \eng{recipe} (recette de cuisine) — \textbf{faux ami}. \\
$\square$ Je prononce correctement \eng{reservation} [re-zer-vei-chn] et \eng{queue} [kiou]. \\
$\square$ Je gère mon stress : je ne panique pas si je rate un mot, j'écoute le contexte global.
\end{aretenir}

\findecours

\end{document}
